\documentclass[11pt,a4paper]{article}
\usepackage[utf8]{inputenc}
\usepackage[T1]{fontenc}
\usepackage[czech]{babel}
\usepackage[margin=2.2cm]{geometry}
\usepackage{amsmath,amssymb,amsthm}
\usepackage{parskip}
\usepackage{enumitem}
\setlist{nosep,leftmargin=1.4em}
\usepackage{booktabs}
\usepackage{array}
\usepackage{xcolor}
\definecolor{linkblue}{RGB}{20,60,150}
\definecolor{defbg}{RGB}{240,244,252}
\definecolor{defframe}{RGB}{100,130,190}
\definecolor{markcol}{RGB}{176,84,0}
\definecolor{markbg}{RGB}{253,246,236}
\usepackage[most]{tcolorbox}
\newtcolorbox{defbox}[1][]{breakable,colback=defbg,colframe=defframe,boxrule=0.6pt,arc=1.5mm,left=2mm,right=2mm,top=1mm,bottom=1mm,title={#1},fonttitle=\bfseries}
\newtcolorbox{poznbox}[1][]{breakable,colback=markbg,colframe=markcol,boxrule=0.6pt,arc=1.5mm,left=2mm,right=2mm,top=1mm,bottom=1mm,title={#1},fonttitle=\bfseries}
\newtheorem{veta}{Věta}[section]
\theoremstyle{definition}
\newtheorem{definice}{Definice}[section]
\usepackage[colorlinks=true,linkcolor=linkblue,urlcolor=linkblue]{hyperref}
\usepackage{algorithm}
\usepackage{algpseudocode}

% Značka pro látku, která není ve slidech NAIL106 (mas-18en.pdf).
\newcommand{\mimo}{\textcolor{markcol}{\textsf{\footnotesize[$\blacklozenge$~nad rámec slidů]}}}
\newcommand{\mimoq}[1]{\textcolor{markcol}{\textsf{\footnotesize[$\blacklozenge$~nad rámec slidů: #1]}}}

\title{\textbf{Multiagentní systémy}\\[0.3em]\large Učební text ke státní závěrečné zkoušce}
\author{SZZ Umělá inteligence, MFF UK}
\date{srpen 2026}

\begin{document}
\maketitle

\vspace{-1.5em}

\begin{center}
\begin{minipage}{0.93\textwidth}
\small
Text pokrývá okruh \uv{Multiagentní systémy} magisterské SZZ oboru Umělá inteligence.
Kapitoly jsou vedené \textbf{jedna ku jedné} podle vět oficiálního znění okruhu: každé
větě okruhu odpovídá právě jedna kapitola, jak ukazuje mapovací tabulka níže. Podkladem
jsou slidy přednášky NAIL106 \emph{Multiagentní systémy} (Neruda, Pilát, MFF UK),
doplněné o~učebnice M.~Wooldridge: \emph{An Introduction to MultiAgent Systems}
(2.~vyd., 2009), Y.~Shoham, K.~Leyton-Brown: \emph{Multiagent Systems} (2009)
a~S.~Russell, P.~Norvig: \emph{Artificial Intelligence: A~Modern Approach}.
Anglické termíny jsou při prvním výskytu uvedeny kurzívou v~závorce, protože literatura
k~oboru je výhradně anglická a~zkoušející je běžně používá.
\end{minipage}
\end{center}

\begin{center}
\begin{minipage}{0.93\textwidth}
\begin{poznbox}[Rozsah textu a~značení]
\small
Rozsah jednotlivých kapitol odpovídá váze látky. Nejvíce prostoru dostávají architektury
agentů, komunikace a~kooperace s~teorií her, tedy témata, kterým se přednáška věnuje
nejpodrobněji.

Značka \mimo{} upozorňuje na látku, která \textbf{není ve slidech} přednášky. V~textu
přesto zůstává, a~to ze dvou důvodů: buď ji \emph{jmenuje oficiální znění okruhu}, nebo
by bez ní okolní výklad nedával smysl. Je to tedy nabídka ke~škrtnutí: co označeno není,
je ve slidech doloženo.

Ve slidech chybějí celé kapitoly \ref{sec:cesta} (hledání cesty) a~\ref{sec:uceni}
(učení agentů), protože je přednáška explicitně uvádí mezi nepokrytými tématy. Okruh je
ale vyjmenovává, takže je nutné je umět. Jsou proto zpracovány stručně, na úrovni, která
na zkoušce obstojí.
\end{poznbox}
\end{minipage}
\end{center}

\section*{Mapování okruhu na kapitoly}
\addcontentsline{toc}{section}{Mapování okruhu na kapitoly}

\begin{center}
\small
\begin{tabular}{p{0.55\textwidth} p{0.38\textwidth}}
\toprule
\textbf{Věta oficiálního znění okruhu} & \textbf{Kapitola v~tomto textu} \\
\midrule
Architektura autonomního agenta; mechanismus výběru akcí.
  & \ref{sec:architektura}~Architektura autonomního agenta a~mechanismus výběru akcí \\
Metody pro řízení agentů; symbolické a~konekcionistické reaktivní plánování, hybridní
přístupy.
  & \ref{sec:rizeni}~Metody pro řízení agentů \\
Problém hledání cesty.
  & \ref{sec:cesta}~Problém hledání cesty \\
Komunikace a~znalosti v~multiagentních systémech, ontologie, řečové akty, FIPA-ACL,
protokoly.
  & \ref{sec:komunikace}~Komunikace a~znalosti v~MAS \\
Distribuované řešení problémů, kooperace, Nashova ekvilibria, Paretova efektivita, alokace
zdrojů, aukce.
  & \ref{sec:kooperace}~Distribuované řešení problémů, kooperace, teorie her a~aukce \\
Metody pro učení agentů; zpětnovazební učení.
  & \ref{sec:uceni}~Učení agentů a~zpětnovazební učení \\
Metodologie návrhu, jazyky a~prostředí multiagentních systémů.
  & \ref{sec:metodologie}~Metodologie návrhu, jazyky a~prostředí MAS \\
\bottomrule
\end{tabular}
\end{center}

\tableofcontents
\newpage

%=====================================================================
\section{Architektura autonomního agenta a~mechanismus výběru akcí}
\label{sec:architektura}
%=====================================================================

\subsection{Agent a~multiagentní systém}

Jako samostatná oblast informatiky se multiagentní systémy (\emph{multiagent systems}, MAS)
zformovaly v~90.~letech, a~to z~docela praktického důvodu. Výpočetní technika se rozšířila
všude a~začala být trvale propojená a~distribuovaná, takže výpočet přestal být osamocenou
úlohou a~stal se interakcí. Zároveň rostla složitost zadání natolik, že bylo výhodnější úlohu
softwaru delegovat, než ji krok po kroku předepsat. K~tomu se přidal tlak na snadnost použití.

Klasické paradigma \uv{program dostane vstup, spočítá výstup, skončí} zde přestává stačit:
software je do svého prostředí zasazený trvale a~to prostředí se mění nezávisle na něm.

\begin{defbox}[Agent --- pracovní definice]
\textbf{Russell, Norvig (AIMA):} Agent je cokoli, co lze chápat jako entitu \emph{vnímající}
své prostředí prostřednictvím senzorů a~\emph{působící} na toto prostředí pomocí efektorů.

\textbf{Pattie Maes:} Autonomní agenti jsou výpočetní systémy, které obývají nějaké složité
dynamické prostředí, autonomně v~něm vnímají a~jednají, a~tím naplňují množinu cílů či úkolů,
pro které byly navrženy.

\textbf{Wooldridge, Jennings:} Agent je počítačový systém zasazený do nějakého prostředí, schopný
v~tomto prostředí \emph{autonomního jednání} vedoucího ke splnění delegovaných cílů.
\end{defbox}

Jednotná definice agenta neexistuje. Franklin a~Graesser (1996) v~článku \emph{Is it an agent,
or just a program?} shromáždili desítky vzájemně nekompatibilních definic. Podstatnější než
přesné znění je proto to, co všechny sledují: odlišit agenta od běžné procedury či objektu.

\begin{defbox}[Multiagentní systém]
Multiagentní systém je systém složený z~více agentů, kteří spolu \emph{interagují} ---
nejčastěji výměnou zpráv po počítačové síti. Agenti \emph{nemusí sdílet společný cíl}; mohou mít
vlastní, i~protichůdné zájmy. Proto je nutné studovat mechanismy vyjednávání, koordinace
a~dynamického rozdělování zdrojů.
\end{defbox}

\textbf{Autonomie.} Autonomie není vlastnost typu ano/ne, ale škála. Člověk je plně autonomní,
metoda v~Javě není autonomní vůbec a~agent leží někde mezi nimi. Měl by autonomně volit
\emph{způsob}, jak zadaný problém vyřešit, tedy vybírat si podcíle. Hlavní cíl si ale nevolí:
ten je mu delegován.

\paragraph{Vymezení vůči příbuzným oborům.} Šíři okruhu vysvětluje to, že MAS lze číst
z~pěti různých stran.
\begin{itemize}
\item \textbf{Softwarové inženýrství.} Agent je dalším stupněm abstrakce v~řadě
  strojový kód $\to$ strukturované programy $\to$ objekty $\to$ distribuované objekty (CORBA)
  $\to$ agenti. Posun proti objektům vystihuje klasický bonmot: \uv{objekty to dělají zadarmo,
  agenti proto, že chtějí}.
\item \textbf{Distribuované a~paralelní systémy.} Desítky let výzkumu tu vyřešily mutex,
  deadlock i~konsensus, hotová řešení se ale přenést nedají. Agenti jsou totiž \emph{autonomní},
  koordinační mechanismy nejsou dané dopředu a~všechno se musí vyřešit až za běhu.
\item \textbf{Umělá inteligence.} Tradiční pohled řadí MAS mezi podoblasti UI. Russell a~Norvig
  to obracejí: cílem UI je podle nich návrh inteligentních agentů, tedy UI $\subseteq$ MAS.
  Etzioni situaci vyhrotil poznámkou, že \uv{MAS je 1~\% UI a~99~\% informatiky}. MAS techniky
  UI používá (reprezentace znalostí, plánování), zdůrazňuje však \emph{sociální} aspekty, které
  klasická UI přehlíží.
\item \textbf{Teorie her a~ekonomie.} Odtud pochází formální aparát pro racionální rozhodování
  více subjektů. Matematická teorie her přitom řeší hlavně otázky \emph{existence}, zatímco MAS
  zajímají \emph{výpočetní} aspekty: složitost a~praktická použitelnost.
\item \textbf{Sociologie.} Klíčovým pojmem MAS jsou \emph{společnosti agentů}. Platí tu stejná
  výhrada jako u~vztahu UI a~lidské inteligence, totiž že inspirace je vítaná, ale ne povinná
  (letadla nemávají křídly). Opačným směrem je vazba pevnější: sociologie, ekonomie i~ekologie
  používají agenty jako stavební kámen svých simulačních modelů, tedy \emph{agentové modelování}
  (\emph{agent-based modelling}, ABM).
\end{itemize}

\subsection{Vlastnosti inteligentního agenta}

Inteligentního agenta od pouhého programu odlišují tři vlastnosti, které musí mít současně.
\begin{itemize}
\item \textbf{Reaktivita} (\emph{reactive}) znamená, že agent vnímá prostředí a~dokáže na jeho
  změny reagovat v~rozumném, tedy reálném čase.
\item \textbf{Proaktivita} (\emph{proactive}) znamená, že agent má vlastní cíle a~aktivně je
  naplňuje z~vlastní iniciativy, nikoli až jako odpověď na vnější podnět.
\item \textbf{Sociálnost} (\emph{social ability}) znamená, že agent umí komunikovat s~jinými
  agenty i~s~lidmi.
\end{itemize}

Dosáhnout \emph{čistě} reaktivního nebo \emph{čistě} cílově orientovaného chování je snadné.
Obtížné je najít \textbf{rovnováhu mezi reaktivitou a~proaktivitou}: agent nesmí slepě
následovat plán, který se mezitím stal nesmyslným, ale ani nesmí cíle přehodnocovat tak často,
že nikam nedojde. Tento kompromis se vrací v~celém textu, nejzřetelněji u~odhodlání agenta
(odst.~\ref{ssec:commitment}) a~u~hybridních architektur (odst.~\ref{ssec:hybridni}).

\subsection{Intencionální systémy a~intencionální postoj}

Mluvíme-li o~agentech, přisuzujeme jim \emph{mentální stavy}: \uv{agent věří, že\dots},
\uv{agent chce\dots}. Daniel Dennett takový způsob popisu pojmenoval. \textbf{Intencionální
systém} je podle něj entita, jejíž chování lze predikovat přisouzením vlastností jako
přesvědčení (\emph{beliefs}), touhy (\emph{desires}) a~racionalita. Intencionální systémy navíc
tvoří hierarchii. Systém \emph{prvního řádu} má přesvědčení a~touhy, systém \emph{druhého řádu}
má přesvědčení a~touhy o~přesvědčeních a~touhách, svých i~cizích, a~tak dále.

Dennett rozlišuje tři \emph{postoje} (\emph{stances}), z~nichž lze chování systému predikovat.
\emph{Fyzikální} postoj vystačí s~fyzikálními zákony: hodím-li kamenem, nemusím mluvit o~jeho
touhách, stačí hmotnost, rychlost a~gravitace. \emph{Návrhový} postoj (\emph{design stance})
nastupuje tam, kde všechny fyzikální zákony neznám, vím ale, k~jakému účelu byl systém navržen;
budík umím nastavit, aniž bych rozuměl jeho elektronice. \emph{Intencionální} postoj popisuje
systém pomocí přesvědčení, tužeb a~racionality.

Shoham uvádí provokativní příklad: \uv{Vypínač v~místnosti je velmi kooperativní agent, který
je z~vlastní vůle ochoten vést elektrický proud, ale činí tak jen tehdy, věří-li, že si to
skutečně přejeme.} Popis je konzistentní a~elegantní, přesto většině lidí připadá absurdní,
\emph{protože pro vypínač máme jednodušší a~stejně přesný popis}. Intencionální postoj se tedy
vyplatí tam, kde fyzikální či návrhový popis neexistuje nebo je nepraktický. I~se schématem
počítače v~ruce je těžké odvodit, proč se po kliknutí na ikonu otevře okno.

\textbf{Z~pohledu informatiky je intencionální postoj především abstrakce pro zvládnutí
složitosti}, a~právě to je pro mnoho programátorů hlavní přínos MAS. Dennett sám k~tomu dodává
filozofickou otázku: intencionalita podle něj neexistuje sama o~sobě, vzniká teprve
\emph{v~procesu interpretace} pozorovatelem.

\subsection{Prostředí a~jeho vlastnosti}

Chování agenta zásadně určuje typ prostředí, ve kterém se nachází. Prostředí se proto třídí
podle čtyř dichotomií; následující klasifikace pochází od Russella a~Norviga.

\begin{center}
\begin{tabular}{p{0.28\textwidth} p{0.64\textwidth}}
\toprule
\textbf{Dichotomie} & \textbf{Význam} \\
\midrule
plně / částečně pozorovatelné \newline (\emph{fully / partially observable})
  & Senzory agenta mu dávají úplný stav prostředí --- nebo jen jeho část
    (skryté informace, šum, omezený dosah). \\
statické / dynamické \newline (\emph{static / dynamic})
  & Prostředí se mění pouze v~důsledku akcí agenta --- nebo i~jinak (jiní agenti, přírodní děje).
    Ve statickém prostředí si agent může \uv{v~klidu rozmyslet} další tah. \\
deterministické / nedeterministické
  & Každá akce má právě jeden možný výsledek --- nebo více (se známým či neznámým rozdělením
    pravděpodobnosti; pak mluvíme o~stochastickém prostředí). \\
diskrétní / spojité
  & Konečný počet stavů a~akcí --- nebo spojité veličiny (poloha, rychlost). \\
\bottomrule
\end{tabular}
\end{center}

\textbf{Nepříjemné konstatování:} většina zajímavých prostředí, tedy reálný svět i~internet,
padne pokaždé do té horší varianty. Jsou spojitá, nedeterministická, dynamická a~jen částečně
pozorovatelná. Právě proto nestačí naprogramovat agenta jako pevnou tabulku pravidel.

\medskip
\noindent\emph{Příklad --- termostat.} Termostat je agent, byť ne příliš inteligentní, zasazený
do prostředí, kterým je místnost. Rozeznává dva vjemy, \uv{je zima} a~\uv{teplota v~pořádku},
a~má dvě akce, zapnout a~vypnout. Jeho rozhodovací jednotka se vejde na jeden řádek:
zima $\Rightarrow$ zapni, OK $\Rightarrow$ vypni. Prostředí je přesto nedeterministické, protože
třeba někdo otevře dveře.

\noindent\emph{Příklad --- softwarový démon.} Vezměme unixový \texttt{xbiff}. Jeho prostředím je
operační systém, vnímá voláním systémových služeb a~jeho akcemi jsou softwarové operace, jako
spustit poštovního klienta nebo změnit ikonu. Rozhodovací algoritmus zůstává na úrovni
termostatu.

\subsection{Abstraktní architektura agenta}

Dosavadní popis byl neformální; teď pohled na agenta a~jeho interakci s~prostředím
zformalizujeme. Předpokládáme přitom, že prostředí je diskrétní. Není to omezující předpoklad,
protože každé spojité prostředí lze diskretizovat s~libovolnou přesností.

\begin{defbox}[Prostředí, agent, běh]
\textbf{Prostředí} může být v~jednom z~konečně mnoha diskrétních stavů,
$E = \{e, e', \dots\}$. \textbf{Agent} má k~dispozici konečnou množinu akcí
$A = \{a, a', \dots\}$.

\textbf{Běh} (\emph{run}) agenta v~prostředí je posloupnost střídajících se stavů prostředí
a~akcí:
\[
r = e_0 \xrightarrow{a_0} e_1 \xrightarrow{a_1} e_2 \xrightarrow{a_2} \cdots
\xrightarrow{a_{n-1}} e_n .
\]
Označme $R$ množinu všech konečných běhů (nad $E$ a~$A$), $R^A \subseteq R$ ty, které končí akcí,
a~$R^E \subseteq R$ ty, které končí stavem prostředí.
\end{defbox}

\begin{defbox}[Přechodová funkce prostředí (\emph{state transformer function})]
\[
\tau : R^A \to 2^E
\]
Funkce $\tau$ zobrazuje běh končící akcí na množinu stavů, do nichž se prostředí může
po této akci dostat. Z~této definice plynou dva důsledky:
\begin{itemize}
\item \textbf{Prostředí má historii.} Následující stav nezávisí jen na poslední akci, ale na
  celé historii běhu.
\item \textbf{Prostředí je nedeterministické.} Výsledkem je celá množina stavů a~dopředu nevíme,
  který z~nich nastane.
\end{itemize}
Je-li $\tau(r) = \emptyset$, běh skončil. Dále předpokládáme, že všechny běhy skončí.

\textbf{Prostředí} je trojice $\mathit{Env} = \langle E, e_0, \tau\rangle$, kde $e_0$ je počáteční stav.
\end{defbox}

\begin{defbox}[Agent a~systém]
\textbf{Agent} je funkce zobrazující běhy končící stavem prostředí na akce:
\[
Ag : R^E \to A .
\]
Agent tedy rozhoduje na základě \emph{celé historie} systému a~rozhoduje se \emph{deterministicky}.
Označme $\mathcal{AG}$ množinu všech agentů.

\textbf{Systém} je dvojice $\langle Ag, \mathit{Env}\rangle$. Posloupnost
$(e_0, a_0, e_1, a_1, \dots)$ je \emph{běh agenta $Ag$ v~prostředí $\mathit{Env}$}, právě když
\begin{enumerate}
\item $e_0$ je počáteční stav $\mathit{Env}$,
\item $a_0 = Ag(e_0)$,
\item pro každé $u > 0$ platí $e_u \in \tau\big((e_0,a_0,\dots,a_{u-1})\big)$
  a~$a_u = Ag\big((e_0,a_0,\dots,e_u)\big)$.
\end{enumerate}
Množinu všech běhů agenta $Ag$ v~$\mathit{Env}$ značíme $R(Ag,\mathit{Env})$.
\end{defbox}

\begin{defbox}[Funkční ekvivalence agentů]
Agenti $Ag_1$ a~$Ag_2$ jsou \textbf{funkčně ekvivalentní vzhledem k~$\mathit{Env}$}, právě když
$R(Ag_1,\mathit{Env}) = R(Ag_2,\mathit{Env})$.

Jsou \textbf{prostě funkčně ekvivalentní} (\emph{simply functionally equivalent}), pokud jsou
funkčně ekvivalentní vzhledem ke \emph{všem} prostředím.
\end{defbox}

Funkční ekvivalence slouží ke srovnávání \emph{vyjadřovací síly} architektur. Říká totiž, že na
vnitřní implementaci nezáleží, pokud vede ke stejným pozorovatelným běhům.

\subsection{Čistě reaktivní agent a~agent s~vnitřním stavem}

Obecná definice agenta počítá s~celou historií běhu, a~to je hodně silný předpoklad. Nabízejí se
proto dvě úspornější varianty. U~obou se ptáme na totéž: kolik vyjadřovací síly ta úspora stojí.

\begin{defbox}[Čistě reaktivní agent (\emph{purely reactive / tropistic agent})]
\[
Ag : E \to A
\]
Reaguje pouze na aktuální stav prostředí, historii ignoruje.

\textbf{Vztah k~obecnému agentovi:} ke každému čistě reaktivnímu agentovi existuje funkčně
ekvivalentní standardní agent; opačně to \emph{neplatí}.

\emph{Termostat:} $Ag(e) = \mathrm{off}$ pro $e = $ \uv{teplota OK}, jinak $Ag(e) = \mathrm{on}$.
\end{defbox}

\begin{defbox}[Agent s~vnitřním stavem]
Nechť $I$ je množina vnitřních stavů agenta a~$\mathit{Per}$ množina jeho vjemů. Agent je určen trojicí funkcí:
\begin{align*}
\mathit{see} &: E \to \mathit{Per} && \text{(vnímání)}\\
\mathit{next} &: I \times \mathit{Per} \to I && \text{(aktualizace vnitřního stavu)}\\
\mathit{action} &: I \to A && \text{(výběr akce)}
\end{align*}
\textbf{Cyklus:} agent startuje ve stavu $i_0$; vnímá prostředí ($\mathit{see}(e)$); aktualizuje
stav ($i \leftarrow \mathit{next}(i, \mathit{see}(e))$); vybere akci ($a \leftarrow \mathit{action}(i)$);
vykoná ji a~opakuje.

\textbf{Věta (bez důkazu):} Ke každému agentovi s~vnitřním stavem existuje funkčně ekvivalentní
standardní agent. Vnitřní stav tedy \emph{nezvyšuje vyjadřovací sílu}, ale je zásadní pro
\emph{efektivitu} a~přehlednost implementace.
\end{defbox}

Funkce $\mathit{see}$ přitom zároveň formalizuje \textbf{pozorovatelnost}. Zobrazí-li dva různé
stavy $e_1 \ne e_2$ na tentýž vjem, agent je od sebe nedokáže odlišit a~prostředí je pro něj jen
částečně pozorovatelné. Krajní případy jsou dva: buď platí $|\mathit{Per}| = |E|$ a~funkce
$\mathit{see}$ je prostá, což znamená plnou pozorovatelnost, nebo je $|\mathit{Per}| = 1$
a~agent nevidí vůbec nic.

\subsection{Utilita: jak agentovi říct, co má dělat}

Agenty s~natvrdo zadrátovaným chováním nechceme. Chceme agentovi říct, \emph{co} má dělat,
nikoli \emph{jak}. Standardní postup UI je definovat cíl \emph{nepřímo}, totiž mírou úspěšnosti.

\begin{defbox}[Utilita]
\textbf{Utilita stavu:} $u : E \to \mathbb{R}$. Cílem agenta je dostat se do stavů s~vysokou
hodnotou utility.

Hodnocení jednotlivých stavů je krátkozraké. Proto zavádíme \textbf{utilitu běhu}:
\[
u : R \to \mathbb{R} .
\]
Ta lépe odpovídá agentům, kteří běží dlouhodobě. Utilitu běhu lze odvodit ze stavů např.\ jako
utilitu \emph{nejhoršího} navštíveného stavu nebo jako \emph{průměrnou} utilitu navštívených
stavů; co je lepší, závisí na úloze.
\end{defbox}

\begin{defbox}[Optimální agent]
Nechť $P(r \mid Ag, \mathit{Env})$ je pravděpodobnost běhu $r$ agenta $Ag$ v~prostředí $\mathit{Env}$,
přičemž $\sum_{r \in R(Ag,\mathit{Env})} P(r \mid Ag,\mathit{Env}) = 1$. Optimální agent maximalizuje
\emph{očekávanou utilitu}:
\[
Ag_{\mathrm{opt}} = \arg\max_{Ag \in \mathcal{AG}} \sum_{r \in R(Ag,\mathit{Env})}
u(r)\, P(r \mid Ag, \mathit{Env}).
\]
\end{defbox}

Definice je elegantní, má ale dvě slabiny. Především je \textbf{nekonstruktivní}: neříká,
jak takového agenta navrhnout. Kromě toho bývá obtížné utilitní funkci vůbec zadat. Poznamenejme,
že peníze nejsou pro člověka dobrou utilitou, protože jejich utilita je nelineární a~liší se pro
bohaté a~chudé.

\medskip
\noindent\textbf{Příklad: Tileworld.} Tileworld je klasické testovací prostředí pro agentní
architektury. Agenti se pohybují po šachovnici do čtyř směrů, na desce jsou díry, překážky
a~bloky a~cílem je zatlačit bloky do děr. Prostředí je \emph{dynamické}, protože díry, překážky
i~bloky náhodně vznikají a~zanikají. Utilita běhu se měří jako
\[
u(r) = \frac{N_f}{N_a},
\]
kde $N_f$ je počet děr zaplněných agentem během běhu $r$ a~$N_a$ počet všech děr, které se
během $r$ objevily. Agent tu musí zvládnout obojí: \emph{reagovat} na změny, když mu zmizí blok,
který právě tlačí, a~zároveň \emph{využívat příležitosti}, když se vedle něj objeví blok nový.
Tileworld se historicky používal právě ke zkoumání kompromisu mezi reaktivitou a~proaktivitou
(viz odst.~\ref{ssec:commitment}).

\subsection{Predikátová specifikace úlohy}

Utilita nemusí být jemně odstupňovaná. Speciálním případem je zobrazení do $\{0,1\}$, kdy se
hodnocení scvrkne na jediné rozhodnutí: běh je buď úspěšný ($u(r)=1$), nebo ne. Tím se popis
úlohy zjednoduší na logickou podmínku.

\begin{defbox}[Prostředí úlohy (\emph{task environment})]
Nechť $\Psi : R \to \{0,1\}$ je \textbf{predikátová specifikace}: $\Psi(r)$ platí, právě když
$u(r) = 1$. \textbf{Prostředí úlohy} je dvojice $\langle \mathit{Env}, \Psi \rangle$; specifikuje
jak vlastnosti systému (přes $\mathit{Env}$), tak kritérium splnění úlohy (přes $\Psi$).

Označme $R_\Psi(Ag,\mathit{Env}) = \{ r \in R(Ag,\mathit{Env}) \mid \Psi(r) \}$.
\end{defbox}

Zbývá dohodnout se, kdy vlastně agent úlohu \emph{řeší}. Existují tři výklady a~liší se
přísností.
\begin{itemize}
\item \textbf{Pesimistický výklad} žádá, aby podmínku splnily \emph{všechny} běhy, tedy
  $R_\Psi(Ag,\mathit{Env}) = R(Ag,\mathit{Env})$. Agent úlohu řeší jen tehdy, když nemůže
  selhat ani v~jednom běhu.
\item \textbf{Optimistický výklad} se spokojí s~tím, že $\Psi$ splní \emph{alespoň jeden} běh:
  $\exists r \in R(Ag,\mathit{Env}): \Psi(r)$.
\item \textbf{Realistický výklad} rozšíří $\tau$ o~rozdělení pravděpodobnosti a~měří úspěch
  pravděpodobností
  $P(\Psi \mid Ag,\mathit{Env}) = \sum_{r \in R_\Psi(Ag,\mathit{Env})} P(r \mid Ag,\mathit{Env})$.
\end{itemize}

\begin{defbox}[Typy úloh]
\textbf{Úlohy dosažení} (\emph{achievement tasks}) --- \uv{něco najít}. Je dána množina cílových
stavů $G \subseteq E$; $\Psi(r)$ platí, pokud se v~běhu $r$ vyskytne alespoň jeden stav z~$G$.
Typicky každá klasická úloha UI (prohledávání).

\textbf{Úlohy udržení} (\emph{maintenance tasks}) --- \uv{něčemu se vyhnout}. Je dána množina
\uv{špatných} stavů $B \subseteq E$; $\Psi(r)$ neplatí, pokud se v~$r$ vyskytne stav z~$B$.
Typicky hry, kde $B$ jsou prohrané pozice a~prostředím je soupeř.

\textbf{Kombinace:} dosáhnout stavu z~$G$ a~přitom se vyhnout stavům z~$B$.
\end{defbox}

\subsection{Mechanismus výběru akcí --- přehled}
\label{ssec:vyber-akci}

\begin{defbox}[Mechanismus výběru akcí (\emph{action selection mechanism})]
Procedura, kterou agent v~každém okamžiku rozhoduje, kterou z~dostupných akcí provede,
na základě vjemů z~prostředí a~svého vnitřního stavu. \textbf{Architektura agenta} je
softwarová architektura, která takový mechanismus realizuje --- podle Maes (1991)
\uv{konkrétní metodologie stavby agentů, která určuje, jak lze agenta rozložit na množinu
komponent a~jak mají tyto komponenty interagovat, aby dohromady odpověděly na otázku, jak data
ze senzorů a~aktuální vnitřní stav určují akce a~budoucí vnitřní stav agenta.}
\end{defbox}

Výběr akce je \textbf{klíčový problém návrhu agenta}. Celá kapitola~\ref{sec:rizeni} proto není
nic jiného než přehled jeho realizací a~následující tabulka shrnuje, jak se dělí do rodin.

\begin{center}
\small
\begin{tabular}{p{0.20\textwidth} p{0.34\textwidth} p{0.38\textwidth}}
\toprule
\textbf{Rodina} & \textbf{Mechanismus výběru akce} & \textbf{Typický zástupce} \\
\midrule
symbolická / deliberativní
  & logická dedukce nad databází formulí; plánování
  & deduktivní agent, STRIPS, BDI, PRS \\
reaktivní (symbolická)
  & sada pravidel \uv{situace $\to$ akce} s~pevnou prioritou
  & subsumpční architektura (Brooks) \\
reaktivní (konekcionistická)
  & šíření aktivace v~síti modulů, vyhrává nejaktivnější
  & agent network architecture (Maes) \\
hybridní
  & vrstvy různé abstrakce + arbitráž mezi nimi
  & TouringMachines, InteRRaP \\
koaliční
  & soutěž dynamicky vznikajících koalic procesů o~\uv{vědomí}
  & IDA / globální pracovní prostor \\
naučená
  & maximalizace naučené $Q$-funkce / stochastická politika
  & Q-learning, actor-critic (kap.~\ref{sec:uceni}) \\
generativní
  & stav se zakóduje do promptu, výstup jazykového modelu se čte jako akce
  & LLM agenti (\S\ref{ssec:llm-agenti}) \\
\bottomrule
\end{tabular}
\end{center}

\subsection{Shrnutí ke~zkoušce}

\begin{itemize}
\item Agent je autonomní systém zasazený do prostředí, který v~něm vnímá a~jedná, aby splnil
  delegované cíle. Multiagentní systém je systém interagujících agentů, kteří \emph{nemusí
  sdílet společný cíl}.
\item Inteligentního agenta charakterizují tři vlastnosti: reaktivita, proaktivita a~sociálnost.
  Těžké přitom není dosáhnout jedné z~nich, ale vyvážit je.
\item Dennett zavádí intencionální systém a~tři postoje k~predikci chování: fyzikální, návrhový
  a~intencionální. Intencionální postoj je pro informatiku abstrakce pro zvládnutí složitosti.
\item Prostředí se třídí podle čtyř dichotomií: pozorovatelnost, statické nebo dynamické,
  determinismus a~diskrétnost. Reálná prostředí spadají vždy do nejtěžší varianty.
\item Formalismus tvoří stavy $E$, akce $A$, běh $r = e_0 a_0 e_1 \dots e_n$, jeho podmnožiny
  $R^A$ a~$R^E$, přechodová funkce $\tau : R^A \to 2^E$, z~níž plyne historie i~nedeterminismus,
  prostředí $\mathit{Env}=\langle E,e_0,\tau\rangle$, agent $Ag : R^E \to A$ a~systém
  $\langle Ag,\mathit{Env}\rangle$.
\item Funkční ekvivalence se definuje vzhledem k~danému prostředí, nebo jako prostá. Čistě
  reaktivní agent $Ag: E \to A$ je slabší než obecný agent, kdežto agent s~vnitřním stavem
  (funkce $\mathit{see}$, $\mathit{next}$, $\mathit{action}$) slabší není; oba se dají převést
  na standardního agenta.
\item Utilita se zavádí buď nad stavy jako $u : E \to \mathbb{R}$, nebo nad běhy jako
  $u : R \to \mathbb{R}$. Optimální agent maximalizuje \emph{očekávanou} utilitu, jenže tato
  definice je nekonstruktivní. Školním příkladem je Tileworld s~utilitou $u(r)=N_f/N_a$.
\item Prostředí úlohy je dvojice $\langle \mathit{Env}, \Psi\rangle$; kdy agent úlohu řeší, na
  to odpovídá pesimista, optimista a~realista. Úlohy se dělí na úlohy dosažení (množina $G$)
  a~úlohy udržení (množina $B$).
\item \textbf{Mechanismus výběru akcí} je jádrem celé otázky. U~zkoušky se očekává, že student
  vyjmenuje rodiny mechanismů i~jejich typické zástupce podle tabulky výše.
\end{itemize}

%=====================================================================

\section{Metody pro řízení agentů: symbolické a~konekcionistické reaktivní
plánování, hybridní přístupy}
\label{sec:rizeni}
%=====================================================================

Tato kapitola je jádrem celého okruhu. Probírá jednotlivé \emph{metody řízení agenta}, tedy
konkrétní realizace mechanismu výběru akcí, který zavádí odst.~\ref{ssec:vyber-akci}. Výklad
postupuje po ose od symbolického (deliberativního) pólu směrem k~pólu reaktivnímu a~končí
u~hybridních architektur, které se snaží oba póly spojit:

\begin{center}
\small
\begin{tabular}{@{}p{0.26\textwidth} p{0.66\textwidth}@{}}
\toprule
\textbf{Přístup} & \textbf{Kde je v~této kapitole} \\
\midrule
symbolické (deliberativní) řízení
  & deduktivní agent, BDI, STRIPS, PRS (odst.\ 2.1--2.6) \\
symbolické reaktivní plánování
  & subsumpční architektura (Brooks), odst.~\ref{ssec:subsumpce} \\
konekcionistické reaktivní plánování
  & agent network architecture (Maes), odst.~\ref{ssec:maes} \\
hybridní přístupy
  & vrstvené architektury, TouringMachines, InteRRaP, IDA, odst.~\ref{ssec:hybridni} a~dále \\
\bottomrule
\end{tabular}
\end{center}
\subsection{Klasický přístup UI a~deduktivní agent}

Klasický způsob, jakým UI staví \uv{inteligentní systém}, stojí na dvou pilířích:
\begin{enumerate}
\item \textbf{Symbolická reprezentace} prostředí a~chování, zde konkrétně reprezentace pomocí
  logických formulí.
\item \textbf{Syntaktická manipulace} s~touto reprezentací, tedy logická dedukce neboli
  dokazování vět.
\end{enumerate}
Teorie o~chování agenta je pak zároveň jeho programem. Je to \emph{spustitelná specifikace},
ze které se přímo generují konkrétní akce, a~právě v~tom je půvab celého přístupu. Naráží
ovšem na dva klasické problémy.
\begin{itemize}
\item \textbf{Problém transdukce} (\emph{transduction problem}) se ptá, jak reálný svět převést
  do symbolické reprezentace. Spadá sem počítačové vidění, zpracování přirozeného jazyka
  i~učení. Goethem řečeno: \uv{Šedá, milý příteli, je všechna teorie, a~zelený je zlatý strom
  života.}
\item \textbf{Problém reprezentace a~uvažování} se ptá, jak znalosti zapsat a~jak nad nimi
  potom efektivně odvozovat. Odpovědí jsou dokazovače a~plánovače.
\end{itemize}

\medskip
\noindent\textbf{Odbočka: dedukce vs.\ indukce.}
\emph{Dedukce} je odvození závěrů, které jsou pravdivé, jsou-li pravdivé předpoklady; postupuje
se při ní od obecného ke konkrétnímu. Školním příkladem je sylogismus \uv{Všichni lidé jsou
smrtelní, Sókratés je člověk, tedy Sókratés je smrtelný}. \emph{Indukce} naproti tomu znamená,
že při pravdivých předpokladech je závěr spíše pravdivý než nepravdivý, a~postupuje se při ní
od konkrétního k~obecnému. Sherlock Holmes tak ve skutečnosti používá indukci, i~když jí sám
říká dedukce. Matematická indukce je navzdory svému jménu naopak dedukce.

\begin{defbox}[Deduktivní agent (agent jako dokazovač vět)]
Nechť $L$ je množina formulí predikátové logiky prvního řádu a~$D = 2^L$ množina databází těchto
formulí. Vnitřní stav agenta tvoří databáze $\mathit{DB} \in D$. Uvažování je realizováno
odvozovacími pravidly $\rho$:
\[
\mathit{DB} \vdash_\rho \varphi \quad\text{---\ formuli } \varphi \text{ lze odvodit z } \mathit{DB}
\text{ pomocí pravidel } \rho .
\]
Agent je pak trojicí funkcí
$\mathit{see} : E \to \mathit{Per}$, $\mathit{next} : D \times \mathit{Per} \to D$,
$\mathit{action} : D \to A$, přičemž $\mathit{action}$ je definována procedurou níže.
\end{defbox}

\noindent\textbf{Výběr akce jako dokazování.}
\begin{algorithmic}[1]
\Function{action}{$\mathit{DB} : D$} : $A$
  \ForAll{$a \in A$}
    \If{$\mathit{DB} \vdash_\rho \mathit{Do}(a)$} \Return $a$ \EndIf
  \EndFor
  \ForAll{$a \in A$}
    \If{$\mathit{DB} \not\vdash_\rho \neg \mathit{Do}(a)$} \Return $a$ \EndIf
  \EndFor
  \State \Return $\mathit{null}$
\EndFunction
\end{algorithmic}

První cyklus hledá akci, kterou lze \emph{dokázat} jako žádoucí, tedy takovou, pro kterou se
z~databáze odvodí $\mathit{Do}(a)$. Pokud žádná taková akce neexistuje, agent ze svého nároku
sleví: druhý cyklus se spokojí s~akcí, která je s~databází alespoň \emph{konzistentní}, protože
o~ní nelze dokázat, že by se dělat neměla. Selže-li i~to, agent nedělá nic.

\medskip
\noindent\textbf{Příklad: vysavačový svět 3$\times$3.}
Svět popisují tři predikáty: $\mathit{In}(x,y)$ říká, že agent je na pozici $(x,y)$,
$\mathit{Dirt}(x,y)$, že na $(x,y)$ je špína, a~$\mathit{Facing}(d)$, že agent je otočen směrem
$d$. Akce se pak odvozují z~pravidel:
\begin{align*}
\mathit{In}(x,y) \wedge \mathit{Dirt}(x,y) &\Rightarrow \mathit{Do}(\mathit{suck}) \\
\mathit{In}(0,0) \wedge \mathit{Facing}(\mathit{north}) \wedge \neg\mathit{Dirt}(0,0)
  &\Rightarrow \mathit{Do}(\mathit{forward}) \\
\mathit{In}(0,2) \wedge \mathit{Facing}(\mathit{north}) \wedge \neg\mathit{Dirt}(0,2)
  &\Rightarrow \mathit{Do}(\mathit{turn}) \qquad \dots
\end{align*}
Agent tak systematicky projde celý svět ve tvaru \uv{hada} 00--01--02--12--11--10--\dots
Stojí za povšimnutí, že úklid má vyšší prioritu než pohyb: pravidlo pro $\mathit{suck}$ je
uvedeno jako první a~všechna ostatní pravidla mají v~předpokladu negaci $\mathit{Dirt}$.

\medskip
\noindent\textbf{Klady a~zápory.}
\begin{itemize}
\item[$+$] Přístup je elegantní a~má jasnou (a~krásnou) sémantiku; specifikace je zároveň programem.
\item[$-$] Pro netriviální úlohy je prakticky nepoužitelný, protože dokazování je pomalé a~nemusí
  vůbec skončit.
\item[$-$] \textbf{Problém vypočitatelné racionality} (\emph{calculative rationality}): agent
  odvodí optimální akci pro ten stav prostředí, jaký byl \emph{na začátku} odvozování. Prostředí
  se ale mezitím změnilo a~vypočtená akce už optimální není. V~rychle se měnícím prostředí je to
  katastrofa.
\item[$-$] Obtížně se navrhuje funkce $\mathit{see}$: není zřejmé, jak z~obrázku udělat formule
  ani jak reprezentovat čas.
\end{itemize}
Přesto jednotlivé prvky tohoto přístupu žijí dál v~pozdějších architekturách, nejvýrazněji v~BDI.

\subsection{Praktické uvažování: architektura BDI}

Deduktivní agent ztroskotal na tom, že se snažil o~všem rozhodovat dokazováním. Architektura BDI
proto vychází z~úplně jiné inspirace, totiž z~toho, jak se rozhoduje člověk.

\begin{defbox}[Praktické uvažování (\emph{practical reasoning})]
Je inspirováno lidským rozhodováním. \emph{Teoretické} uvažování (Sókratés) vede jen k~tomu,
co si o~světě myslíme, kdežto \emph{praktické} uvažování vede k~\textbf{akcím}. Probíhá ve
dvou fázích:
\begin{itemize}
\item \textbf{Deliberace} (\emph{deliberation}) rozhoduje o~tom, \emph{čeho chceme dosáhnout}
  (\uv{chci dodělat magisterské studium}).
\item \textbf{Uvažování o~prostředcích} (\emph{means-ends reasoning}), tedy plánování, rozhoduje
  o~tom, \emph{jak toho dosáhneme} (vytvořit plán, jak dostudovat).
\end{itemize}
A~obojí nesmí trvat příliš dlouho.
\end{defbox}

\begin{defbox}[Beliefs --- Desires --- Intentions]
\textbf{Beliefs} (přesvědčení) tvoří informační stav agenta, tedy jeho znalosti o~světě.
V~MAS jim záměrně neříkáme \uv{znalosti}, abychom zdůraznili, že jsou \emph{subjektivní},
že \emph{nemusí být pravdivé} a~že se \emph{mohou v~budoucnu změnit}. Mohou obsahovat
i~odvozovací pravidla umožňující dopředné řetězení.

\textbf{Desires} (touhy) tvoří motivační stav agenta. Jsou to cíle nebo situace, kterých by
agent chtěl dosáhnout, a~na rozdíl od záměrů \emph{mohou být vzájemně nekonzistentní}
(chci se učit i~jít na pivo).

\textbf{Intentions} (záměry) jsou stavy světa, kterých se agent \emph{zavázal} dosáhnout.
Právě záměry vedou k~akcím, protože agent kvůli nim plánuje. Záměry \textbf{přetrvávají},
dokud nenastane jedna ze tří věcí:
\begin{enumerate}
\item agent jich dosáhne, nebo
\item začne věřit, že jsou nedosažitelné, nebo
\item pominou důvody, které k~nim vedly.
\end{enumerate}
Kromě toho záměry \emph{omezují další deliberaci}, protože agent už znovu nezvažuje možnosti
neslučitelné s~přijatým záměrem, a~\emph{ovlivňují budoucí přesvědčení}, protože agent počítá
s~tím, že se mu záměr podaří.
\end{defbox}

\noindent\textbf{Rozdíl mezi touhou a~záměrem} ilustruje Bratman (1990):
\uv{Moje touha zahrát si dnes odpoledne basketbal je pouze potenciálním vlivem na moje odpolední
chování. Musí soupeřit s~mými dalšími relevantními touhami\dots{} Naproti tomu jakmile
\emph{zamýšlím} hrát dnes odpoledne basketbal, je věc rozhodnuta: obvykle už nemusím dál vážit
pro a~proti; až odpoledne přijde, prostě svůj záměr vykonám.}

\begin{defbox}[Funkce deliberace]
Označme $B \subseteq \mathit{Bel}$ aktuální přesvědčení, $D \subseteq \mathit{Des}$ aktuální touhy
a~$I \subseteq \mathit{Int}$ aktuální záměry agenta. Deliberace se skládá ze tří funkcí:
\begin{itemize}
\item $\mathit{brf} : 2^{\mathit{Bel}} \times \mathit{Per} \to 2^{\mathit{Bel}}$ je
  \emph{belief revision function}, která aktualizuje přesvědčení podle vjemu;
\item $\mathit{options} : 2^{\mathit{Bel}} \times 2^{\mathit{Int}} \to 2^{\mathit{Des}}$
  generuje možnosti (touhy) z~přesvědčení a~ze stávajících záměrů;
\item $\mathit{filter} : 2^{\mathit{Bel}} \times 2^{\mathit{Des}} \times 2^{\mathit{Int}}
  \to 2^{\mathit{Int}}$ z~nabídnutých možností vybírá záměry.
\end{itemize}
Rozdělení deliberace na \emph{generování možností} a~\emph{filtrování} je klíčové: generování
je levné a~široké, zatímco filtrování je vlastní rozhodování, které bere v~úvahu i~stávající
závazky.
\end{defbox}

\subsection{Plánování: STRIPS}

Druhou fází praktického uvažování je plánování. Zatímco deliberace vybere cíl, plánování hledá
cestu k~němu.

\begin{defbox}[Uvažování o~prostředcích / plánování]
Je to proces rozhodování, jak dosáhnout cíle (záměru) pomocí dostupných prostředků (akcí).
\begin{itemize}
\item \textbf{Vstup:} cíl (= záměr $I$); aktuální stav prostředí (= přesvědčení $B$);
  akce dostupné agentovi ($\mathit{Ac}$).
\item \textbf{Výstup:} plán, tedy posloupnost akcí, po jejímž provedení je cíl splněn.
\end{itemize}
\end{defbox}

\textbf{STRIPS} (Fikes, Nilsson, 1971) modeluje svět množinou formulí logiky prvního řádu.
Každá akce je popsána \emph{schématem} s~předpoklady a~efekty, tedy s~přidávanými a~mazanými
fakty. Plánovací algoritmus najde rozdíl mezi cílem a~aktuálním stavem a~vhodně zvolenou akcí
ho zmenší; tomuto postupu se říká \emph{means-ends analysis}. Je to hezké, ale ne příliš
praktické, protože algoritmus se často utápí v~detailech nízké úrovně.

\begin{defbox}[Deskriptor akce, plánovací problém, plán]
\textbf{Deskriptor akce} $a$ je trojice $\langle P_a, D_a, A_a\rangle$, kde
$P_a$ jsou předpoklady (\emph{preconditions}), $D_a$ mazaná množina (\emph{delete list})
a~$A_a$ přidávaná množina (\emph{add list}); pro jednoduchost jde o~uzemněné atomické formule
(bez spojek a~proměnných).

\textbf{Plánovací problém} je trojice $\langle B_0, O, G\rangle$, kde $B_0$ jsou počáteční
přesvědčení, $O = \{\langle P_a,D_a,A_a\rangle : a \in \mathit{Ac}\}$ deskriptory všech akcí
a~$G$ množina formulí popisujících cíl.

\textbf{Plán} $\pi = (a_1,\dots,a_n)$, $a_i \in \mathit{Ac}$, určuje posloupnost databází
\[
B_i = (B_{i-1} \setminus D_{a_i}) \cup A_{a_i}, \qquad i = 1,\dots,n .
\]
Plán je \textbf{přípustný} (\emph{admissible}), pokud $B_{i-1} \models P_{a_i}$ pro každé $i$.
Plán je \textbf{korektní/validní} (\emph{correct}), pokud je přípustný a~navíc $B_n \models G$.

\textbf{Plánování} znamená pro zadaný $\langle B_0,O,G\rangle$ najít korektní plán, nebo prohlásit,
že žádný neexistuje.
\end{defbox}

\medskip
\noindent\textbf{Příklad: svět kostek (\emph{blocks world}).}
Situaci popisují predikáty $\mathit{On}(x,y)$ ($x$ je na $y$), $\mathit{OnTable}(x)$,
$\mathit{Clear}(x)$ (na $x$ nic není), $\mathit{Holding}(x)$ (rameno drží $x$)
a~$\mathit{ArmEmpty}$. Agent má k~dispozici čtyři akce:
\begin{center}
\small
\begin{tabular}{p{0.20\textwidth} p{0.36\textwidth} p{0.36\textwidth}}
\toprule
\textbf{Akce} & \textbf{Pre / Del} & \textbf{Add} \\
\midrule
$\mathit{Stack}(x,y)$
  & Pre: $\{\mathit{Clear}(y), \mathit{Holding}(x)\}$ \newline
    Del: $\{\mathit{Clear}(y), \mathit{Holding}(x)\}$
  & $\{\mathit{ArmEmpty}, \mathit{On}(x,y)\}$ \\
$\mathit{UnStack}(x,y)$
  & Pre: $\{\mathit{On}(x,y),$ $\mathit{Clear}(x),\mathit{ArmEmpty}\}$ \newline
    Del: $\{\mathit{On}(x,y),\mathit{ArmEmpty}\}$
  & $\{\mathit{Holding}(x), \mathit{Clear}(y)\}$ \\
$\mathit{Pickup}(x)$
  & Pre: $\{\mathit{OnTable}(x),$ $\mathit{Clear}(x),\mathit{ArmEmpty}\}$ \newline
    Del: $\{\mathit{OnTable}(x),\mathit{ArmEmpty}\}$
  & $\{\mathit{Holding}(x)\}$ \\
$\mathit{PutDown}(x)$
  & Pre: $\{\mathit{Holding}(x)\}$ \newline
    Del: $\{\mathit{Holding}(x)\}$
  & $\{\mathit{ArmEmpty}, \mathit{OnTable}(x)\}$ \\
\bottomrule
\end{tabular}
\end{center}
Vyjdeme z~počátečního stavu
\[
B_0 = \{\mathit{Clear}(A), \mathit{On}(A,B), \mathit{OnTable}(B),
\mathit{OnTable}(C), \mathit{Clear}(C), \mathit{ArmEmpty}\}
\]
a~chceme dosáhnout cíle $G = \{\mathit{OnTable}(A), \mathit{OnTable}(B), \mathit{OnTable}(C)\}$.

\noindent Korektní plán je $\pi = (\mathit{UnStack}(A,B),\ \mathit{PutDown}(A))$. Ověřit se dá
krok za krokem:
\begin{itemize}
\item Platí $B_0 \models \{\mathit{On}(A,B),\mathit{Clear}(A),\mathit{ArmEmpty}\}$, takže první
  akce je přípustná; po jejím provedení dostáváme
  \[
  B_1 = \big(B_0 \setminus \{\mathit{On}(A,B),\mathit{ArmEmpty}\}\big) \cup
  \{\mathit{Holding}(A),\mathit{Clear}(B)\}.
  \]
\item Platí $B_1 \models \{\mathit{Holding}(A)\}$, takže i~druhá akce je přípustná; po ní
  dostáváme
  \[
  B_2 = \big(B_1 \setminus \{\mathit{Holding}(A)\}\big) \cup
  \{\mathit{ArmEmpty},\mathit{OnTable}(A)\}.
  \]
\item Konečně $B_2 \models G$, takže plán je korektní.
\end{itemize}

\begin{defbox}[Pomocné funkce nad plány]
$\mathit{Plan}$ je množina všech plánů nad $\mathit{Ac}$;
$\mathit{pre}(\pi)$ je předpoklad plánu; $\mathit{body}(\pi)$ tělo plánu, tedy posloupnost akcí;
$\mathit{empty}(\pi)$ se ptá, zda je plán prázdný;
$\mathit{execute}(\pi)$ provede všechny akce plánu;
$\mathit{hd}(\pi)$ vrací první akci a~$\mathit{tail}(\pi)$ zbytek;
$\mathit{sound}(\pi,I,B)$ říká, že plán $\pi$ je korektní vzhledem k~záměrům $I$
a~přesvědčením $B$.

Plánovací funkce agenta má tvar
$\mathit{plan} : 2^{\mathit{Bel}} \times 2^{\mathit{Int}} \times 2^{\mathit{Ac}} \to \mathit{Plan}$.
\end{defbox}

\textbf{Knihovna plánů.} Konstruovat plány online je drahé, a~agent to naštěstí dělat nemusí.
Velmi často je proto $\mathit{plan}()$ implementována jako \emph{knihovna hotových plánů}:
stačí ji jednou projít a~zkontrolovat, zda předpoklady plánu odpovídají aktuálním přesvědčením
a~zda efekty plánu odpovídají cíli. Tento nápad je jádrem architektury PRS (sekce~\ref{ssec:prs}).

\subsection{Implementace BDI agenta}

Jednotlivé díly už máme pohromadě: revizi přesvědčení, deliberaci i~plánování. Teď je stačí složit
do jedné řídicí smyčky.

\noindent\textbf{Hlavní řídicí smyčka BDI agenta:}
\begin{algorithmic}[1]
\State $B \gets B_0;\quad I \gets I_0$
\While{$\mathbf{true}$}
  \State $v \gets \mathit{see}()$
  \State $B \gets \mathit{brf}(B,v)$
  \State $D \gets \mathit{options}(B,I)$
  \State $I \gets \mathit{filter}(B,D,I)$
  \State $\pi \gets \mathit{plan}(B,I,\mathit{Ac})$
  \While{$\neg\big(\mathit{empty}(\pi) \vee \mathit{succeeded}(I,B) \vee \mathit{impossible}(I,B)\big)$}
    \State $a \gets \mathit{hd}(\pi);\quad \mathit{execute}(a);\quad \pi \gets \mathit{tail}(\pi)$
    \State $v \gets \mathit{see}();\quad B \gets \mathit{brf}(B,v)$
    \If{$\mathit{reconsider}(I,B)$}
      \State $D \gets \mathit{options}(B,I);\quad I \gets \mathit{filter}(B,D,I)$
    \EndIf
    \If{$\neg\,\mathit{sound}(\pi,I,B)$}
      \State $\pi \gets \mathit{plan}(B,I,\mathit{Ac})$ \Comment{přeplánování}
    \EndIf
  \EndWhile
\EndWhile
\end{algorithmic}

Dva predikáty ve~vnitřním cyklu si zaslouží vysvětlení. Predikát $\mathit{succeeded}(I,B)$ říká,
že záměry $I$ platí za předpokladu přesvědčení $B$, kdežto $\mathit{impossible}(I,B)$ říká, že
záměry $I$ za předpokladu $B$ platit nemohou. Vnitřní cyklus tedy skončí ve třech případech:
když je plán vyčerpán, když je cíl splněn, anebo když se cíl stal nedosažitelným.

\subsection{Odhodlání agenta (\emph{commitment})}
\label{ssec:commitment}

\begin{defbox}[Strategie odhodlání vůči záměrům]
Kdy má agent opustit cíl? Nabízejí se tři odpovědi:
\begin{itemize}
\item \textbf{Slepé odhodlání} (\emph{blind commitment}, též \emph{fanatical}) drží záměr tak
  dlouho, dokud agent \emph{nezačne věřit, že ho dosáhl}. Takový agent se nikdy nevzdá.
\item \textbf{Jednosměrné odhodlání} (\emph{single-minded commitment}) drží záměr, dokud agent
  nezačne věřit, že ho dosáhl, \emph{nebo} že už není dosažitelný.
\item \textbf{Otevřené odhodlání} (\emph{open-minded commitment}) nechává záměr přetrvávat tak
  dlouho, dokud agent věří, že je stále dosažitelný a~stále chtěný. Záměr proto může být opuštěn
  i~při pouhé změně motivace.
\end{itemize}
\end{defbox}

Odlišit je ale potřeba \textbf{odhodlání vůči plánu} od \textbf{odhodlání vůči cíli}. Agent
v~algoritmu výše je odhodlán vůči jednomu plánu a~opustí ho tehdy, když věří, že cíl je splněn, že
je nedosažitelný, anebo když je plán prázdný. Zbývá tím ovšem ta těžší otázka: kdy má agent
\emph{přehodnotit sám záměr}?

Klasické dilema zní, že deliberace \textbf{trvá čas} a~prostředí se během ní mění. Obě krajní
polohy jsou špatně.
\begin{itemize}
\item Agent, který záměry \emph{nepřehodnocuje}, se může upnout k~cílům, jichž už nelze dosáhnout.
\item Agent, který přehodnocuje \emph{příliš často}, je zase tak zaneprázdněn uvažováním, že nic
  nestihne.
\end{itemize}

Východiskem je \textbf{metařízení} (Wooldridge, Parsons, 1999). Zavádí se funkce
$\mathit{reconsider}()$, která je \emph{výpočetně jednodušší} než samotné přehodnocení a~jen
odhaduje, zda se přehodnocení vyplatí. Kritérium kvality je přirozené: $\mathit{reconsider}()$
funguje dobře, pokud pokaždé, když navrhne přehodnocení, agent po deliberaci \emph{skutečně}
záměr změní, a~naopak.

\begin{defbox}[Odvážný vs.\ opatrný agent]
\textbf{Odvážný agent} (\emph{bold agent}) přehodnocuje záměry teprve po dokončení celého plánu;
kvůli deliberaci plán nikdy nepřeruší.

\textbf{Opatrný agent} (\emph{cautious agent}) přehodnocuje po každé akci plánu.

\textbf{Míra odvahy} (\emph{level of boldness}) udává, kolik akcí agent provede mezi dvěma
deliberacemi.

\textbf{Dynamičnost prostředí} (\emph{rate of world change}) udává, kolikrát se prostředí změní
během jednoho cyklu agenta.

\textbf{Efektivita agenta} $=$ počet dosažených záměrů $/$ počet všech záměrů.
\end{defbox}

\textbf{Experimentální výsledek (Kinny, Georgeff, v~Tileworldu):} který z~obou typů agentů vyhraje,
závisí na prostředí.
\begin{itemize}
\item Mění-li se svět \emph{pomalu}, jsou efektivnější \textbf{odvážní} agenti, protože deliberace
  by pro ně byla jen zbytečnou režií.
\item Mění-li se svět \emph{rychle}, překonávají je \textbf{opatrní} agenti, protože plány rychle
  zastarávají.
\end{itemize}
Je to konkrétní kvantitativní podoba obecného kompromisu reaktivita vs.\ proaktivita.

\subsection{Procedurální odvozovací systém (PRS)}
\label{ssec:prs}

\begin{defbox}[PRS --- \emph{Procedural Reasoning System}]
Vznikl v~80.~letech na Stanfordu (Georgeff a~kol.). Jde o~\textbf{první implementaci architektury
BDI} a~zároveň o~jednu z~prakticky nejúspěšnějších agentních architektur vůbec; byla mnohokrát
reimplementována jako
\textbf{AgentSpeak/Jason}, \textbf{JAM}, \textbf{JACK} nebo \textbf{Jadex}.

Z~praktických nasazení jmenujme OASIS (systém řízení letového provozu v~Sydney), SPOC (organizace
podnikových procesů) a~SWARMM (letecký simulátor australského letectva).
\end{defbox}

\textbf{Komponenty PRS agenta.} Agent se skládá z~databáze \emph{beliefs} (atomy FOL prologovského
typu), z~\emph{desires} (cíle), z~\emph{knihovny plánů} (\emph{plan library}), ze zásobníku
\emph{intentions} a~z~\emph{interpretu}, který vše propojuje. Data ze senzorů vstupují do beliefs
a~interpret na druhém konci vydává akce.

\begin{defbox}[Plán v~PRS]
Agent má knihovnu hotových plánů, které reprezentují jeho \emph{procedurální znalosti}.
\textbf{Neplánuje se od nuly}, pouze se z~knihovny \emph{vybírá}. Plán se skládá ze tří částí:
\begin{itemize}
\item \textbf{Goal} (též \emph{invocation condition} či \emph{cue}) je podmínka, která má platit
  po vykonání plánu; určuje, na jaký cíl se plán hodí.
\item \textbf{Context} je podmínka nutná pro spuštění plánu a~musí být konzistentní
  s~přesvědčeními.
\item \textbf{Body} jsou akce, které se mají vykonat.
\end{itemize}
Tělo plánu \emph{není} jen lineární posloupnost akcí, protože může obsahovat i~další \emph{cíle}
tvaru \uv{dosáhni $f$}, \uv{dosáhni $f$ nebo $g$} nebo \uv{dosahuj $f$, dokud nenastane $g$}.
Tím vzniká hierarchické rozkládání cílů na podcíle.
\end{defbox}

\textbf{Běh interpretu} probíhá takto:
\begin{enumerate}
\item Nejprve se agent inicializuje počátečními přesvědčeními $B_0$ a~cílem nejvyšší úrovně
  (\emph{top-level goal}).
\item Zásobník záměrů obsahuje aktuální cíle v~různých stavech rozpracovanosti.
\item Interpret vezme záměr z~vrcholu zásobníku a~v~knihovně plánů hledá plány s~odpovídajícím
  cílem.
\item Z~nalezených plánů vybere ty, jejichž \emph{kontext} je konzistentní s~aktuálními
  přesvědčeními; ty tvoří aktuální \emph{options} neboli \emph{desires}.
\item \textbf{Deliberace} je pak výběr jednoho záměru z~těchto možností. Nabízejí se dvě cesty:
  \begin{itemize}
  \item Původní PRS používal \emph{metaplány}, tedy plány o~plánech, které modifikovaly záměry
    agenta. Ukázalo se to jako příliš složité.
  \item Praktičtější varianta dává každému plánu číselné \emph{ohodnocení} (očekávanou utilitu)
    a~vybere plán s~nejvyšší hodnotou.
  \end{itemize}
\item Vybraný plán se vykoná, což může vést k~přidání dalších záměrů na zásobník.
\item Selže-li plán, agent zvolí jiný záměr z~možností a~pokračuje.
\end{enumerate}

\medskip
\noindent\textbf{Ukázka --- Jason/AgentSpeak} (agent nosí majiteli pivo z~ledničky):
\begin{quote}\small\ttfamily
\noindent available(beer,fridge).\ \ limit(beer,10).\ \ \ \ /* počáteční beliefs */\\
too\_much(B) :- .count(consumed(\_,\_,\_,B),Q) \& limit(B,L) \& Q > L.\ \ /* pravidlo */\\[0.3em]
+!has(owner,beer)\ \ \ \ \ \ \ \ \ \ \ \ \ \ \ \ \ \ \ \ \ \ \ /* spouštěcí událost */\\
\hspace*{1em}: available(beer,fridge) \& not too\_much(beer)\ \ \ \ /* kontext */\\
\hspace*{1em}<- !at(robot,fridge); open(fridge); get(beer); close(fridge);\ \ /* tělo */\\
\hspace*{2em} !at(robot,owner); hand\_in(beer).\\[0.3em]
+!at(robot,P) : not at(robot,P) <- move\_towards(P); !at(robot,P).
\end{quote}
Zápis \texttt{+!g : c <- b} se čte: \uv{nastane-li událost přidání cíle \texttt{g} a~platí-li
kontext \texttt{c}, vykonej tělo \texttt{b}}. Symbol \texttt{!} označuje \emph{achievement goal}
(dosáhni), symboly \texttt{+} a~\texttt{-} přidání a~odebrání přesvědčení nebo cíle
a~\texttt{.send} a~\texttt{.count} jsou vestavěné akce.

\subsection{Reaktivní přístup a~jeho východiska}

V~80.\ a~90.\ letech se ukázalo, že drobnými úpravami symbolického přístupu se jeho potíže
odstranit nedají. Vznikla proto alternativní paradigmata UI, která se od symbolické tradice
neodchýlila v~detailu, ale rovnou v~základu. Jejich společným jmenovatelem je
\textbf{odmítnutí symbolické reprezentace} i~uvažování postaveného na syntaktické manipulaci
s~ní. Opřít se dají o~tři klíčové teze.

\begin{defbox}[Tři principy reaktivního přístupu]
\textbf{Interakce} (\emph{situatedness}). Inteligentní chování závisí na prostředí, ve kterém
je agent \emph{zasazen}, a~nelze ho proto studovat izolovaně od tohoto prostředí.

\textbf{Ztělesnění} (\emph{embodiment}). Inteligentní chování není jen logika: je produktem
agenta, který \emph{má tělo} a~fyzicky interaguje se světem.

\textbf{Emergence}. Inteligentní chování \emph{vzniká} (emerguje) interakcí mnoha jednoduchých
chování a~není nikde explicitně naprogramováno.
\end{defbox}

Z~těchto tezí plynou i~charakteristiky reaktivních agentů. Jsou \emph{behaviorální}, protože
těžiště návrhu leží ve vývoji a~kombinaci jednotlivých chování, a~právě jejich repertoár
agenta definuje. Jsou \emph{situovaní}, tedy neoddělitelní od svého prostředí. A~jsou
\emph{reaktivní}: agent v~podstatě jen reaguje na prostředí a~nedeliberuje. Tomu odpovídá
i~reprezentace, se kterou pracuje, protože ta je \textbf{subsymbolická}. Bývají to
konekcionistické sítě, konečné automaty nebo jednoduchá pravidla typu IF--THEN.

\subsection{Symbolické reaktivní plánování: subsumpční architektura}
\label{ssec:subsumpce}

Reaktivní program dotáhl nejdál robotik, který chtěl ukázat, že se i~poměrně složité chování
obejde bez modelu světa. Konflikty mezi chováními u~něj rozhoduje pevně dané pořadí priorit.

\begin{defbox}[Subsumpční architektura (\emph{subsumption architecture})]
Autorem je Rodney Brooks, 1986--1991, a~jde o~nejúspěšnější reaktivní přístup vůbec. Vychází
ze tří Brooksových tezí:
\begin{enumerate}
\item Inteligentní chování lze generovat \emph{bez symbolické reprezentace}.
\item Inteligentní chování lze generovat \emph{bez explicitního abstraktního uvažování}.
\item Inteligence je \emph{emergentní vlastnost} určitých složitých systémů.
\end{enumerate}

Inteligence agenta je realizována množinou jednoduchých cílově orientovaných \textbf{chování}
(\emph{behaviours}), pro něž platí:
\begin{itemize}
\item Každé chování je samo o~sobě mechanismem výběru akce, totiž dvojicí
  \emph{situace $\to$ akce}.
\item Každé chování přijímá vjemy a~transformuje je na akce. Odpovídá vždy za jeden dílčí cíl
  a~má jednoduchou strukturu pravidla.
\item Chování běží \emph{paralelně} a~\emph{soupeří} spolu o~kontrolu nad agentem.
\item Chování jsou uspořádána do \textbf{subsumpční hierarchie}, která definuje jejich
  priority: vyšší vrstva nižší \uv{pohlcuje}, tedy potlačuje ji.
\end{itemize}
\end{defbox}

\noindent\textbf{Mechanismus výběru akce v~subsumpční architektuře:}
\begin{algorithmic}[1]
\Function{action}{vjem $p$}
  \State $\mathit{fired} \gets \{\, b \in \mathit{Behaviours} \mid p \text{ splňuje podmínku } b \,\}$
    \Comment{chování, která \uv{vystřelí}}
  \ForAll{$b \in \mathit{fired}$}
    \If{neexistuje $b' \in \mathit{fired}$ takové, že $b' \prec b$}
      \Comment{$b$ není nikým inhibováno}
      \State \Return akce příslušná chování $b$
    \EndIf
  \EndFor
  \State \Return \textbf{null} \Comment{nevystřelilo nic, agent nedělá nic}
\EndFunction
\end{algorithmic}

Formálně je chování dvojice $\langle c, a\rangle$, kde $c \subseteq \mathit{Per}$ je podmínka
a~$a \in A$ akce. Na množině chování $R$ je definována relace inhibice
$\prec\ \subseteq R \times R$, která tvoří striktní totální uspořádání. Pozor na to, kterým
směrem se relace čte, protože právě tady se dá snadno chybovat. U~Wooldridge i~na přednášce
znamená $b_1 \prec b_2$, že \textbf{$b_1$ inhibuje $b_2$}, takže $b_1$ leží \emph{níže}
v~subsumpční hierarchii a~má \emph{přednost}. Zápis $b_1 \prec b_2 \prec \dots \prec b_n$ proto
vyjmenovává chování od nejvyšší priority k~nejnižší.

Celý mechanismus má dvě přednosti. Je \emph{jednoduchý}, byť při desítkách chování se
nastavení priorit stává záludným, a~je \emph{velmi rychlý}: dá se implementovat hardwarově
a~má konstantní složitost.

\medskip
\noindent\textbf{Příklad: Steelsův průzkumník Marsu.}
Cílem je sbírat vzácné vzorky hornin na vzdálené planetě, a~to rojem robotických průzkumníků.
Základna k~tomu vysílá navigační signál, takže roboti nemusí umět nic víc než detekovat
\emph{gradient} tohoto signálu. Každý robot navíc veze radioaktivní drobky, kterými
\emph{nepřímo komunikuje} s~ostatními, tedy pomocí stigmergie. Přímá komunikace tak vůbec není
potřeba.

\emph{První iterace} zvládne jen náhodnou procházku a~návrat jednoho robota. Priorita pravidel
klesá v~tabulce shora dolů:
\begin{center}\small
\begin{tabular}{ll}
\toprule
R1 & \textbf{if} vidím překážku \textbf{then} změň směr \\
R2 & \textbf{if} nesu vzorek \textbf{and} jsem na základně \textbf{then} polož vzorek \\
R3 & \textbf{if} nesu vzorek \textbf{and} nejsem na základně \textbf{then} jdi po gradientu nahoru \\
R4 & \textbf{if} vidím vzorek \textbf{then} seber ho \\
R5 & \textbf{if} true \textbf{then} pohybuj se náhodně \\
\bottomrule
\end{tabular}
\end{center}

\emph{Druhá iterace} zlepší průzkum pomocí stigmergie. Vzorky se totiž často vyskytují ve
shlucích, a~tak robot, který jeden našel, o~tom ostatním \uv{řekne} tím, že cestou domů
rozhazuje drobky:
\begin{center}\small
\begin{tabular}{ll}
\toprule
R6 & \textbf{if} nesu vzorek \textbf{and} jsem na základně \textbf{then} polož vzorek \\
R7 & \textbf{if} nesu vzorek \textbf{and} nejsem na základně \textbf{then} polož 2 drobky
     a~jdi po gradientu nahoru \\
R8 & \textbf{if} cítím drobek \textbf{then} seber 1~drobek a~jdi po gradientu dolů \\
\bottomrule
\end{tabular}
\end{center}
Priority vypadají v~notaci \uv{levé inhibuje pravé}, tedy seřazeny od nejvyšší priority, takto:
$\mathrm{R1} \prec \mathrm{R6} \prec \mathrm{R7} \prec \mathrm{R4} \prec
\mathrm{R8} \prec \mathrm{R5}$. Robot, který sleduje stopu, z~ní bere vždy po jednom drobku,
takže se stopa postupně \emph{vypaří}, jakmile je shluk vytěžen. Roj tedy vykazuje kooperativní
chování, přestože žádný jednotlivý robot nemá model světa ani neposílá zprávy.

\subsection{Konekcionistické reaktivní plánování: agent network architecture}
\label{ssec:maes}

Subsumpční architektura řeší konflikty \emph{pevnou} prioritou, a~to je rigidní. Konekcionistický
přístup proto pevné priority nahradí \textbf{spojitou aktivací šířenou sítí}, takže výběr akce
vzejde z~dynamiky sítě, a~ne z~pořadí zapsaného předem. Výsledku se říká \uv{reaktivní
plánování}: agent se chová, jako by plánoval, ale žádný explicitní plán nikdy nevytvoří.

\begin{defbox}[Agent network architecture / \emph{spreading activation networks}]
Autorkou je Pattie Maes, 1989--1991. Agent je množina \textbf{kompetenčních modulů}
(\emph{competence modules}), které se podobají Brooksovým chováním. Každý modul $i$ má:
\begin{itemize}
\item \textbf{předpoklady} $c_i$ (\emph{preconditions}), tedy to, co musí platit, aby byl modul
  \emph{proveditelný} (\emph{executable});
\item \textbf{efekty}, kterými jsou \emph{add list} $a_i$ a~\emph{delete list} $d_i$, stejně
  jako je tomu u~STRIPS;
\item \textbf{úroveň aktivace} $\alpha_i$ a~globální \textbf{práh aktivace} $\theta$.
\end{itemize}
Moduly jsou propojeny v~orientovaném grafu podle svých podmínek: hrany tvoří shoda pre-
a~postpodmínek a~k~nim přibývají hrany reprezentující časovou následnost a~konflikty.
V~každém kroku se aktivace šíří a~\textbf{vybrán je proveditelný modul s~nejvyšší aktivací},
pokud jeho aktivace překročí práh $\theta$.
\end{defbox}

Aktivace přitéká do sítě ze tří zdrojů:
\begin{itemize}
\item \textbf{Ze stavu světa} (\emph{state}) ji dostávají moduly, jejichž předpoklady jsou
  momentálně splněny, tedy ty, o~nichž platí \uv{teď se dá udělat tohle}.
\item \textbf{Z~cílů} (\emph{goals}) ji dostávají moduly, jejichž add list obsahuje cíl, tedy
  ty, o~nichž platí \uv{tohle vede k~cíli}.
\item \textbf{Z~ochranných cílů} (\emph{protected goals}) naopak plyne \emph{inhibice}, a~to
  u~modulů, jejichž delete list obsahuje již splněný cíl. O~těch totiž platí \uv{tohle by
  zbořilo, co už mám}.
\end{itemize}
Mezi moduly navzájem se pak aktivace šíří třemi typy spojů:
\begin{itemize}
\item \textbf{Následnické spoje} (\emph{successor links}) vedou od modulu $x$ k~modulu $y$
  tehdy, když $y$ přidává něco, co $x$ potřebuje jako předpoklad. Jde o~\emph{dopředné}
  řetězení ve smyslu \uv{až budu hotov, půjde to dál}.
\item \textbf{Předchůdcovské spoje} (\emph{predecessor links}) používá modul $x$, který není
  proveditelný: posílá aktivaci modulům, jež by mu splnily chybějící předpoklad. Jde
  o~\emph{zpětné} řetězení ve smyslu \uv{pomoz mi, ať můžu běžet}.
\item \textbf{Konfliktní spoje} (\emph{conflictor links}) znamenají, že modul $x$
  \emph{inhibuje} modul $y$, který ruší jeho předpoklady.
\end{itemize}
Chování celé sítě řídí pět globálních parametrů. Práh aktivace $\theta$ je jediným prahem
v~systému, $\phi$ je aktivace vpravená do sítě stavem světa, $\gamma$ aktivace vpravená cíli,
$\delta$ aktivace odebraná chráněnými cíli a~$\pi$ \emph{průměrná} úroveň aktivace, na kterou
se síť po každém kroku normalizuje, aby se celková energie neměnila. Dva z~parametrů se dají
číst přímo jako návrhové kompromisy. Poměr $\gamma/\phi$ přesně vyjadřuje kompromis
\textbf{proaktivita vs.\ reaktivita}, zatímco $\theta$ určuje \textbf{uvážlivost vs.\ rychlost}:
vysoký práh $\to$ agent déle váhá, ale rozhoduje se lépe. Pokud v~daném kroku práh nepřekročí
žádný proveditelný modul, sníží se $\theta$ o~10~\% a~zkouší se znovu. Síť tak nemůže
\uv{zamrznout}.

\textbf{Jak se aktivace v~jednom kroku počítá.} Aktualizace probíhá ve čtyřech fázích.
Nejprve přijde na řadu (1)~\emph{vstup z~prostředí}, kdy každý pravdivý fakt rozdělí kvantum
$\phi$ rovným dílem mezi moduly, které jej mají ve svých předpokladech. Následuje
(2)~\emph{vstup od cílů}: každý cíl rozdělí $\gamma$ mezi moduly, které ho přidávají, a~každý
chráněný cíl naopak \emph{odebere} $\delta$ modulům, jež by ho rušily. Ve třetí fázi nastane
(3)~\emph{šíření mezi moduly}. Proveditelný modul posílá část své aktivace dopředu následníkům,
opět rozdělenou rovným dílem mezi příjemce, neproveditelný modul ji posílá zpět modulům, které
mu mohou splnit chybějící předpoklad, a~konfliktní spoje odpovídající díl aktivace odečítají.
Krok uzavírá (4)~\emph{normalizace}, při níž se aktivace všech modulů přeškálují tak, aby jejich
průměr byl $\pi$, a~celková energie sítě proto zůstává konstantní. Teprve poté se vybere
proveditelný modul s~aktivací nad prahem $\theta$; provedený modul svou aktivaci vynuluje.

\emph{Poznámka k~přednášce:} slidy uvádějí práh aktivace jako vlastnost jednotlivého modulu
a~mluví o~něm jako o~prioritě. V~Maesově původním modelu je ovšem práh $\theta$
\emph{globální} a~roli \uv{priority} hraje aktuální úroveň aktivace $\alpha_i$, která se
dynamicky mění.

\medskip\noindent
\textbf{Srovnání se subsumpcí:} pořadí chování zde není pevné, ale vzniká \emph{dynamicky}
z~aktuálního stavu světa a~z~cílů. Síť proto dokáže vygenerovat i~vícekrokové \uv{plánovité}
sekvence, protože se aktivace řetězí zpět od cíle přes předchůdcovské spoje, a~přitom nikdy
nevznikne explicitní datová struktura plánu. Odtud pochází název \emph{reaktivní plánování}.
Cenou za to je citlivost na nastavení parametrů a~nemožnost zaručit optimalitu či úplnost.

\subsection{Omezení reaktivních architektur}

Reaktivní architektury si za svou rychlost a~jednoduchost platí, a~to hned ve čtyřech ohledech:

\begin{itemize}
\item Reaktivní agenti \textbf{nemají model světa} a~všechno musí odvodit z~aktuálních vjemů.
  Někdy proto musí \emph{měnit prostředí}, aby si informaci \uv{zapamatovali}. Přesně to dělají
  radioaktivní drobky, které fungují jako externalizovaná paměť, tedy jako stigmergie.
\item Mají jen \textbf{krátkodobý pohled}, protože jednají podle aktuálního stavu a~podle
  lokální informace. Zohlednit globální podmínky a~dlouhodobé cíle je pak obtížné.
\item Emergence \textbf{není dobrý inženýrský postup}: chování je nutné vyladit experimentálně
  a~chybí k~tomu jak metodologie, tak možnost verifikace.
\item Návrh je zvládnutelný pro několik vrstev, ale \textbf{s~desítkami chování se stává
  neřešitelným}, protože počet interakcí mezi nimi roste kvadraticky.
\end{itemize}

\subsection{Hybridní architektury}
\label{ssec:hybridni}

Ani čistě reaktivní, ani čistě deliberativní architektura tedy není ideální. Nabízí se proto
obojí zkombinovat, a~právě to dělají hybridní architektury:
\begin{itemize}
\item \textbf{Deliberativní neboli plánovací komponenta} pracuje na symbolické úrovni a~vytváří
  reprezentace a~plány.
\item \textbf{Reaktivní komponenta} zajišťuje okamžité reakce bez složitých výpočtů.
\end{itemize}
Komponenty bývají uspořádány do \textbf{vrstev} (\emph{layered architectures}) a~reaktivní
vrstvy mají zpravidla \emph{přednost} před deliberativními, protože bezpečnost je důležitější
než optimalita. Uspořádat vrstvy vůči senzorům a~efektorům jde dvěma způsoby.

\begin{defbox}[Horizontální vs.\ vertikální vrstvení]
\textbf{Horizontální vrstvení:} všechny vrstvy jsou \emph{paralelně} připojeny k~senzorům
i~k~efektorům. Každá vrstva se chová jako samostatný agent a~navrhuje akci.
\begin{itemize}
\item[$+$] Koncepčně je to jednoduché, protože pro $n$ druhů chování stačí $n$ vrstev.
\item[$-$] Vrstvy se navzájem ovlivňují a~jejich návrhy si odporují, takže je nutná
  \textbf{mediační funkce} (\emph{mediator function}), která konflikt rozhodne.
  Právě ta je úzkým hrdlem: musí uvažovat $m^n$ kombinací, kde $n$ je počet vrstev a~$m$ počet
  možných návrhů, a~zároveň je kritickým bodem selhání.
\end{itemize}

\textbf{Vertikální vrstvení:} vrstvy jsou zapojeny \emph{sériově}.
\begin{itemize}
\item Ve variantě \emph{one-pass} vjem vstoupí do nejnižší vrstvy, prochází nahoru a~akce
  vystupuje z~nejvyšší. Uspořádání i~hierarchie chování jsou tak přirozené.
\item Ve variantě \emph{two-pass} putují vjemy \emph{zdola nahoru} a~řídicí povely, tedy akce,
  \emph{shora dolů}. Tok připomíná řízení ve skutečné firmě.
\item[$+$] Odpadá potřeba mediace a~složitost interakcí je lineární, protože rozhraní je
  jen $n-1$.
\item[$-$] Chybí robustnost: selhání kterékoli vrstvy shodí celého agenta.
\end{itemize}
\end{defbox}

\begin{defbox}[TouringMachines (Ferguson, 1992)]
\textbf{Horizontální} třívrstvá architektura pro simulovaného agenta v~dopravním prostředí.
\begin{itemize}
\item \textbf{Modelovací vrstva} (\emph{modelling layer}) drží symbolické modely agenta,
  prostředí a~ostatních agentů. Předvídá a~řeší konflikty, stanovuje cíle a~posílá je
  plánovací vrstvě.
\item \textbf{Plánovací vrstva} (\emph{planning layer}) zajišťuje proaktivní chování: vybírá
  z~předpřipravených plánů podobně jako PRS a~sestavuje z~nich plán cesty.
\item \textbf{Reaktivní vrstva} (\emph{reactive layer}) obsahuje klasická reaktivní pravidla
  \emph{situace $\to$ akce} a~stará se o~rychlou okamžitou reakci, třeba o~vyhýbání překážkám.
\end{itemize}
Vrstvy pracují paralelně a~jejich výstupy sjednocuje sada \textbf{řídicích pravidel}
(\emph{control rules}), což je konkrétní podoba mediační funkce. Tato pravidla mohou vstupy
a~výstupy vrstev potlačovat: slouží k~tomu pravidla \emph{censor} a~\emph{suppressor}.
\end{defbox}

\begin{defbox}[InteRRaP (Müller, 1994)]
\textbf{Vertikální dvouprůchodová} (\emph{two-pass}) třívrstvá architektura.
\begin{itemize}
\item \textbf{Vrstva kooperativního plánování} (\emph{cooperative planning layer}) obstarává
  sociální interakce a~společné plány s~ostatními agenty.
\item \textbf{Vrstva lokálního plánování} (\emph{local planning layer}) zajišťuje běžné
  individuální plánování.
\item \textbf{Behaviorální vrstva} (\emph{behaviour-based layer}) obsahuje reaktivní chování.
\end{itemize}
Každá vrstva má \emph{vlastní znalostní bázi} na příslušné úrovni abstrakce, tedy znalosti
o~světě, plány, případně společné plány a~modely ostatních agentů. Mezi sebou vrstvy komunikují
dvěma operacemi: \emph{aktivací zdola nahoru}, kdy nižší vrstva situaci nezvládne a~předá ji
výš, a~\emph{zásahem shora dolů}, kdy vyšší vrstva používá nižší k~vykonání svých rozhodnutí.
\end{defbox}

\medskip
\noindent\textbf{Reálný příklad: Stanley} (Volkswagen Touareg R5, Stanford). Tento autonomní
vůz vyhrál DARPA Grand Challenge 2005, tedy závod na 212~km Mohavskou pouští, a~je přímým
předchůdcem dnešních autonomních vozidel. Jeho architektura je vrstvená:
\begin{itemize}
\item vrstva senzorů $\to$ vrstva abstraktní percepce $\to$ vrstva plánování a~řízení,
  která sestaví plán trasy a~řídí vůz, $\to$ vrstva rozhraní vozidla;
\item k~tomu přistupuje vrstva uživatelského rozhraní (panel, start) a~vrstva globálních služeb
  (souborový systém, komunikace, hodiny).
\end{itemize}

\subsection{Komplexní architektury: IDA a~globální pracovní prostor}
\label{ssec:ida}

Poslední rodina architektur jde opačným směrem než reaktivní minimalismus. Nesnaží se agenta
osekat na nezbytné minimum, ale integrovat do něj \emph{všechno}: výběr akce, paměť, deliberaci
i~emoce.

\begin{defbox}[IDA --- \emph{Intelligent Distribution Agent} (S.~Franklin)]
Jedna z~nejsložitějších agentních architektur, jaké kdy byly postaveny. Vznikla pro reálnou
úlohu amerického námořnictva, totiž pro přidělování personálu na nová služební místa, odkud
pochází ono \uv{distribution}; distribuovaná je ostatně i~sama architektura.

\textbf{Východiskem je teorie globálního pracovního prostoru} (\emph{global workspace theory},
Baars, 1988--97), jedna z~teorií vědomí. Stojí na těchto předpokladech:
\begin{itemize}
\item Mysl \emph{je} multiagentní systém, protože se skládá z~mnoha jednoduchých, převážně
  nevědomých procesů, které provádějí specializované operace.
\item Procesy spolu komunikují \emph{zřídka} a~pouze přes sdílenou paměť typu \emph{blackboard},
  organizovanou jako asociativní pole (srov.\ kap.~\ref{sec:kooperace}).
\item Procesy dynamicky vytvářejí \textbf{koalice} vyšších řádů.
\item \textbf{Mechanismus výběru akce} spočívá v~tom, že koalice, která nejlépe odpovídá
  aktuálním vstupům, se dostane do \uv{vědomí}, tedy do globálního pracovního prostoru,
  a~je vykonána.
\item Existuje \emph{hierarchie kontextů} reprezentujících svět na různých úrovních: kontext
  cílů, vjemů, smyslů, konceptů a~kontext kulturní.
\end{itemize}
\textbf{Hodnocení:} architektura kombinuje mnoho přístupů z~MAS i~ze zbytku UI, tedy výběr
akce, paměť, deliberaci a~emoce, a~právě v~tom je zároveň její slabina. Řada rozhodnutí je
totiž \emph{ad hoc} a~sladit tak heterogenní kolekci modelů do optimálně spolupracujícího
celku je velmi obtížné.
\end{defbox}

Při srovnávání mechanismů výběru akce je IDA užitečný jako \emph{čtvrtý} typ arbitráže.
Subsumpce rozhoduje \textbf{pevnou prioritou}, síť Maes \textbf{úrovní aktivace}, vrstvené
architektury \textbf{mediační funkcí} a~IDA \textbf{soutěží dynamicky vznikajících koalic}
o~přístup do globálního pracovního prostoru.

\subsection{Srovnání architektur}

Následující tabulka staví vedle sebe všechny rodiny architektur probrané v~této kapitole
a~ukazuje, že jednotlivé vlastnosti se vzájemně vyvažují: co jedna rodina získá na rychlosti,
ztrácí na proaktivitě, a~naopak.

\begin{center}
\small
\begin{tabular}{@{}p{0.15\textwidth}
  >{\raggedright\arraybackslash}p{0.19\textwidth}
  >{\raggedright\arraybackslash}p{0.19\textwidth}
  >{\raggedright\arraybackslash}p{0.17\textwidth}
  >{\raggedright\arraybackslash}p{0.17\textwidth}@{}}
\toprule
& \textbf{Deduktivní} & \textbf{BDI / PRS} & \textbf{Reaktivní} & \textbf{Hybridní} \\
\midrule
Reprezentace & symbolická (FOL) & symbolická (beliefs, plány) & subsymbolická / pravidla
  & smíšená po vrstvách \\
Výběr akce & dokazování $\mathit{Do}(a)$ & deliberace + plán z~knihovny
  & priorita / aktivace & vrstvy + arbitráž \\
Model světa & úplný & částečný, revidovatelný & žádný & po vrstvách \\
Rychlost & velmi nízká & střední & vysoká (konst.) & vysoká v~reflexu \\
Reaktivita & špatná & střední (laditelná) & výborná & dobrá \\
Proaktivita & dobrá & výborná & žádná & dobrá \\
Návrh & obtížný, ale formální & systematický & experimentální & složitý \\
\bottomrule
\end{tabular}
\end{center}


\subsection{Shrnutí ke~zkoušce}

\begin{itemize}
\item \textbf{Deduktivní agent} pracuje s~databází $\mathit{DB} \in 2^L$ a~akci vybírá
  dokazováním $\mathit{DB} \vdash_\rho \mathit{Do}(a)$. Výběr akce je dvoufázový: hledá se
  nejprve dokazatelná akce, pak akce konzistentní, a~když neuspěje ani to, agent neudělá nic.
  Přístup je elegantní, ale nepraktický, protože naráží na problém transdukce a~na
  vypočitatelnou racionalitu.
\item \textbf{BDI} má dvě fáze. První je deliberace, kterou tvoří $\mathit{brf}$,
  $\mathit{options}$ a~$\mathit{filter}$, druhou means-ends reasoning, tedy plánování. Beliefs
  jsou subjektivní a~omylné, desires mohou být nekonzistentní a~intentions jsou závazky, které
  přetrvávají a~omezují další deliberaci.
\item \textbf{STRIPS} popisuje akci deskriptorem $\langle P_a,D_a,A_a\rangle$, plánovací problém
  trojicí $\langle B_0,O,G\rangle$ a~stav aktualizuje předpisem
  $B_i = (B_{i-1}\setminus D_{a_i})\cup A_{a_i}$. Plán je \emph{přípustný}, pokud platí
  $B_{i-1}\models P_{a_i}$, a~\emph{korektní}, pokud navíc $B_n \models G$. U~zkoušky je třeba
  umět projít svět kostek.
\item Zvlášť dobře je potřeba znát \textbf{BDI smyčku} a~roli funkcí $\mathit{reconsider}$
  a~$\mathit{sound}$.
\item \textbf{Commitment} má tři podoby: blind, single-minded a~open-minded. Podle ochoty
  přehodnocovat se agenti dělí na bold a~cautious, efektivita se počítá jako podíl dosažených
  a~všech záměrů a~platí, že v~pomalém prostředí vyhrává bold agent, kdežto v~rychlém cautious.
\item \textbf{PRS} je první implementací BDI. Plán se skládá z~částí Goal, Context a~Body, místo
  plánování se používá knihovna plánů a~záměry drží zásobník. Deliberace se řeší metaplány nebo
  utilitou a~potomky systému jsou AgentSpeak/Jason, JAM, JACK a~Jadex.
\item Reaktivní přístup stojí na trojici \emph{situatedness --- embodiment --- emergence}
  a~odmítá symbolickou reprezentaci.
\item \textbf{Subsumpce} (Brooks) představuje \emph{symbolické} reaktivní plánování. Množina
  jednoduchých chování \emph{situace~$\to$~akce} běží paralelně a~je uspořádána do pevné
  priority. U~zkoušky je potřeba umět Steelsova průzkumníka Marsu a~roli stigmergie, tedy
  drobků.
\item \textbf{Agent network architecture} (Maes) je naopak \emph{konekcionistické} reaktivní
  plánování. Moduly mají předpoklady a~efekty (add, delete), aktivace se šíří následnickými,
  předchůdcovskými a~konfliktními spoji a~jejími zdroji jsou stav světa, cíle a~inhibice
  od chráněných cílů. Vybere se proveditelný modul s~nejvyšší aktivací nad prahem a~poměr
  parametrů $\gamma/\phi$ ladí proaktivitu vůči reaktivitě.
\item Omezení reaktivních agentů jsou čtyři: nemají model světa, mají jen krátkodobý pohled,
  jejich návrh není inženýrský a~množství vrstev se neškáluje.
\item \textbf{Hybridní} architektury se dělí na horizontální a~vertikální. U~horizontálních běží
  vrstvy paralelně a~rozhoduje mediátor, který musí uvažovat $m^n$ kombinací, a~je proto úzkým
  hrdlem. Vertikální jsou sériové ve variantě one-pass nebo two-pass, vystačí si s~$n-1$
  rozhraními, ale jsou křehké.
\item \textbf{TouringMachines} jsou horizontální 3~vrstvy (modelling, planning, reactive)
  doplněné řídicími pravidly, kdežto \textbf{InteRRaP} je vertikální two-pass se 3~vrstvami
  (kooperativní plánování, lokální plánování, behaviorální) a~s~vlastní KB na každé vrstvě.
  \textbf{Stanley} slouží jako příklad reálného nasazení.
\item \textbf{IDA} (Franklin) staví na \emph{global workspace theory} (Baars) a~chápe mysl jako
  MAS jednoduchých nevědomých procesů, které komunikují přes blackboard. Koalice procesů soutěží
  o~vstup do \uv{vědomí} a~vítězná se vykoná, přičemž existuje hierarchie kontextů. Jde
  o~nejsložitější architekturu a~je v~ní mnoho ad hoc rozhodnutí.
\item Arbitráž mezi chováními má čtyři typy: \textbf{pevnou prioritu} u~subsumpce,
  \textbf{úroveň aktivace} u~Maes, \textbf{mediační funkci} u~vrstvených architektur
  a~\textbf{soutěž koalic} u~IDA.
\end{itemize}

%=====================================================================

\section{Problém hledání cesty}
\label{sec:cesta}
%=====================================================================

\noindent\mimoq{celá kapitola} Přednáška NAIL106 řadí plánování i~robotiku mezi témata, která
\emph{nepokrývá}, a~ve slidech ke~hledání cesty proto nic není. Okruh ho ale výslovně jmenuje,
takže se mu vyhnout nejde. Kapitola je držena na minimu, které u~zkoušky obstojí:
algoritmus A* i~s~vlastnostmi heuristiky, a~pak už jen to, co je na hledání cesty vysloveně
\emph{multiagentního}, tedy dynamické prostředí a~koordinace více agentů.

\subsection{Formulace a~zařazení}

Hledání cesty (\emph{pathfinding}) je nejčastější konkrétní podúloha, kterou musí situovaný
agent řešit. Objevuje se v~plánovací vrstvě hybridní architektury, vystupuje jako akce
\texttt{move\_towards(P)} v~BDI plánu a~agent ji také může nabízet ostatním jako službu.
V~multiagentním kontextu k~běžnému zadání přibývají dvě specifika. Prostředí je
\textbf{dynamické}, tedy mapa se za běhu mění, a~ostatní agenti jsou \textbf{pohyblivé
překážky}, se kterými se agent musí koordinovat.

\begin{defbox}[Problém hledání cesty]
Je dán ohodnocený graf $G = (V,E)$ s~nezáporným ohodnocením hran $w : E \to \mathbb{R}_0^+$,
počáteční vrchol $s \in V$ a~cílový vrchol $g \in V$. Hledáme cestu $s = v_0, v_1, \dots, v_k = g$
minimalizující $\sum_{i} w(v_{i-1},v_i)$.

Graf přitom obvykle není vypsaný, ale zadaný \emph{implicitně} nástupnickou funkcí; typicky jde
o~4- nebo 8-souvislou mřížku, o~navigační síť nebo o~graf viditelnosti.
\end{defbox}

\paragraph{Odkud se vezmou vrcholy.} Prostředí je spojité, kdežto prohledávání potřebuje graf.
Musí tedy proběhnout diskretizace a~právě její volba rozhoduje o~kvalitě výsledné cesty víc než
volba samotného algoritmu. Nabízí se několik možností:
\begin{itemize}
\item \textbf{Mřížka} (4- nebo 8-souvislá) se implementuje triviálně, ale cesta po ní vychází
  \uv{schodovitě} a~vrcholů je obrovské množství.
\item \textbf{Graf viditelnosti} má za vrcholy rohy mnohoúhelníkových překážek a~jeho hrany
  spojují dvojice vzájemně viditelných rohů. V~rovině s~mnohoúhelníkovými překážkami dá
  \emph{skutečně optimální} cestu, jenže jeho velikost roste kvadraticky s~počtem rohů.
\item \textbf{Navigační síť} (\emph{navmesh}) pokryje prostor konvexními plochami, uvnitř nichž
  je pohyb volný. Je to standard v~herních enginech a~v~robotice.
\item \textbf{Stavová mřížka} bere jako vrchol nejen polohu, ale i~orientaci a~případně
  rychlost. Bez toho se neobejdou neholonomní vozidla, která se nemohou otáčet na místě.
\item \textbf{Vzorkovací metody} graf nestaví celý předem.
  \emph{PRM} (\emph{probabilistic roadmap}) si předpočítá graf náhodných konfigurací,
  zatímco \emph{RRT} a~\emph{RRT*} nechávají strom inkrementálně růst. Pro vysokodimenzionální
  konfigurační prostory, jako je konfigurační prostor robotického ramene, jsou jedinou
  praktickou volbou.
\item \textbf{Potenciálová pole} nechají cíl přitahovat a~překážky odpuzovat, agent pak jde
  po gradientu. Jsou velmi levná a~reaktivní, takže patří spíš do reaktivní vrstvy
  (\S\ref{ssec:subsumpce}), zato \emph{uvíznou v~lokálním minimu} a~nedávají žádnou záruku
  úplnosti.
\end{itemize}

\subsection{A*}

Prohledávání typu \emph{best-first} si ke~každému vrcholu $n$ vede tři hodnoty:
\begin{align*}
g(n) &= \text{cena nejlepší dosud nalezené cesty } s \to n, \\
h(n) &= \text{heuristický odhad ceny } n \to g_{\mathrm{cil}}, \\
f(n) &= g(n) + h(n) \quad \text{--- odhad ceny nejlepší cesty } s \to g_{\mathrm{cil}}
  \text{ vedoucí přes } n .
\end{align*}

\noindent Algoritmus vždy expanduje ten z~otevřených vrcholů, který má nejmenší $f$; narazí-li
přitom na levnější cestu k~vrcholu, který už zná, jeho hodnoty opraví.

\noindent\textbf{Algoritmus A*} (Hart, Nilsson, Raphael, 1968):
\begin{algorithmic}[1]
\State $\mathit{OPEN} \gets \{s\}$ (prioritní fronta podle $f$), $\mathit{CLOSED} \gets \emptyset$
\State $g(s) \gets 0$, $f(s) \gets h(s)$, ostatní $g \gets \infty$
\While{$\mathit{OPEN} \ne \emptyset$}
  \State $n \gets \arg\min_{x \in \mathit{OPEN}} f(x)$; odeber $n$ z~$\mathit{OPEN}$
  \If{$n$ je cíl} \Return rekonstruovaná cesta \EndIf
  \State přidej $n$ do $\mathit{CLOSED}$
  \ForAll{$m \in \mathit{succ}(n)$}
    \State $g_{\mathrm{new}} \gets g(n) + w(n,m)$
    \If{$g_{\mathrm{new}} < g(m)$}
      \State $g(m) \gets g_{\mathrm{new}}$; $\mathit{parent}(m) \gets n$;
             $f(m) \gets g(m) + h(m)$
      \State vlož/aktualizuj $m$ v~$\mathit{OPEN}$ (a~odeber z~$\mathit{CLOSED}$, byl-li tam)
    \EndIf
  \EndFor
\EndWhile
\State \Return \uv{cesta neexistuje}
\end{algorithmic}

\medskip
\noindent Jaké záruky A* dává, závisí na vlastnostech heuristiky $h$. Rozlišují se u~ní dvě,
slabší a~silnější.

\begin{defbox}[Přípustnost a~konzistence heuristiky]
Heuristika $h$ je \textbf{přípustná} (\emph{admissible}), pokud nikdy nenadhodnocuje:
$h(n) \le h^*(n)$ pro všechna $n$, kde $h^*(n)$ je skutečná cena nejlevnější cesty z~$n$ do cíle.

Heuristika je \textbf{konzistentní} (\emph{consistent}, monotónní), pokud pro každou hranu
$(n,m)$ platí trojúhelníková nerovnost
\[
h(n) \le w(n,m) + h(m), \qquad h(g_{\mathrm{cil}}) = 0 .
\]
\textbf{Vlastnosti.} Konzistence je silnější požadavek než přípustnost, tedy konzistence
$\Rightarrow$ přípustnost, opačně to neplatí. Je-li heuristika přípustná, je A*
\emph{optimální}, tedy najde nejlevnější cestu. Konzistentní heuristika navíc zaručí, že $f$ je
podél každé cesty neklesající a~že se každý vrchol expanduje \emph{nejvýše jednou}.
\end{defbox}

\textbf{Speciální případy a~varianty.} Položíme-li $h \equiv 0$, degeneruje A* na
\textbf{Dijkstrův algoritmus}. Na mřížce se typicky používá heuristika manhattanská
$|\Delta x| + |\Delta y|$ pro 4~směry, oktilová pro 8~směrů a~eukleidovská. \textbf{Weighted A*}
počítá $f = g + \varepsilon h$ s~$\varepsilon > 1$: doběhne rychleji, ale nalezená cesta může být
až $\varepsilon\times$ dražší než optimum. \textbf{IDA*} zase šetří paměť, protože si vystačí
s~$O(d)$ místo exponenciálního nároku.

\paragraph{Zařazení mezi prohledávací algoritmy.} A* je jeden bod na škále, na jejíchž koncích
stojí dvě krajnosti. \emph{BFS} a~\emph{Dijkstra/uniform-cost} berou v~úvahu jen dosavadní cenu
$g$, takže jsou sice úplné a~optimální, ale prohledávají všemi směry stejně. Naopak
\emph{greedy best-first} bere jen odhad $h$, což je rychlé, zato \emph{neoptimální} a~snadno se
dá svést. A* obojí kombinuje a~při přípustné $h$ zůstává optimální. Prohledaný objem lze ještě
řádově zmenšit obousměrným prohledáváním, které postupuje současně od startu a~od cíle.

\paragraph{Zrychlení bez ztráty optimality.} Na mřížkách se používá \textbf{JPS}
(\emph{jump point search}). Mřížka obsahuje obrovské množství stejně dlouhých cest, které se
liší jen pořadím kroků, a~JPS tuto symetrii odstraní: \uv{nezajímavé} vrcholy přeskakuje
a~expanduje jen \emph{skokové body}. Jinou cestou jdou \textbf{Theta*} a~další \emph{any-angle}
algoritmy, které uvolní požadavek, aby cesta vedla po hranách mřížky. Při expanzi zkusí napojit
vrchol rovnou na \emph{dědečka}, je-li mezi nimi přímá viditelnost, a~dají tak podstatně kratší
a~přirozenější cesty. \textbf{HPA*} (\emph{hierarchical}) rozdělí mapu na clustery a~předpočítá
cesty mezi jejich vstupy; hledá pak nejdřív hrubě mezi clustery a~teprve potom uvnitř nich.

\medskip
\noindent\textbf{Malý příklad.} Mějme mřížku $3\times3$ s~pohybem do 4~směrů a~cenou hrany $1$,
start v~$(0,0)$, cíl v~$(2,2)$ a~překážku na $(1,1)$. Manhattanská heuristika
$h(x,y) = |2-x| + |2-y|$ je zde přípustná i~konzistentní. Start má $f=0+4=4$, jeho nástupci
$(1,0)$ a~$(0,1)$ mají $g=1$, $h=3$, $f=4$. Dále přijdou na řadu $(2,0)$ a~$(0,2)$
s~$g=2,h=2,f=4$, pak $(2,1)$ a~$(1,2)$ s~$g=3,h=1,f=4$ a~konečně cíl $(2,2)$ s~$g=4,h=0,f=4$.
Všechny vrcholy na optimální cestě tedy mají $f=4$, což je typický projev konzistentní
heuristiky, která zbytek cesty odhaduje přesně.

\subsection{Dynamické prostředí: D*~Lite}

Dosud jsme předpokládali, že mapa se během hledání nemění. Skutečný agent se ale pohybuje
a~průběžně zjišťuje, že skutečnost mapě neodpovídá: objeví se nová překážka, změní se cena
průchodu. Plánovat A* od nuly po každé takové zprávě je drahé. Řešením jsou \textbf{inkrementální}
algoritmy, které při lokální změně přepočítají jen tu část výpočtu, které se změna opravdu
dotkla.

\begin{defbox}[LPA* a~D*~Lite]
\textbf{LPA*} (\emph{Lifelong Planning A*}, Koenig, Likhachev, 2001) je inkrementální A* pro
pevný start a~cíl. Ke~každému vrcholu si kromě $g(n)$ udržuje ještě \emph{lookahead} hodnotu
\[
\mathit{rhs}(n) = \begin{cases}
0 & n = s,\\
\min_{p \in \mathit{pred}(n)} \big(g(p) + w(p,n)\big) & \text{jinak.}
\end{cases}
\]
Vrchol je \textbf{lokálně konzistentní}, pokud $g(n) = \mathit{rhs}(n)$. Změní-li se cena hrany,
stanou se některé vrcholy nekonzistentními; ty se zařadí do prioritní fronty a~změna se propaguje
jen tam, kde je to potřeba.

\textbf{D*~Lite} (Koenig, Likhachev, 2002) je LPA* prohledávající \emph{od cíle ke startu}, aby
spočtené hodnoty zůstaly použitelné i~poté, co se robot posune; pohyb se dorovnává korekcí
prioritní fronty. Funkčně je ekvivalentní staršímu algoritmu \textbf{D*} (Stentz, 1994), je však
výrazně jednodušší a~stal se \emph{de facto} standardem v~mobilní robotice.
\end{defbox}

V~neznámém prostředí se k~tomu přidává \emph{optimistický předpoklad}
(\emph{freespace assumption}): neznámé buňky se považují za průchodné. Agent tedy vyrazí po
naplánované cestě, a~jakmile narazí, přeplánuje. To je přesně scénář, pro který byl D* navržen.

\subsection{Hledání cest pro více agentů (MAPF)}
\label{ssec:mapf}

\noindent\mimoq{není ve slidech, ale je to jediná vysloveně multiagentní část okruhu}

Dosud cestu plánoval jediný agent a~ostatní pro něj byli nanejvýš pohyblivé překážky. Jakmile se
ale plánuje pro celou skupinu naráz, mění se zadání i~jeho obtížnost.

\begin{defbox}[Multi-Agent Path Finding, MAPF]
Vstupem je graf $G=(V,E)$ a~$k$ agentů, kde agent $i$ má start $s_i$ a~cíl $g_i$. Čas je
diskrétní a~v~každém kroku každý agent buď přejde po hraně, nebo počká na místě.

\textbf{Konflikty.} Řešení musí být \emph{bezkonfliktní}, a~to ve dvou ohledech.
\emph{Vrcholový konflikt} nastane, když se dva agenti ocitnou ve stejném čase ve stejném vrcholu,
\emph{hranový, takzvaný \uv{swap} konflikt} tehdy, když si dva agenti v~jednom kroku vymění
pozice.

\textbf{Účelová funkce} se volí buď jako \emph{sum-of-costs} $\sum_i \mathit{cost}_i$, nebo jako
\emph{makespan} $\max_i \mathit{cost}_i$. Nalezení \emph{optimálního} řešení je NP-těžké.
\end{defbox}

\textbf{Prioritizované plánování} je oddělený přístup k~úloze. Agenti se uspořádají podle
priority a~pro $i = 1,\dots,k$ se cesta agenta $i$ hledá algoritmem A* v~\emph{prostoru--čase},
kde stavem je dvojice $\langle$vrchol, čas$\rangle$, tak, aby se agent vyhnul už naplánovaným
trajektoriím uloženým v~\textbf{rezervační tabulce}. Tato varianta se nazývá \emph{Cooperative A*}
(Silver, 2005). Je rychlá a~dobře škáluje, zaplatí se to však \textbf{neúplností}: existují
řešitelné instance, které nevyřeší žádné pořadí priorit. Typickým příkladem jsou dva agenti
v~úzkém koridoru, kde jeden musí uhnout do výklenku.

Úplnost i~optimalitu vrací zpět jiný přístup. Nesnaží se konfliktům předejít pevným pořadím
agentů, ale řeší je až ve chvíli, kdy se opravdu objeví.

\begin{defbox}[Conflict-Based Search (CBS)]
Sharon a~kol., 2012--2015. Dvouúrovňový \emph{optimální} a~\emph{úplný} algoritmus, který se
vyhýbá exponenciálnímu sdruženému stavovému prostoru $V^k$.

\textbf{Nízká úroveň} hledá pomocí A* optimální cestu jednoho agenta v~prostoru--čase, a~to při
respektování množiny \emph{omezení} tvaru $\langle a_i, v, t\rangle$
(\uv{agent $i$ nesmí být ve vrcholu $v$ v~čase $t$}).

\textbf{Vysoká úroveň} prohledává \emph{strom omezení} podle ceny. Kořen má prázdnou množinu
omezení a~individuálně optimální cesty. Vezme se uzel s~nejnižší cenou; jsou-li jeho cesty
bezkonfliktní, je řešení optimální. Jinak se vybere první konflikt a~uzel se \emph{rozvětví}
na dva potomky, přičemž do každého z~nich se přidá omezení pro jednoho z~kolidujících agentů
a~přeplánuje se \emph{jen} on.
\end{defbox}

\textbf{Škálovatelné varianty.} Optimální CBS zvládne desítky agentů. Jde-li o~stovky až tisíce,
jak je tomu ve skladech Amazon/Kiva, obětuje se optimalita kontrolovaným způsobem. \emph{ECBS} je
\emph{bounded suboptimal}: na obou úrovních používá \emph{focal search} a~jeho výsledek je
zaručeně nejvýše $w\times$ horší než optimum. \emph{SIPP} (\emph{safe interval path planning})
nepracuje s~jednotlivými časovými kroky, ale se \emph{souvislými intervaly}, v~nichž je vrchol
volný, takže velikost stavového prostoru nezávisí na délce horizontu. \emph{MAPF-LNS}
(\emph{large neighborhood search}) opakovaně zahodí a~znovu naplánuje cesty malé podmnožiny
agentů.

\subsection{Shrnutí ke~zkoušce}

\begin{itemize}
\item Algoritmus A* pracuje s~$f = g+h$ a~jeho záruky plynou z~vlastností heuristiky:
  \emph{přípustnost} $h \le h^*$ $\Rightarrow$ optimalita, \emph{konzistence}
  $h(n) \le w(n,m)+h(m)$ $\Rightarrow$ každý vrchol se expanduje jen jednou. Pro $h\equiv 0$
  dostaneme Dijkstrův algoritmus, weighted A* je $\varepsilon$-suboptimální.
\item V~dynamickém prostředí se plánuje inkrementálně: LPA* si vede hodnoty $g$ a~\emph{rhs}
  a~sleduje lokální konzistenci $g=\mathit{rhs}$, D*~Lite prohledává od cíle a~při změně opravuje
  jen zasaženou část. V~neznámém prostředí se k~tomu přidává freespace assumption.
\item V~úloze MAPF se rozlišují konflikty vrcholové a~hranové (swap) a~účelová funkce se volí
  jako sum-of-costs, nebo jako makespan. Optimální MAPF je NP-těžký.
\item Prioritizované plánování, tedy Cooperative A* s~rezervační tabulkou, je rychlé, ale
  neúplné. \textbf{CBS} má nízkou úroveň v~podobě A* s~omezeními a~vysokou úroveň, která větví
  strom omezení podle nalezeného konfliktu; je optimální a~úplný. Škálovat lze pomocí ECBS,
  SIPP a~LNS.
\item O~kvalitě cesty rozhoduje diskretizace prostoru víc než algoritmus. Na výběr je mřížka,
  graf viditelnosti (optimální v~rovině s~mnohoúhelníky), navmesh, stavová mřížka, PRM a~RRT pro
  vysoké dimenze a~potenciálová pole, která jsou levná, ale uvíznou v~lokálním minimu.
\item Co se zařazení týče, BFS a~Dijkstra berou jen $g$, greedy best-first jen $h$ a~A* obojí.
  Bez ztráty optimality zrychlují JPS (odstraní symetrii v~mřížce), Theta* (any-angle) a~HPA*
  (hierarchie).
\end{itemize}

%=====================================================================

\section{Komunikace a~znalosti v~MAS: ontologie, řečové akty, FIPA-ACL, protokoly}
\label{sec:komunikace}
%=====================================================================
\subsection{Od agenta k~multiagentnímu systému}

Navrhnout jednoho agenta už umíme. Aby ale z~několika agentů vznikl systém, musí spolu
komunikovat, a~to je víc než jen protáhnout mezi nimi drát. Vyřešit se musí pět věcí najednou:

\begin{itemize}
\item \textbf{Reprezentace znalostí}, tedy v~čem vůbec jsou znalosti vyjádřeny.
\item \textbf{Společný slovník}, aby stejné slovo znamenalo u~obou stran totéž. Přesně to je
  úloha \emph{ontologií}.
\item \textbf{Standardní zprávy}, které komunikaci dávají pevnou syntaxi a~sémantiku
  (\emph{ACL}).
\item \textbf{Předvídatelné chování při komunikaci} --- konverzace se vede podle
  \emph{interakčních protokolů}, takže je dopředu zřejmé, co může přijít dál.
\item \textbf{Technické problémy}, které zůstanou i~poté: systém je distribuovaný, agenti se
  mohou pohybovat, běží asynchronně a~komunikační kanály nejsou spolehlivé.
\end{itemize}

První tři body jsou obsahem této kapitoly a~začínají u~slovníku, protože bez sdílených pojmů
nemá smysl řešit ani formát zpráv, ani jejich pořadí.

\subsection{Co je ontologie}

Slovo ontologie má dva různě staré významy a~oba se hodí znát, protože informatický pojem si
z~toho filozofického půjčuje ambici popsat, co ve světě existuje.

\begin{defbox}[Ontologie]
\emph{Filozoficky} je ontologie nauka o~bytí. Jméno vzniklo z~řeckého \emph{to on} (jsoucí)
a~\emph{logos} (slovo); Aristotelés ji nazval \uv{první filozofií} a~řadil ji do metafyziky.

\emph{V~informatice} znamená ontologie \textbf{explicitní a~formalizovaný popis části reality}.
Obvykle má podobu formálního deklarativního popisu, který obsahuje \emph{glosář} (definice
konceptů) a~\emph{tezaurus} (definice vztahů mezi koncepty). Ontologie je jakýsi slovník pro
ukládání a~výměnu znalostí o~doméně standardním způsobem.

\textbf{Gruber (1995):} \uv{Ontologie je popis (podobný formální specifikaci programu) konceptů
a~vztahů, které mohou pro agenta nebo komunitu agentů formálně existovat.}
Často citovaná zkrácená verze: \emph{ontologie je explicitní specifikace konceptualizace}.
\end{defbox}

\textbf{Motivační příklad} (dialog dvou lidí, nebo stejně tak dvou agentů). \uv{Slyšel jsi 7777?
--- Ne, co to je? --- To je nové CD kapely SRPR, takový alt.rock s~elektronikou a~skvělými
texty. --- Aha.} Aby si oba rozuměli, musí sdílet celou malou teorii hudebního světa.
7777 je album a~spadá do žánrů alt.rock a~electronica. Album má skladby, interpreta (zde SRPR)
a~autory hudby a~textů. SRPR je kapela, kapela má členy a~ti hrají na nástroje\dots

Aby se takový popis dal zapsat strojově, potřebuje pevnou kostru několika stavebních prvků.

\begin{defbox}[Formalismus ontologie]
\begin{itemize}
\item \textbf{Třídy} (\emph{classes}, koncepty) sdružují věci, které mají něco společného.
\item \textbf{Instance} (\emph{individuals}, objekty) jsou konkrétní jedinci patřící do tříd.
\item \textbf{Vlastnosti} (\emph{properties}) popisují atributy tříd a~instancí.
\item \textbf{Relace mezi třídami} spojují třídy navzájem; nejdůležitější z~nich je
  \textbf{podtřída} (\emph{is-a}), která je \emph{tranzitivní}.
\item Zbytek tvoří další relace a~vlastnosti, tedy základní znalosti o~doméně.
\end{itemize}
\textbf{Strukturálním částem} (třídám, relacím a~axiomům) se obvykle říká vlastní
\emph{ontologie}; teprve spolu s~\textbf{fakty o~konkrétních jedincích} tvoří \textbf{znalostní bázi}
(\emph{knowledge base}).
\end{defbox}

\textbf{Ontologie ontologií podle síly formalizace.} Ontologie se navzájem liší tím, jak přísně
jsou formalizované, a~mezi nejslabší a~nejsilnější podobou vede plynulá škála
(od slabých k~silným):
slovník (řízený slovník vybraných termínů) $\to$ glosář (definice významů v~přirozeném jazyce)
$\to$ tezaurus (synonyma) $\to$ neformální hierarchie (Amazon, wiki) $\to$
formální \emph{is-a} hierarchie (subsumpce tříd) $\to$ třídy s~vlastnostmi $\to$
omezení hodnot (\uv{každý člověk má právě jednu matku}) $\to$ libovolná logická omezení.
Čím dál po škále postoupíme, tím víc se dá vyjádřit a~odvodit, ale tím dráž: poslední stupeň
už vede na složité odvozovací algoritmy.

\textbf{Ontologie ontologií podle rozsahu.} Druhé dělení se řídí tím, jak široký výsek světa
ontologie popisuje:
\begin{itemize}
\item \textbf{Aplikační} ontologie jsou v~praxi nejběžnější, zato se těžko používají znovu jinde.
\item \textbf{Doménové} ontologie popisují celý obor a~jsou oblíbeným předmětem výzkumu
  sémantického webu.
\item \textbf{Horní} ontologie (\emph{upper ontologies}) by měly popisovat \uv{všechno}
  (Thing, Living thing, \dots) a~sloužit jako základ doménových ontologií. Patří sem BFO, GFO
  a~UFO; WordNet a~IDEAS naopak nejsou pro formální odvozování vhodné. Proti samotné snaze
  vybudovat univerzální ontologii ostatně stojí metodologické námitky
  (Wittgenstein, \emph{Tractatus logico-philosophicus}).
\end{itemize}

\subsection{Ontologické jazyky}

Ontologii je potřeba v~něčem zapsat. K~tomu slouží ontologické jazyky: formální deklarativní
jazyky pro reprezentaci znalostí, které obsahují fakta i~odvozovací pravidla. Nejčastěji se
opírají o~logiku prvního řádu nebo o~\emph{deskripční logiku}.

Historickým předchůdcem těchto jazyků jsou \textbf{rámce} (\emph{frames}) od Minského, oblíbené
v~klasické UI a~v~expertních systémech. Jejich stavba téměř přesně odpovídá objektovému
programování:

\begin{center}\small
\begin{tabular}{ll}
\toprule
\textbf{Terminologie rámců} & \textbf{Objektová terminologie} \\
\midrule
Frame (rámec) & třída objektů \\
Slot & vlastnost / atribut objektu \\
Trigger, accessor, mutator & metoda \\
\bottomrule
\end{tabular}
\end{center}

\textbf{XML} pro ontologie vytvořen nebyl, přesto se k~nim používá. Jeho hlavní výhodou je
možnost definovat vlastní značky, které samy o~sobě přirozeně reprezentují slovník domény:
\begin{quote}\small\ttfamily
<product typ=\textquotedbl CD\textquotedbl>\\
\hspace*{1em}<title>7777</title>\\
\hspace*{1em}<artist>SRPR</artist>\\
</product>
\end{quote}
Chybí mu ovšem definovaná sémantika: popisuje strukturu dokumentu, nikoli význam toho, co
v~dokumentu stojí.

\begin{defbox}[RDF --- \emph{Resource Description Framework}]
Standardní nástroj pro reprezentaci znalostí (nejen) pro web. Je záměrně \textbf{jednoduchý},
tedy málo expresivní, zato s~rychlými odvozovacími algoritmy.

Vše se vyjadřuje \textbf{trojicemi} \emph{subjekt -- predikát -- objekt}, tedy
$\mathit{Predikát}(\mathit{Subjekt},\mathit{Objekt})$. Například
$\mathit{MarriedTo}(\mathit{Karel},\mathit{Jája})$,
$\mathit{FatherOf}(\mathit{Karel},\mathit{Péťa})$.
Množina trojic tvoří orientovaný ohodnocený graf (\emph{RDF graph}) a~zdroje se v~něm
identifikují pomocí IRI. Serializovat se dá do RDF/XML, Turtle, N-Triples nebo JSON-LD,
dotazuje se jazykem SPARQL. Nadstavba \textbf{RDFS} přidává třídy a~konstrukce
\texttt{subClassOf}, \texttt{domain} a~\texttt{range}.
\end{defbox}

\begin{defbox}[OWL --- \emph{Web Ontology Language}]
Rovněž pochází ze (sémantického) webu a~existuje ve verzích 1 a~2. Není to jeden jazyk, ale
\emph{kolekce několika formalismů} pro popis ontologií; zapsat se dají v~XML, v~RDF, funkcionální
notací, Manchester syntaxí nebo v~Turtle. Za vším stojí snaha o~jazyk, který je formální,
a~přesto prakticky použitelný.

Základem je \textbf{deskripční logika}, tedy \emph{rozhodnutelný fragment} logiky prvního řádu.

Platí \textbf{předpoklad otevřeného světa} (\emph{open world assumption}, OWA): co nevíme, je
\emph{nedefinované}. Databáze a~Prolog naopak stojí na předpokladu uzavřeného světa, kde je vše,
co nevíme, \emph{nepravdivé}.
\end{defbox}

\begin{defbox}[Deskripční logika]
Pracuje se třemi druhy entit: \textbf{koncepty} (třídy), \textbf{role} (vlastnosti, predikáty)
a~\textbf{individua} (objekty). \textbf{Axiom} je logický výraz o~rolích a~konceptech. Zatímco
rámce a~objektové programování popisují třídu \emph{úplně}, axiomy ji popisují jen
\emph{částečně} a~mohou navíc přicházet z~více zdrojů najednou.

Nejde o~jeden systém, ale o~celou rodinu logických systémů podle toho, jaké konstruktory se
povolí. Základní $\mathcal{ALC}$ nabízí negaci, konjunkci a~disjunkci konceptů a~omezené
existenční a~univerzální kvantifikátory. Rozšíření se značí písmeny: hierarchie rolí
($\mathcal{H}$), tranzitivita ($\mathcal{S}$), inverzní role ($\mathcal{I}$), kardinalitní
omezení ($\mathcal{N}$, $\mathcal{Q}$), nominály ($\mathcal{O}$).

Znalostní báze se dělí na dvě části: \textbf{T-Box} je terminologie, tedy axiomy o~konceptech
a~rolích, a~\textbf{A-Box} je asertivní část, tedy fakta o~individuích.
\end{defbox}

\textbf{Profily OWL.} Verze 1 se dodává ve třech odstupňovaných profilech. \emph{OWL~1 Lite}
($\mathcal{SHIF}$) je nejjednodušší a~stojí blízko RDF; má mnoho omezení, aby šel rychle
strojově zpracovat. \emph{OWL~1 DL} ($\mathcal{SHOIN}$) odpovídá deskripční logice, umožňuje
tedy vyjádřit například disjunktnost tříd, a~přitom zůstává úplný a~rozhodnutelný.
\emph{OWL~1 Full} má největší vyjadřovací sílu, jenže mnoho problémů je v~něm nerozhodnutelných.
\emph{OWL~2} ($\mathcal{SROIQ}$) k~tomu přidává tři profily šité na míru typickým úlohám:
\emph{EL} odvozuje v~polynomiálním čase a~používají ho velké lékařské ontologie typu SNOMED,
\emph{QL} cílí na dotazy nad znalostní bází a~napojení na relační databáze a~\emph{RL} zapisuje
axiomy ve tvaru pravidel.

\begin{defbox}[KIF --- \emph{Knowledge Interchange Format}]
Reprezentuje znalosti v~logice prvního řádu, a~to LISPovskou prefixovou syntaxí. Navržen byl
v~rámci projektu DARPA \emph{Knowledge Sharing Effort} jako \textbf{vnitřní jazyk} zpráv, tedy
jazyk jejich obsahu, zatímco KQML byl jazykem \emph{vnějším}, tedy obálkou. Příklady:
\begin{quote}\small\ttfamily
(salary 015-46-3946 widgets 72000)\\
(> (* (width chip1) (length chip1)) (* (width chip2) (length chip2)))\\
(forall (?F ?T) (=> (and (instance ?F Farmer) (instance ?T Tractor)) (likes ?F ?T)))\\
(exists (?F ?T) (and (instance ?F Farmer) (instance ?T Tractor) (likes ?F ?T)))
\end{quote}
\end{defbox}

Přímým předchůdcem OWL byl ještě \textbf{DAML+OIL}, tedy spojení DARPA Agent Markup Language
a~Ontology Interchange Language; opuštěn byl v~roce 2006.

\subsection{Uvažování o~znalostech: epistemická logika}

\noindent\mimoq{slidy modální logiky výslovně vynechávají; okruh ale mluví o~\uv{znalostech
v~MAS}, takže minimum se hodí umět}

Ontologie říká, co znalost znamená. Epistemická logika řeší něco jiného: jak o~znalosti mluvit
zvenčí, tedy jak zapsat, že jeden agent něco ví, a~jak z~toho vyvozovat důsledky pro celou
skupinu.

\begin{defbox}[Epistemická logika a~společná znalost]
Znalost agenta se modeluje modálním operátorem $K_i\varphi$ (\uv{agent $i$ \emph{ví},
že $\varphi$}) s~\textbf{kripkovskou sémantikou možných světů}: $K_i\varphi$ platí
ve světě $w$, právě když $\varphi$ platí ve \emph{všech} světech, které agent $i$
neumí od $w$ odlišit. \textbf{Znalost} se axiomatizuje systémem \textbf{S5} a~klíčový je v~něm
axiom \textbf{T}: $K_i\varphi \Rightarrow \varphi$, tedy co vím, to platí. \textbf{Přesvědčení}
(operátor $B_i$, srov.\ BDI) se naproti tomu axiomatizuje systémem \textbf{KD45}, kde axiom~T
nahrazuje slabší \textbf{D}: $\neg B_i\bot$. Přesvědčení tedy musí být konzistentní, ale
\emph{může být mylné}. Formálně je rozdíl mezi znalostí a~přesvědčením právě axiom~T.

Pro skupinu $G$ se zavádí $E_G\varphi$, tedy že každý v~$G$ ví $\varphi$. Silnější je
\textbf{společná znalost} $C_G\varphi \equiv E_G\varphi \wedge E_G E_G\varphi \wedge \dots$,
tedy že všichni vědí, že všichni vědí, \dots

\textbf{Koordinovaný útok} (problém dvou generálů) ukazuje, proč na tom záleží. Dva generálové
se přes nespolehlivého posla domlouvají na současném útoku, jenže žádný konečný počet potvrzení
nevytvoří společnou znalost: po $k$ zprávách platí nejvýše $E_G^k\varphi$, nikdy $C_G\varphi$.
Poučení pro MAS je nepříjemné a~obecné, totiž že \emph{nespolehlivá komunikace nikdy nevytvoří
společnou znalost}, a~tedy ani dokonalou koordinaci.
\end{defbox}

\subsection{Proč je komunikace agentů jiná}

V~objektovém programování znamená \uv{komunikace} volání metody objektu s~parametry a~volaný
objekt \emph{musí} vyhovět. U~agentů to takhle nefunguje: nikdo nemůže jiného agenta přinutit,
aby něco udělal, ani mu nesmí přímo přepsat vnitřní proměnné. Zbývá jediná cesta, totiž
\textbf{provést komunikační akt} s~jedním ze dvou cílů:
\begin{itemize}
\item vyměnit si s~ostatními informace, nebo
\item \emph{ovlivnit} jiné agenty, aby něco udělali.
\end{itemize}
Ostatní agenti mají vlastní cíle a~je jen na nich, jak s~přijatou informací, žádostí či dotazem
naloží. To je zásadní rozdíl: \textbf{komunikace v~MAS je akce, ne volání}.

\subsection{Teorie řečových aktů}

Pokud je komunikace akce, potřebujeme teorii, která to řekne přesně. Multiagentní systémy si ji
vypůjčily z~filozofie jazyka.

\begin{defbox}[Řečové akty (\emph{speech acts}) --- Austin, 1962]
Některé části užívání jazyka mají charakter \textbf{akcí}, protože mění stav světa podobně jako
akce fyzické. Právě těm se říká řečové akty:
\begin{itemize}
\item \uv{Prohlašuji vás mužem a~ženou.}
\item \uv{Vyhlašuji válku.}
\end{itemize}
Je to \emph{pragmatická} teorie jazyka: zajímá ji, jak se řeč používá k~dosahování cílů.
Typickými performativními slovesy jsou \emph{request}, \emph{inform} a~\emph{promise}.
\end{defbox}

Jeden a~týž výrok se přitom dá popsat na třech úrovních a~pro agenty je použitelná jen jedna
z~nich.

\begin{defbox}[Tři vrstvy řečového aktu]
\begin{itemize}
\item \textbf{Lokuční akt} (\emph{locutionary act}) zachycuje, \emph{co bylo řečeno}, tedy
  samotný výrok. \uv{Udělej mi čaj.}
\item \textbf{Ilokuční akt} (\emph{illocutionary act}) zachycuje, \emph{co bylo míněno}, tedy
  lokuci plus performativní význam (dotaz, žádost, \dots). \uv{Požádal mě o~čaj.}
\item \textbf{Perlokuční akt} (\emph{perlocutionary act}) zachycuje, \emph{co se skutečně
  stalo}, tedy efekt řečového aktu. \uv{Přiměla mě, abych jí uvařil čaj.}
\end{itemize}
Pro MAS je klíčová vrstva \textbf{ilokuční} --- ta se stává obsahem zprávy (performativa).
\end{defbox}

Austin popsal, co řečové akty jsou. \textbf{Searle} k~tomu rozpracoval podmínky, za nichž je
řečový akt \uv{zdařilý}. Pro akt \emph{MLUVČÍ požádá POSLUCHAČE o~akci} vypadají takto:
\begin{itemize}
\item \textbf{Standardní vstupně-výstupní podmínky} zajišťují, že komunikace vůbec proběhla:
  posluchač slyší, nejde o~scénu ve filmu\dots
\item \textbf{Předpoklady} shrnují, co musí platit, aby mluvčí mohl tento akt zvolit. Posluchač
  musí být schopen akci provést, mluvčí tomu musí věřit a~nesmí být zřejmé, že by ji posluchač
  provedl i~bez požádání.
\item \textbf{Upřímnost} (\emph{sincerity}) znamená, že mluvčí opravdu chce, aby se akce
  provedla.
\end{itemize}

Searle rovněž setřídil řečové akty samotné.

\begin{defbox}[Searlovy kategorie řečových aktů]
Starší dělení bylo jen výčtem druhů: \emph{request}, \emph{advice}, \emph{statement},
\emph{promise}, \dots

Novější, a~dnes standardní, dělení rozlišuje pět tříd:
\begin{itemize}
\item \textbf{Asertiva} (\emph{assertives}, reprezentativa) informují posluchače
  (\emph{inform}).
\item \textbf{Direktiva} (\emph{directives}) žádají posluchače o~akci (\emph{request}).
\item \textbf{Komisiva} (\emph{commissives}) zavazují mluvčího k~něčemu do budoucna
  (\emph{promise}).
\item \textbf{Expresiva} (\emph{expressives}) vyjadřují mentální stav či emoci mluvčího
  (\uv{děkuji}).
\item \textbf{Deklarace} (\emph{declarations}) samy mění stav věcí (válka, sňatek).
\end{itemize}
Každý řečový akt má vždy dvě složky: \textbf{performativní sloveso} (request, query, inform)
a~\textbf{propoziční obsah} (\uv{okno je zavřené}).
\end{defbox}

Aby se s~řečovými akty dalo počítat uvnitř agenta, musely dostat sémantiku srozumitelnou
plánovači.

\begin{defbox}[Plánovací teorie řečových aktů (Cohen, Perrault, 1979)]
\uv{\dots modelování [řečových aktů] v~plánovacím systému jako operátorů definovaných\dots{}
pomocí přesvědčení a~cílů mluvčích a~posluchačů. S~řečovými akty se tak zachází stejně jako
s~fyzickými akcemi.}

Sémantika řečových aktů je tedy dána v~duchu \textbf{STRIPS}, totiž trojicí
předpoklady -- delete -- add. Rozdíl je jen v~tom, že se v~ní používají modální operátory pro
\emph{beliefs}, \emph{abilities} a~\emph{wants}.
\end{defbox}

\noindent\textbf{Příklad: Request a~Inform.}
\begin{center}\small
\begin{tabular}{p{0.16\textwidth} p{0.38\textwidth} p{0.38\textwidth}}
\toprule
& \textbf{Request(S,H,A)} & \textbf{Inform(S,H,$\varphi$)} \\
\midrule
Předpoklad \emph{cando}
  & $(S\ \mathrm{believe}\ (H\ \mathrm{cando}\ A))\ \wedge$ \newline
    $(S\ \mathrm{believe}\ (H\ \mathrm{believe}\ (H\ \mathrm{cando}\ A)))$
  & $(S\ \mathrm{believe}\ \varphi)$ \\
Předpoklad \emph{want}
  & $(S\ \mathrm{believe}\ (S\ \mathrm{want}\ \mathrm{requestInstance}))$
  & $(S\ \mathrm{believe}\ (S\ \mathrm{want}\ \mathrm{informInstance}))$ \\
Efekt
  & $(H\ \mathrm{believe}\ (S\ \mathrm{believe}\ (S\ \mathrm{want}\ A)))$
  & $(H\ \mathrm{believe}\ (S\ \mathrm{believe}\ \varphi))$ \\
\bottomrule
\end{tabular}
\end{center}
Za pozornost stojí, že efektem \emph{Inform} \emph{není} \uv{posluchač uvěří $\varphi$}, ale jen
\uv{posluchač uvěří, že mluvčí věří $\varphi$}. Autonomní agent totiž není povinen věřit tomu,
co mu kdo řekne, a~právě v~tom se komunikační akt liší od volání metody.

\subsection{KQML a~FIPA-ACL}
\label{ssec:kqml-acl}

Teorie řečových aktů se do praxe promítla dvěma jazyky. Starší z~nich je KQML.

\begin{defbox}[KQML --- \emph{Knowledge Query and Manipulation Language}]
Vznikl v~90.~letech v~rámci projektu DARPA \emph{Knowledge Sharing Effort} (KSE) společně s~KIF.
Oba jazyky si rozdělily role:
\begin{itemize}
\item \textbf{KQML} je \emph{vnější} komunikační jazyk, tedy \uv{obálka dopisu}; nese
  \emph{ilokuční} část zprávy.
\item \textbf{KIF} je jazyk \emph{vnitřní}: propoziční obsah zprávy a~reprezentace znalostí.
\end{itemize}
Zpráva má podobu seznamu, v~jehož čele stojí \emph{performativum} a~za ním následují pojmenované
parametry:
\begin{quote}\small\ttfamily
(ask-one\\
\hspace*{1em}:content (PRICE IBM ?price)\\
\hspace*{1em}:receiver stock-server\\
\hspace*{1em}:language LPROLOG\\
\hspace*{1em}:ontology NYSE-TICKS)
\end{quote}
\textbf{Parametry:} content, force, reply-with, in-reply-to, sender, receiver, language, ontology.

\textbf{Performativa:} achieve, advertise, ask-about, ask-one, ask-all, ask-if,
break, sorry, error, broadcast, forward, recruit-one, recruit-all, reply, subscribe, tell, \dots
\end{defbox}

\textbf{Kritika KQML} se soustředí do čtyř bodů. Chybělo performativum \emph{commit}, tedy závazek.
Sémantika nebyla přesně definovaná. Množina performativ byla nesystematická a~zbytečně velká.
A~transportní mechanismus vůbec nebyl standardizován. Na tuto kritiku odpovídá druhý jazyk.

\begin{defbox}[FIPA-ACL --- \emph{Agent Communication Language}]
FIPA (\emph{Foundation for Intelligent Physical Agents}, dnes standardizační komise IEEE)
navrhla ACL jako \textbf{zjednodušení a~zpřesnění KQML}: má menší a~systematičtější množinu
performativ a~\emph{formálně definovanou sémantiku}. Praktickou implementací je platforma JADE.

\textbf{Struktura zprávy} čítá 13 polí, totiž performativum a~12 parametrů; povinné je přitom
jediné, \texttt{performative}:
\begin{center}\small
\begin{tabular}{ll}
\toprule
\textbf{Pole} & \textbf{Význam} \\
\midrule
\texttt{performative} & typ komunikačního aktu (ilokuční síla) --- \emph{povinné} \\
\texttt{sender}, \texttt{receiver}, \texttt{reply-to} & identifikátory agentů (AID) \\
\texttt{content} & vlastní obsah zprávy \\
\texttt{language} & jazyk obsahu (SL, KIF, \dots) \\
\texttt{encoding}, \texttt{ontology} & kódování a~ontologie pro interpretaci obsahu \\
\texttt{protocol} & interakční protokol, jehož je zpráva součástí \\
\texttt{conversation-id}, \texttt{reply-with}, \texttt{in-reply-to} & párování zpráv do konverzace \\
\texttt{reply-by} & termín, do kdy se očekává odpověď \\
\bottomrule
\end{tabular}
\end{center}
\begin{quote}\small\ttfamily
(inform\\
\hspace*{1em}:sender agent1\\
\hspace*{1em}:receiver agent2\\
\hspace*{1em}:content (price good2 150)\\
\hspace*{1em}:language sl\\
\hspace*{1em}:ontology hpl-auction)
\end{quote}
\end{defbox}

\textbf{Performativa FIPA-ACL.} Je jich celkem 22 a~přehledně se dělí podle toho, k~čemu
slouží; níže jsou ta nejdůležitější:
\begin{center}\small
\begin{tabular}{p{0.24\textwidth} p{0.70\textwidth}}
\toprule
\textbf{Skupina} & \textbf{Performativa} \\
\midrule
předávání informace & \texttt{inform}, \texttt{inform-if}, \texttt{inform-ref},
  \texttt{confirm}, \texttt{disconfirm} \\
žádost o~informaci & \texttt{query-if}, \texttt{query-ref}, \texttt{subscribe} \\
žádost o~akci & \texttt{request}, \texttt{request-when}, \texttt{request-whenever},
  \texttt{cancel} \\
vyjednávání & \texttt{cfp} (\emph{call for proposals}), \texttt{propose},
  \texttt{accept-proposal}, \texttt{reject-proposal} \\
provádění akcí & \texttt{agree}, \texttt{refuse}, \texttt{failure} \\
chyby, směrování & \texttt{not-understood}, \texttt{proxy}, \texttt{propagate} \\
\bottomrule
\end{tabular}
\end{center}

\noindent\mimoq{formální sémantika SL ve slidech není, performativa ano}

Tím, čím se FIPA-ACL nejvíc odlišuje od KQML, je právě sémantika: každé performativum má
předepsáno, kdy se smí použít a~proč se používá.

\begin{defbox}[Sémantika FIPA-ACL: FP a~RE]
Sémantika je definována v~jazyce SL (\emph{Semantic Language}), což je modální logika
s~operátory $B_i\varphi$ (agent $i$ věří $\varphi$), $U_i\varphi$ (je nejistý),
$I_i\varphi$ (má záměr), $\mathrm{Feasible}$ a~$\mathrm{Done}$. Ke každému performativu jsou
dány dvě věci:
\begin{itemize}
\item \textbf{FP} (\emph{feasibility precondition}) říká, co musí platit, aby odesílatel mohl
  akt provést;
\item \textbf{RE} (\emph{rational effect}) je efekt, kvůli němuž akt provádí. Příjemce
  \emph{není povinen} RE naplnit, protože je autonomní --- RE je jen důvod odesílatele.
\end{itemize}
Například pro $\langle i, \mathrm{inform}(j,\varphi)\rangle$:
\[
\text{FP: } B_i\varphi \wedge \neg B_i\big(\mathit{Bif}_j\varphi \vee \mathit{Uif}_j\varphi\big),
\qquad
\text{RE: } B_j\varphi .
\]
Čte se to takto: \uv{já $\varphi$ věřím a~nevěřím, že už $j$ o~$\varphi$ ví (že by věděl, zda
platí, či neplatí --- $\mathit{Bif}$ --- ani že by o~tom měl mínění s~nejistotou ---
$\mathit{Uif}$); chci, aby $j$ uvěřil $\varphi$.}
Obdobně pro $\langle i, \mathrm{request}(j,\alpha)\rangle$:
$\text{FP: } \mathit{FP}(\alpha)[i \backslash j] \wedge B_i \mathrm{Agent}(j,\alpha)
\wedge \neg B_i I_j \mathrm{Done}(\alpha)$, $\text{RE: } \mathrm{Done}(\alpha)$.

\textbf{Praktický problém} je v~tom, že tato sémantika je neverifikovatelná: zvenčí se nedá
ověřit, že agent skutečně věří tomu, co tvrdí (\emph{sincerity assumption}). Je to známá slabina
mentalistické sémantiky ACL a~alternativou k~ní jsou \emph{sociální} sémantiky založené na
veřejných závazcích (\emph{commitments}).
\end{defbox}

\subsection{Interakční protokoly}

Jednotlivá zpráva sama o~sobě nestačí, protože agenti spolu vedou celé \textbf{konverzace}
s~předepsanou strukturou. Interakční protokol (\emph{interaction protocol}) je vzor přípustných
posloupností zpráv; FIPA jich standardizuje kolem tuctu ve specifikaci FIPA-IP a~zapisuje je
notací AUML.

\begin{center}\small
\begin{tabular}{p{0.26\textwidth} p{0.68\textwidth}}
\toprule
\textbf{Protokol} & \textbf{Průběh} \\
\midrule
\textbf{FIPA-Request}
  & $I \to P$: \texttt{request}. $P \to I$: \texttt{refuse} (konec), nebo \texttt{agree};
    po vykonání \texttt{inform-done} / \texttt{inform-result}, případně \texttt{failure}.
    Odpovědí může být i~\texttt{not-understood}. \\
\textbf{FIPA-Query}
  & $I \to P$: \texttt{query-if} / \texttt{query-ref}. $P$: \texttt{refuse} nebo \texttt{agree},
    dále \texttt{inform} s~výsledkem nebo \texttt{failure}. \\
\textbf{FIPA-Contract-Net}
  & viz níže: \texttt{cfp} $\to$ $n\times$(\texttt{propose} / \texttt{refuse}) $\to$
    \texttt{accept-proposal} vítězi + \texttt{reject-proposal} ostatním $\to$
    \texttt{inform-done} / \texttt{failure}. \\
\textbf{FIPA-Iterated-Contract-Net}
  & jako CNET, ale iniciátor může po přijetí nabídek vyhlásit \emph{další kolo}
    \texttt{cfp} (vyjednávání o~lepších podmínkách). \\
\textbf{FIPA-Subscribe}
  & $I$: \texttt{subscribe}; $P$ posílá \texttt{inform} pokaždé, když se sledovaná
    hodnota změní, dokud nepřijde \texttt{cancel}. \\
\textbf{FIPA-Brokering / Recruiting}
  & $I$ požádá zprostředkovatele (\emph{broker}), ten najde vhodné poskytovatele, deleguje
    na ně požadavek a~výsledky vrátí (příp.\ je nechá odpovědět přímo). \\
\textbf{Aukční protokoly}
  & FIPA-English-Auction, FIPA-Dutch-Auction (viz odst.~\ref{ssec:aukce}). \\
\bottomrule
\end{tabular}
\end{center}

\textbf{Poznámka k~návrhu.} Protokol určuje, jaké zprávy jsou v~daném stavu konverzace legální.
Tím vnáší \emph{předvídatelnost} do jinak autonomního chování a~zároveň dovoluje implementovat
agenta jako konečný automat nad stavy konverzace. Přesně to dělá JADE svými třídami typu
\texttt{AchieveREInitiator} nebo \texttt{ContractNetResponder}.

\subsection{Shrnutí ke~zkoušce}

\begin{itemize}
\item Ontologie je explicitní formalizovaný popis části reality, tedy glosář a~tezaurus
  dohromady; standardně se cituje Gruberova definice. Formalismus tvoří třídy, instance,
  vlastnosti, relace (zejména tranzitivní \emph{is-a}) a~fakta, přičemž ontologie spolu s~fakty
  dává znalostní bázi.
\item Formalizace má škálu od pouhého slovníku až po libovolná logická omezení; podle rozsahu se
  ontologie dělí na aplikační, doménové a~horní.
\item RDF stojí na trojicích \emph{subjekt--predikát--objekt}, je jednoduché a~rychlé; RDFS
  k~němu přidává třídy.
\item OWL je postaveno na deskripční logice, tedy rozhodnutelném fragmentu FOL, platí v~něm
  \emph{open world assumption} a~znalostní báze se dělí na T-Box a~A-Box. Verze 1 má profily
  Lite, DL a~Full ($\mathcal{SHIF}$, $\mathcal{SHOIN}$), OWL~2 je $\mathcal{SROIQ}$ a~přidává
  profily EL, QL a~RL.
\item KIF je FOL v~LISPovské syntaxi a~tvoří obsah zprávy, KQML je obálka.
\item Epistemická logika zapisuje znalost operátorem $K_i\varphi$ nad možnými světy; znalost se
  axiomatizuje systémem S5 s~axiomem~T, přesvědčení systémem KD45. Společná znalost $C_G\varphi$
  nespolehlivou komunikací nevznikne, což ukazuje koordinovaný útok.
\item Komunikace agentů je \emph{akce}, ne volání metody, takže příjemce není povinen vyhovět.
\item Austin zavedl řečové akty a~jejich tři vrstvy, lokuci, ilokuci a~perlokuci. Searle k~tomu
  přidal podmínky zdařilosti a~pět kategorií (asertiva, direktiva, komisiva, expresiva,
  deklarace); každý akt se skládá z~performativního slovesa a~propozičního obsahu.
\item Cohen a~Perrault dali řečovým aktům sémantiku ve stylu STRIPS, tedy předpoklady a~efekt,
  jen s~modálními operátory. Request(S,H,A) a~Inform(S,H,$\varphi$) patří k~tomu, co se
  u~zkoušky píše zpaměti; efektem Inform je \emph{$H$ věří, že $S$ věří $\varphi$}.
\item KQML tvoří obálku a~KIF obsah; k~tomu patří parametry, performativa a~kritika KQML
  (chybí commit, nejasná sémantika).
\item FIPA-ACL má danou strukturu zprávy (performative, sender/receiver, content, language,
  ontology, protocol, conversation-id, reply-by), 22 performativ a~sémantiku přes FP a~RE
  v~jazyce SL; její slabinou je neverifikovatelnost.
\item Z~protokolů se jmenují FIPA-Request, FIPA-Query, Contract Net i~jeho iterovaná varianta,
  Subscribe, Brokering a~aukční protokoly.
\end{itemize}

%=====================================================================

\section{Distribuované řešení problémů, kooperace, teorie her a~aukce}
\label{sec:kooperace}
%=====================================================================

Nejobsáhlejší věta okruhu spojuje dvě situace, které spolu na první pohled souvisejí, ve
skutečnosti však vedou k~docela jiné matematice. Kapitola má proto dvě poloviny. V~té první
(odst.\ 5.1--5.6) agenti \emph{sdílejí cíl} a~zbývá vyřešit jedinou otázku, totiž jak si
rozdělit práci; to je \emph{distribuované řešení problémů} a~\emph{kooperace}. Ve druhé
polovině (odst.\ 5.7--5.15) sleduje každý agent \emph{vlastní zájmy}, a~tak je potřeba teorie
her s~Nashovými ekvilibrii a~Paretovou efektivitou, spolu s~mechanismy odolnými vůči
manipulaci, tedy alokací zdrojů a~aukcemi.
\subsection{Koherence a~koordinace}

Agenti mají různé cíle, jednají autonomně, pracují v~čase a~rozhodují se za běhu. Jejich
součinnost proto nelze naplánovat dopředu, systém musí zvládnout \emph{dynamickou koordinaci}.
Sdílejí při ní dvě věci: \textbf{úlohy} (\emph{task sharing}) a~\textbf{informace}
(\emph{result sharing}).

\begin{defbox}[Koherence a~koordinace]
\textbf{Koherence} (\emph{coherence}) říká, jak dobře funguje systém \emph{jako celek}: jakou
kvalitu má nalezené řešení, jak efektivně se využívají zdroje a~jak systém degraduje při
poruchách.

\textbf{Koordinace} (\emph{coordination}) měří něco jiného, totiž jak dobře se agentům daří
\emph{minimalizovat režii} spojenou se synchronizací. Dobře koordinovaný systém se vyhne
zbytečně duplicitní práci, konfliktům o~zdroje i~zbytečné komunikaci.

Obojí spolu nesouvisí tak těsně, jak by se zdálo. Systém může být koherentní i~při špatné
koordinaci: agenti tutéž práci udělají třikrát, výsledek je ale správný. Zaplatí se to
plýtváním.
\end{defbox}

\begin{defbox}[CDPS --- \emph{Cooperative Distributed Problem Solving}]
Pojem pochází od Lessera a~kol. z~80.~let. Označuje kooperaci jednotlivých agentů při řešení
problému, který přesahuje jejich individuální schopnosti, ať už jde o~informace, zdroje, nebo
výpočetní kapacitu.

Celý přístup stojí na jednom předpokladu: agenti \emph{implicitně sdílejí společný cíl}
a~konflikty mezi nimi nevznikají. Této vlastnosti se říká \textbf{benevolence}. Měřítkem
úspěchu pak není prospěch jednotlivce, nýbrž \emph{výkon systému jako celku}, takže agent celku
pomůže i~tehdy, když je to pro něj samotného nevýhodné. Benevolence \emph{enormně zjednodušuje}
návrh systému.
\end{defbox}

\begin{center}\small
\begin{tabular}{p{0.14\textwidth} p{0.38\textwidth} p{0.40\textwidth}}
\toprule
& \textbf{Charakteristika} & \textbf{Rozdíl} \\
\midrule
\textbf{PPS} \newline (\emph{parallel problem solving})
  & Bond, Gasser, 80.~léta; důraz na paralelní řešení, homogenní a~jednoduché procesory
  & agenti nejsou autonomní, jde v~podstatě o~paralelní algoritmus \\
\textbf{CDPS}
  & heterogenní agenti se sofistikovanými schopnostmi, sdílený cíl, benevolence
  & konflikty se neřeší, protože nevznikají \\
\textbf{MAS}
  & společnost agentů s~\emph{vlastními} cíli, které \emph{nesdílejí}
  & je nutné řešit: proč a~jak kooperovat, jak identifikovat a~řešit konflikty, jak
    vyjednávat a~smlouvat \\
\bottomrule
\end{tabular}
\end{center}

\subsection{Sdílení úloh a~sdílení výsledků}

\textbf{Obecný postup CDPS} probíhá ve třech fázích:
\begin{enumerate}
\item \textbf{Dekompozice problému} je hierarchická a~rekurzivní. Ptáme se přitom na dvě věci:
  \emph{jak} problém rozložit a~\emph{kdo} rozklad provede. Krajní variantou je model ACTORS,
  v~němž pro každý podproblém vznikne nový agent, a~to až na úroveň jednotlivých instrukcí.
\item \textbf{Řešení podproblémů} obvykle neprobíhá izolovaně: agenti si během něj typicky
  sdílejí informace a~mohou potřebovat synchronizovat akce.
\item \textbf{Syntéza řešení} je opět hierarchická, skládají se při ní dílčí výsledky.
\end{enumerate}

Obě sdílené věci fungují jinak. \textbf{Sdílení úloh} vyžaduje \emph{dohodu} agentů o~tom, kdo
co udělá, a~typickým mechanismem takové dohody je Contract Net. \textbf{Sdílení výsledků}
naproti tomu může být buď \emph{proaktivní}, kdy agent pošle výsledek, o~němž si myslí, že se
bude hodit, nebo \emph{reaktivní}, kdy jej pošle až na vyžádání.

\begin{defbox}[Přínosy sdílení výsledků]
\begin{itemize}
\item \textbf{Jistota} (\emph{confidence}): nezávislá řešení téhož problému lze porovnat,
  a~shodnou-li se, důvěra ve výsledek roste.
\item \textbf{Úplnost} (\emph{completeness}): agenti sdílejí své \emph{lokální} pohledy a~tím
  vzniká globálnější představa o~problému.
\item \textbf{Přesnost} (\emph{precision}): sdílení může zpřesnit celkové řešení.
\item \textbf{Včasnost} (\emph{timeliness}): zjevná výhoda distribuce, totiž zkrácení času.
\end{itemize}
\end{defbox}
\subsection{Contract Net}

\begin{defbox}[Contract Net Protocol (CNET)]
Smith a~Davis, 1977--1980. Protokol staví na metafoře \textbf{sdílení úloh pomocí smluv} podle
vzoru trhu: agent v~nouzi vystupuje jako \emph{manažer} (\emph{manager}), ostatní jako
\emph{dodavatelé} (\emph{contractors}). Probíhá ve čtyřech fázích:
\begin{enumerate}
\item \textbf{Rozpoznání} (\emph{recognition}): agent zjistí, že má problém, který sám
  nezvládne. Buď mu chybí schopnost, zdroje či informace, nebo by řešení dopadlo hůř, než kdyby
  úlohu předal.
\item \textbf{Ohlášení} (\emph{announcement}): rozešle, typicky broadcastem, oznámení úlohy.
  To obsahuje \emph{popis úlohy}, který může být i~spustitelný, dále \emph{omezení} na termín,
  cenu a~kvalitu a~konečně \emph{meta-informace} o~tom, kdo se smí hlásit.
\item \textbf{Nabídky} (\emph{bidding}): příjemci zváží, zda mají zájem a~schopnosti, a~pošlou
  nabídku (\emph{tender}) s~cenou, časem a~kvalitou.
\item \textbf{Přidělení a~realizace} (\emph{awarding, expediting}): manažer vybere vítěze
  a~uzavře s~ním smlouvu, vítěz úlohu vykoná a~vrátí výsledek.
\end{enumerate}
\textbf{Vlastnosti:} protokol je jednoduchý, decentralizovaný a~robustní vůči výpadkům.
Dodavatel se navíc může sám stát manažerem podúlohy, čímž vznikají \emph{hierarchické kaskády
subkontraktů}. Zátěž se tak rozděluje dynamicky. Jde o~nejčastěji implementovaný koordinační
mechanismus vůbec: FIPA jej standardizuje jako FIPA-Contract-Net a~JADE jej nabízí jako hotové
chování.

\textbf{Limity:} broadcast se při mnoha agentech prodraží a~celý protokol předpokládá
benevolenci, tedy že dodavatel nelže o~svých schopnostech. Jakmile tento předpoklad padne, je
nutné sáhnout po aukčních mechanismech odolných vůči manipulaci (odst.~\ref{ssec:aukce}).
\end{defbox}

\subsection{Systémy s~tabulí (blackboard)}

\begin{defbox}[Blackboard system (BBS)]
Historicky vůbec první schéma pro kooperativní řešení problémů; použil je systém HEARSAY-II pro
rozpoznávání řeči v~70.~letech. Výsledky se sdílejí přes společnou datovou strukturu, které se
říká \textbf{tabule} (\emph{blackboard}).

Metafora je názorná. Kolem tabule sedí několik expertů a~každý z~nich na ni může psát i~z~ní
číst. Na tabuli se dynamicky objevují úlohy, a~jakmile některý expert uvidí úlohu, kterou umí
řešit, zapíše na tabuli částečné řešení. Takto se pokračuje, dokud se na tabuli neobjeví řešení
celé úlohy.

\textbf{Role:}
\begin{itemize}
\item \textbf{Tabule} je sdílená paměť, typicky strukturovaná do několika \emph{úrovní
  abstrakce}, které odpovídají hierarchii cílů a~akcí. Zásadní volbou je formalismus
  reprezentace informace.
\item \textbf{Experti} (\emph{knowledge sources}) jsou agenti řešící konkrétní typ úlohy.
  Reagují na cíle na tabuli, svůj cíl si z~tabule přebírají a~po vybrání vykonají akci. Každý
  z~nich má seznam schopností a~svou efektivitu.
\item \textbf{Arbitr} (\emph{arbiter}, \emph{control component}) vybírá, který expert smí
  k~tabuli. Může být čistě reaktivní, nebo může uvažovat plány a~maximalizovat očekávanou
  utilitu. Odpovídá za řešení problému na vyšší úrovni.
\end{itemize}
Nad tabulí je \textbf{nutné vzájemné vyloučení}, a~právě tam vzniká úzké hrdlo systému.
\end{defbox}

\textbf{Klady a~zápory BBS:}
\begin{itemize}
\item[$+$] Poskytuje jednoduchý mechanismus koordinace a~kooperace, který pokrývá sdílení úloh
  i~sdílení výsledků.
\item[$+$] Experti o~sobě \emph{nemusí vědět} a~přesto spolupracují (\emph{loose coupling});
  přidání dalšího experta nevyžaduje změnu ostatních.
\item[$+$] Zprávy na tabuli lze přepisovat, a~tím delegovat úlohy, vytvářet podúlohy nebo měnit
  experty.
\item[$-$] Všichni agenti musí mít přístup k~tabuli a~sdílet stejnou architekturu, a~kolem
  tabule bývá \uv{narváno}; částečně to řeší distribuované hašovací tabulky.
\end{itemize}

\medskip
\noindent\textbf{Příklad: BBWar.} Tabule je zde hašovací tabulka, která mapuje požadované
schopnosti na úlohy, a~publikují se na ní \uv{otevřené mise}. Experti tvoří hierarchii: agent
\emph{velitel} hledá úlohy typu \uv{dobyj město} (\texttt{ATTACK-CITY}) a~převádí je na
několik úloh \uv{zaútoč na pozici} (\texttt{ATTACK-LOCATION}). Vojáci různých typů si pak
z~tabule vybírají mise, na které mají schopnosti.

\medskip
\noindent\textbf{Příklad: FELINE} (Wooldridge, Jennings, 1990) je distribuovaný expertní systém
z~doby před KQML a~ACL. Zajišťuje sdílení znalostí a~distribuci podúloh. Každý agent je
pravidlový systém, který si udržuje dvě položky: \emph{Skills}, tedy \uv{umím dokázat/vyvrátit
následující\dots}, a~\emph{Interests}, tedy \uv{zajímá mě, zda následující platí\dots}.
Komunikace má tvar trojice (odesílatel, příjemce, obsah), kde obsahem je hypotéza doplněná
řečovým aktem (request, response, inform).

\subsection{Distribuované omezující podmínky (DisCSP a~DCOP)}

\noindent\mimoq{není ve slidech; okruh ale mluví o~\uv{distribuovaném řešení problémů},
kam formální větev patří}

\begin{defbox}[DisCSP a~DCOP]
\textbf{Distribuovaný CSP} vzniká tak, že se proměnné $x_1,\dots,x_n$ s~doménami $D_i$ rozdělí
mezi agenty, typicky po jedné proměnné na agenta, a~omezení propojují proměnné \emph{různých}
agentů. Hledá se přiřazení splňující všechna omezení. Žádný agent přitom nevidí celý problém
a~komunikuje jen se sousedy v~grafu omezení.

\textbf{DCOP} je optimalizační varianta téhož. Omezení jsou \emph{měkká} a~zadaná funkcemi ceny
$f_{ij}(x_i,x_j)$; maximalizuje se $\sum_{ij} f_{ij}$, tedy sociální blahobyt.

\textbf{Algoritmy:} nejznámější je \emph{ABT} (asynchronní backtracking, Yokoo). Agenti jsou
v~něm uspořádáni a~posílají si dva druhy zpráv: \texttt{ok?} se svým přiřazením směrem níž
a~\texttt{nogood} zpět nahoru, jakmile nastane konflikt. Dále sem patří \emph{ADOPT},
\emph{DPOP} a~\emph{Max-Sum}. Typickými aplikacemi jsou rozvrhování schůzek, přidělování úkolů
senzorům a~alokace frekvencí.
\end{defbox}
\subsection{Interakce agentů v~prostředí}

Agenti si mohou navzájem pomáhat i~překážet, aniž by spolu vůbec \emph{komunikovali}. Stačí, že
ovlivňují prostředí:
\begin{itemize}
\item Roboti přenesou cihlu jen tehdy, když na ni tlačí ze dvou stran současně. To je pozitivní
  interakce a~vyžaduje synchronizaci.
\item Roboti se natlačí ke dveřím a~pak je nemohou otevřít. To je negativní interakce
  a~vyžaduje koordinaci.
\end{itemize}
Agenti mohou vytvářet společenství, hierarchie i~nepřátelství. \textbf{Je nutné znát typy
interakcí} v~konkrétním MAS, protože bez této znalosti nelze navrhnout efektivní řídicí
mechanismy. Nejjednodušším modelovým případem, kterým se zabývá následující oddíl, je
interakce dvou racionálních (sobeckých) agentů v~prostředí připomínajícím hru.
\subsection{Sobecký agent}

Proč by měl být agent upřímný ohledně svých schopností? A~proč by měl přidělenou úlohu vůbec
dokončit? Je-li systém homogenní, tedy celý navržený námi, je benevolence dobrá strategie.
Většinou ale systém obsahuje agenty s~různými zájmy a~mezi společným cílem a~cíli jednotlivců
vzniká konflikt. V~řízení letového provozu chtějí bezpečnost všichni, jenže každá aerolinka
chce přistát první. Někdy navíc ani není zřejmé, co je společný zájem. V~takové situaci je lépe
uvažovat \textbf{sobecké} (\emph{self-interested}) agenty.

\begin{defbox}[Sobecký agent]
Množina možných výsledků $O = \{o_1, o_2, \dots\}$ je \emph{společná} pro všechny agenty,
preference už však nikoli: každý agent $i$ má vlastní \textbf{utilitní funkci}
$u_i : O \to \mathbb{R}$, která výsledky uspořádá.

\textbf{Strategické myšlení} pak zní: maximalizuj svou utilitu za předpokladu, že všichni
ostatní jednají rovněž racionálně. Společnou utilitu tento postup \emph{nemaximalizuje}, zato
je \emph{robustní}.

\emph{Poznámka:} peníze nejsou pro člověka dobrou utilitou. Utilitní funkce peněz je nelineární
a~liší se pro bohaté a~chudé.
\end{defbox}

\subsection{Hra v~normální formě}

\begin{defbox}[Hra dvou hráčů v~normální formě]
Máme agenty $i$ a~$j$ s~množinami akcí $A_i$, $A_j$. Výsledek není v~moci ani jednoho z~nich,
protože jej určují obě volby najednou: $e : A_i \times A_j \to O$. Agent volí \textbf{strategii}
$s_i$, a~to ve dvou možných podobách:
\begin{itemize}
\item \textbf{ryzí} (\emph{pure}) strategie je deterministická volba jedné akce;
\item \textbf{smíšená} (\emph{mixed}) strategie je rozdělení pravděpodobnosti nad ryzími
  strategiemi, tedy náhodný výběr.
\end{itemize}
Hru zapisujeme \textbf{výplatní maticí} (\emph{payoff matrix}), v~níž řádky odpovídají
strategiím hráče $i$, sloupce strategiím hráče $j$ a~každá buňka obsahuje dvojici $u_i / u_j$.

Obecně je hra v~normální formě trojice $\langle N, (A_i)_{i \in N}, (u_i)_{i \in N}\rangle$,
kde $N$ je množina hráčů.
\end{defbox}

\begin{defbox}[Dominance]
Strategie $s_i$ \textbf{silně dominuje} strategii $s_i'$, pokud dává agentovi $i$
\emph{striktně} lepší výsledek proti \emph{všem} strategiím protihráče. \textbf{Slabě dominuje},
pokud dává výsledek alespoň stejně dobrý proti všem a~striktně lepší alespoň proti jedné; tak
dominanci definuje i~přednáška. Strategie je \textbf{dominantní}, pokud dominuje všem ostatním
strategiím agenta $i$.

Rozdíl mezi oběma pojmy není akademický: pravdomluvnost ve Vickreyově aukci je jen \emph{slabě}
dominantní (odst.~\ref{ssec:aukce}).

Racionální agent hraje dominantní strategii, existuje-li nějaká. Na tom stojí základní řešicí
technika, \emph{iterované odstraňování dominovaných strategií} (IESDS): dominovanou strategii
nikdy nikdo nezahraje, lze ji tedy z~matice vyškrtnout, což může učinit dominovanými další
strategie. Odstraňujeme-li \emph{silně} dominované strategie, na pořadí nezáleží a~výsledek je
jednoznačný. U~\emph{slabě} dominovaných strategií na pořadí \textbf{záleží} a~lze přijít
o~některá ekvilibria.
\end{defbox}

\begin{defbox}[Nashovo ekvilibrium]
Dvojice strategií $(s_1^*, s_2^*)$ je v~\textbf{Nashově ekvilibriu}, pokud platí obojí:
\begin{itemize}
\item hraje-li agent $i$ strategii $s_1^*$, je pro agenta $j$ nejlepší hrát $s_2^*$;
\item hraje-li agent $j$ strategii $s_2^*$, je pro agenta $i$ nejlepší hrát $s_1^*$.
\end{itemize}
Strategie $s_1^*$ a~$s_2^*$ jsou tedy \emph{vzájemně nejlepší odpovědi} (\emph{mutual best
responses}). Obecně pro $n$ hráčů je profil $s^* = (s_1^*,\dots,s_n^*)$ Nashovým ekvilibriem,
pokud pro každého hráče $i$ a~každou jeho strategii $s_i$ platí
\[
u_i(s_i^*, s_{-i}^*) \ge u_i(s_i, s_{-i}^*),
\]
kde $s_{-i}$ značí strategie všech ostatních. Jinými slovy: \textbf{žádný hráč nemá jednostranný
důvod se odchýlit.}

Naivní hledání ekvilibrií pro $n$ agentů a~$m$ strategií vyžaduje projít $m^n$ profilů. Uvedená
definice se týká \emph{ryzích} strategií, a~v~nich ekvilibrium existovat nemusí: hra
kámen--nůžky--papír žádné nemá, jiné hry jich naopak mají několik.
\end{defbox}

\begin{veta}[Nashova věta, 1950]
Každá hra s~konečným počtem hráčů a~konečným počtem strategií má alespoň jedno Nashovo
ekvilibrium \emph{ve smíšených strategiích}.
\end{veta}

\textbf{Výpočetní složitost.} Hledání NE je \emph{totální} vyhledávací problém: řešení vždy
existuje a~jde jen o~to najít ho. Od roku 2006 víme, že tento problém je \textbf{PPAD-úplný},
velmi zjednodušeně řečeno \uv{o~něco méně beznadějné než NP-úplnost}. NP-úplnost pro problém,
který má vždy řešení, ostatně ani nedává smysl. Pro tři a~více hráčů to dokázali Daskalakis,
Goldberg a~Papadimitriou; těsně poté Chen a~Deng doplnili i~zbývající a~nejtěžší případ
\emph{dvou} hráčů.

\begin{defbox}[Paretova efektivita a~sociální blahobyt]
Profil strategií je \textbf{Pareto-optimální} (\emph{Pareto efficient}), pokud neexistuje jiný
profil, který by zlepšil výsledek jednoho agenta, aniž by zhoršil výsledek některého jiného.
Řešení, které Pareto-optimální není, lze tedy zlepšit \uv{zadarmo}.

\textbf{Sociální blahobyt} (\emph{social welfare}) profilu je součet utilit všech agentů,
$\mathit{sw}(o) = \sum_{i \in N} u_i(o)$. Maximalizovat sociální blahobyt dává smysl
v~kooperativních scénářích, tedy když agenti patří do jednoho týmu, mají jednoho vlastníka nebo
řeší jednu úlohu; čím homogennější systém, tím lépe. Maximum sociálního blahobytu je vždy
Pareto-optimální, obráceně to však neplatí.
\end{defbox}
\subsection{Klasické hry}

\begin{defbox}[Vězňovo dilema (\emph{Prisoner's dilemma})]
\begin{center}
\begin{tabular}{c|cc}
$i \backslash j$ & $C$ (spolupráce) & $D$ (zrada) \\
\hline
$C$ & $3/3$ & $0/4$ \\
$D$ & $4/0$ & $1/1$ \\
\end{tabular}
\end{center}
Zkratka $C$ znamená \emph{cooperate}, tedy spolupracovat se spolupachatelem, zkratka $D$ pak
\emph{defect}, tedy udat ho a~dostat nižší trest. Výplaty mají kanonické značení: $R$
(\emph{reward} za oboustrannou spolupráci), $P$ (\emph{punishment} za oboustrannou zradu), $T$
(\emph{temptation}, tedy zradím spolupracujícího) a~$S$ (\emph{sucker's payoff}, tedy jsem
zrazen). Hra je vězňovým dilematem, právě když platí
\[
T > R > P > S \qquad\text{a často navíc}\qquad 2R > T + S .
\]
\begin{itemize}
\item $R > P$ zaručuje, že oboustranná spolupráce je lepší než oboustranná zrada.
\item $T > R$ a~$P > S$ znamenají, že \emph{zrada je dominantní strategie pro oba agenty}.
\item $(D,D)$ je \textbf{jediné Nashovo ekvilibrium} hry.
\item $(D,D)$ je zároveň \textbf{jediné ne-Pareto-optimální} řešení hry.
\item $(C,C)$ maximalizuje sociální blahobyt ($3+3=6$).
\end{itemize}
\textbf{Racionalita každého jednotlivce tedy vede ke kolektivně nejhoršímu výsledku.}
\end{defbox}

\textbf{Důsledky vězňova dilematu:}
\begin{itemize}
\item \textbf{Tragédie obecní pastviny} (\emph{tragedy of the commons}): sdílený zdroj, ať už
  pastvina, rybolov, nebo atmosféra, se vyčerpá, protože každý jednotlivec maximalizuje svůj
  užitek.
\item Otevřená otázka, \emph{co vlastně znamená být racionální} a~zda jsou lidé racionální;
  experimentálně totiž lidé spolupracují častěji, než teorie předpovídá.
\item \textbf{Stín budoucnosti} (\emph{shadow of the future}): v~\emph{iterovaném} vězňově
  dilematu se celá situace mění.
\end{itemize}

\begin{defbox}[Iterované vězňovo dilema a~Axelrodovy turnaje]
Hraje-li se PD \emph{konečněkrát a~oba to vědí}, vyjde zpětnou indukcí $D$ v~každém kole:
v~posledním kole není důvod spolupracovat, tedy ani v~předposledním, a~tak dál. Hraje-li se
\emph{nekonečně}, případně s~neznámým koncem daným pravděpodobností pokračování, stává se
spolupráce racionální. Přesně to říká \emph{folk theorem}: každý výplatní profil lepší než
minimax lze podepřít jako ekvilibrium opakované hry.

Robert Axelrod uspořádal v~roce 1980 turnaje strategií pro iterované PD. Vyhrála, a~to dvakrát,
nejjednodušší přihlášená strategie \textbf{Tit For Tat} (TFT, \uv{půjčka za oplátku}):
\emph{v~prvním kole spolupracuj, pak vždy zopakuj poslední tah soupeře}. Z~výsledků odvodil
Axelrod čtyři vlastnosti úspěšných strategií:
\begin{itemize}
\item \textbf{slušnost} (\emph{nice}): nezrazuj jako první;
\item \textbf{odvetnost} (\emph{retaliating}): na zradu okamžitě odpověz;
\item \textbf{odpouštění} (\emph{forgiving}): vrátí-li se soupeř ke spolupráci, odpusť mu;
\item \textbf{nezávistivost} (\emph{non-envious}): neusiluj o~to porazit \emph{konkrétního}
  soupeře. TFT nikdy nezískal víc bodů než jeho protihráč, a~přesto vyhrál turnaj.
\end{itemize}
Jako pátou vlastnost Axelrod uvádí \textbf{srozumitelnost} (\emph{clarity}): strategie musí být
pro protihráče čitelná, jinak s~ní nelze navázat spolupráci. TFT má i~slabinu. Objeví-li se šum,
tedy chybně provedený tah, zaseknou se dva hráči TFT ve vzájemné odvetě; proto vznikly varianty
\emph{Generous TFT}, která občas odpustí, a~\emph{Tit for Two Tats}.
\end{defbox}
\begin{defbox}[Koordinační hra]
\begin{center}
\begin{tabular}{c|cc}
$i \backslash j$ & Balet & Zápas \\
\hline
Balet & $1/2$ & $0/0$ \\
Zápas & $0/0$ & $2/1$ \\
\end{tabular}
\end{center}
Hra je známá také jako \emph{Battle of the Sexes} (BoS) nebo \emph{Bach or Stravinsky}. Výplaty
podporují \emph{shodu} strategií, hráči se ale liší v~tom, na které shodě jim záleží víc.
\begin{itemize}
\item Nashova ekvilibria v~ryzích strategiích jsou dvě: (Balet, Balet) a~(Zápas, Zápas).
\item Sociální blahobyt i~Paretovu optimalitu maximalizují opět tyto dvě dvojice.
\item \textbf{Smíšené NE}: každý hráč hraje svou \emph{preferovanější} možnost
  s~pravděpodobností $2/3$ a~méně preferovanou s~$1/3$. Očekávaná utilita ve smíšeném
  ekvilibriu vychází jen $2/3$, tedy \emph{méně} než v~kterémkoli ryzím ekvilibriu.
\end{itemize}

\emph{Výpočet smíšeného NE (obecný recept pro hry $2\times2$):} hráč $j$ volí pravděpodobnost
$q$ hraní Baletu tak, aby byl hráč $i$ \emph{indiferentní} mezi svými akcemi:
\[
u_i(\text{Balet}) = 1\cdot q + 0\cdot(1-q), \qquad u_i(\text{Zápas}) = 0\cdot q + 2\cdot(1-q).
\]
Z~rovnosti $q = 2 - 2q$ plyne $q = 2/3$, hráč $j$ tedy hraje Balet, svou preferovanou volbu,
s~pravděpodobností $2/3$. Symetricky hráč $i$ hraje Zápas s~pravděpodobností $2/3$. Očekávaná
utilita hráče $i$ je pak $1 \cdot \frac{2}{3} = \frac{2}{3}$.
\textbf{Klíčové pravidlo: ve smíšeném NE volí každý hráč pravděpodobnosti tak, aby byl
\emph{protihráč} indiferentní.}
\end{defbox}

\noindent\textbf{Antikoordinační hra} (\emph{Chicken}, \emph{Hawk--Dove}) je opakem té
předchozí, protože výplaty naopak podporují volbu \emph{různých} strategií:
$u(\text{uhnu},\text{uhnu}) = 5/5$, $u(\text{uhnu},\text{nedám se}) = 1/6$,
$u(\text{nedám se},\text{uhnu}) = 6/1$, $u(\text{nedám se},\text{nedám se}) = 0/0$.
Ryzí NE jsou obě \emph{asymetrické} dvojice, zatímco maximum sociálního blahobytu, tedy situace,
kdy oba uhnou ($5+5$), NE \emph{není}. Symetrické smíšené NE je $p = 1/2$ s~očekávanou utilitou
$3$.

\subsection{Sociální volba a~volby}
\label{ssec:volby}

\noindent\mimoq{ve slidech je celá kapitola, oficiální znění okruhu ji ale nejmenuje ---
zařazeno jako doplněk k~agregaci preferencí}

Teorie her řeší, jak se agenti zachovají při daných preferencích. \emph{Teorie sociální volby}
se ptá opačným směrem: jak z~individuálních preferencí \emph{odvodit} kolektivní rozhodnutí.
Slouží k~tomu dva nástroje. \textbf{Funkce sociálního blahobytu} agreguje preference do úplného
uspořádání $>_{sw}$, kdežto \textbf{funkce sociální volby} vybírá jednoho vítěze.

\textbf{Volebních systémů} existuje celá řada. \emph{Většinový} žádá nad 50~\% hlasů.
\emph{Prostá většina} (\emph{plurality}) přiděluje bod za každé první místo. \emph{Prostá
většina s~vyřazováním} (IRV) opakovaně vyřadí nejslabšího kandidáta a~přerozdělí jeho hlasy.
\emph{Sekvenční většinové volby} staví pavouka párových soubojů, takže \emph{pořadí zápasů
ovlivňuje výsledek}. \emph{Bordova metoda} dává $N-1$ bodů za první místo až po $0$ za poslední,
a~je tedy hlasováním konsenzuálním, nikoli většinovým. \emph{Slaterův systém} hledá acyklické
uspořádání minimalizující počet obrácených hran v~turnajovém grafu, což je NP-těžké.

\textbf{Bordovy kombinace} opravují hlavní slabinu Bordovy metody, totiž to, že nemusí zvolit
Condorcetova vítěze. Všechny tři ho zvolí, pokud existuje. \emph{Blackova metoda} nejdřív zkusí
Condorcetovu metodu, a~není-li vítěz, vezme Bordova. \emph{Baldwinova metoda} spočte Bordova
skóre, vyřadí kandidáta s~nejnižším, skóre \emph{přepočte} jen mezi zbylými a~celý postup
opakuje. \emph{Nansonova metoda} dělá totéž, jen v~každém kole vyřadí všechny kandidáty pod
průměrem, a~konverguje proto rychleji.

Vítěz prosté většiny nemusí být většinový. Při rozdělení $o_1 = 40~\%$, $o_2 = 30~\%$,
$o_3 = 30~\%$ vyhrává $o_1$, ačkoli ho 60~\% voličů nechtělo. Odtud pramení \textbf{taktické
hlasování}.

\begin{defbox}[Condorcetův paradox]
Existují situace, kdy \emph{ať zvolíme kterýkoli výsledek, většina voličů s~ním bude
nespokojena}.

\emph{Příklad:} $V_1: A > B > C$, \quad $V_2: B > C > A$, \quad $V_3: C > A > B$.
V~párových soubojích $A$ porazí $B$ (2:1), $B$ porazí $C$ (2:1) a~$C$ porazí $A$ (2:1).
\textbf{Kolektivní preference je cyklická}, přestože všechny individuální preference jsou
lineární, a~\emph{Condorcetův vítěz} tak neexistuje. Motiv je stejný jako u~neexistence ryzího
Nashova ekvilibria v~kámen--nůžky--papír.
\end{defbox}

\begin{defbox}[Kritéria volebních systémů]
\textbf{Paretova podmínka} (jednomyslnost) žádá, aby platilo $X >_{sw} Y$, kdykoli preferují
$X$ před $Y$ \emph{všichni} voliči. \emph{Splňuje} ji majority a~Borda, \emph{nesplňují} ji
sekvenční většinové volby.

\textbf{Podmínka Condorcetova vítěze} žádá, aby vyhrál ten, kdo porazí všechny v~párovém
srovnání. \emph{Splňují} ji sekvenční většinové volby, \emph{nesplňuje} plurality ani Borda.

\textbf{Nezávislost na irelevantních alternativách} (IIA) žádá, aby to, zda $X >_{sw} Y$,
záviselo pouze na relativním pořadí $X$ a~$Y$ u~voličů. \emph{Nesplňuje} ji prakticky žádný
běžný systém.

Zbývají \textbf{neomezený definiční obor} (univerzalita) a~\textbf{nediktátorství}.
\end{defbox}

\begin{veta}[Arrowova věta o~nemožnosti, 1951]
Pro tři a~více kandidátů neexistuje preferenční volební systém, který by převedl individuální
uspořádání na úplné a~tranzitivní společenské uspořádání a~zároveň splňoval neomezený definiční
obor, nediktátorství, Paretovu podmínku a~IIA.
\end{veta}

Jednodušeji řečeno: pro více než dva kandidáty je jedinou procedurou, která splňuje Paretovu
podmínku a~IIA, \emph{diktatura}. Je to \emph{negativní} výsledek o~limitech kolektivního
rozhodování. Neplyne z~něj, že jediným funkčním systémem je diktatura, nýbrž že každý reálný
systém musí některé kritérium obětovat, nejčastěji IIA.

\textbf{Gibbardova--Satterthwaiteova věta} jde ještě dál: každá surjektivní a~nediktátorská
funkce sociální volby pro tři a~více kandidátů je \textbf{manipulovatelná}. Pro MAS je to
zásadní zjištění. Agent je počítač a~taktické hlasování mu nedělá morální problém, a~proto se
hledají mechanismy odolné vůči manipulaci, což je téma aukcí.

\subsection{Aukce jako mechanismus alokace zdrojů}
\label{ssec:aukce}

Zatímco sociální volba řeší, jak z~individuálních preferencí složit jedno kolektivní rozhodnutí,
aukce odpovídají na užší, zato velmi praktickou otázku: komu připadne vzácný zdroj a~za kolik.

\begin{defbox}[Aukce]
Aukce je \textbf{tržní instituce}, v~níž zprávy od obchodníků nesou cenovou informaci: buď
nabídku ke koupi za danou cenu (\emph{bid}), nebo nabídku k~prodeji za danou cenu (\emph{ask}).
Instituce přitom dává přednost vyšším nabídkám ke koupi a~nižším nabídkám k~prodeji.

V~MAS je aukce především \textbf{mechanismus, jak rozdělit vzácné zdroje mezi agenty}.
\emph{Reálně} takto funguje eBay i~Sotheby's a~takto se přidělují těžební povolení nebo
frekvenční pásma pro mobilní operátory. \emph{V~informatice} se aukcí rozděluje procesorový čas,
šířka pásma, výpočetní úlohy nebo reklamní pozice.

Aukce je \textbf{efektivní}, přidělí-li zdroj agentovi, který ho chce nejvíce, tedy tomu
s~nejvyšším oceněním. Zájmy obou stran jdou přitom proti sobě: prodávající chce cenu
\emph{maximalizovat}, kupující \emph{minimalizovat}, a~každý kupující se řídí vlastní
utilitní funkcí.
\end{defbox}

\begin{defbox}[Klasifikace aukcí]
Aukce se třídí podle dvou nezávislých hledisek.

\textbf{Podle ocenění předmětu} se rozlišují tři situace:
\begin{itemize}
\item \emph{soukromá hodnota} (\emph{private value}) závisí jen na kupujícím samotném, jako
  u~lístku na koncert;
\item \emph{společná hodnota} (\emph{common value}) je pro všechny stejná, jenom ji nikdo předem
  nezná, jako u~těžebních práv; odtud pochází \emph{prokletí vítěze} (\emph{winner's curse}),
  protože vyhrává ten, kdo se ve svém odhadu nejvíce spletl směrem nahoru;
\item \emph{korelovaná hodnota} obě předchozí možnosti kombinuje.
\end{itemize}
\textbf{Podle protokolu} se sleduje, zda vítěz platí první, nebo druhou cenu, zda se draží
otevřeným vyvoláváním (\emph{open cry}), nebo obálkovou metodou (\emph{sealed bid}), a~zda
aukce proběhne v~jediném kole, nebo ve~více kolech.
\end{defbox}
\subsection{Čtyři základní typy aukcí}

Kombinacemi předchozích hledisek se v~praxi ustálily čtyři klasické protokoly. Neliší se jen
tím, jak dražba probíhá, ale hlavně tím, jaká strategie je pro kupujícího racionální.

\begin{center}\small
\begin{tabular}{@{}p{0.15\textwidth}
  >{\raggedright\arraybackslash}p{0.23\textwidth}
  >{\raggedright\arraybackslash}p{0.24\textwidth}
  >{\raggedright\arraybackslash}p{0.27\textwidth}@{}}
\toprule
\textbf{Aukce} & \textbf{Protokol} & \textbf{Racionální strategie} & \textbf{Poznámky} \\
\midrule
\textbf{Anglická}
  & otevřené vyvolávání, \emph{rostoucí} cena, platí se první (nejvyšší) cena
  & \emph{dominantní}: přihazovat po $\varepsilon$, dokud cena nedosáhne mého ocenění,
    pak odejít
  & Nejběžnější (starožitnosti, umění, internet --- asi 95~\% aukcí). Typický případ, kdy
    kupující cenu přeplatí. \\
\textbf{Holandská}
  & otevřené vyvolávání, \emph{klesající} cena, první, kdo přijme, platí aktuální cenu
  & dominantní strategie \emph{neexistuje}; čekat, až cena klesne těsně pod mé ocenění ---
    ale je to hazard
  & Rychlá; zboží podléhající zkáze (květiny, ryby). Oblíbená u~prodávajících.
    Neumožňuje odhadnout tržní cenu ani preference ostatních. \\
\textbf{FPSB} \newline (\emph{first-price sealed bid})
  & jedno kolo, uzavřené obálky, vyhrává nejvyšší nabídka a~platí svou cenu
  & dominantní strategie \emph{neexistuje}; nabídnout \emph{méně} než své ocenění
    (\emph{bid shading}) --- optimum závisí na odhadu ostatních
  & Prodej nemovitostí, státních dluhopisů. Nenutí kupující použít svou utilitní funkci,
    neodhalí tržní cenu. \emph{Strategicky ekvivalentní holandské aukci.} \\
\textbf{Vickreyova} \newline (\emph{second-price sealed bid})
  & jedno kolo, uzavřené obálky, vyhrává nejvyšší nabídka, ale platí \emph{druhou} nejvyšší
  & \emph{dominantní}: nabídnout \textbf{pravdivě} své ocenění
  & Teoreticky nejoblíbenější; Google AdWords, filatelistické aukce.
    \emph{Ekvivalentní anglické aukci} (se soukromými hodnotami). \\
\bottomrule
\end{tabular}
\end{center}

\begin{defbox}[Proč je pravdomluvnost dominantní strategií ve Vickreyově aukci]
Nechť moje ocenění je $v$, nabídnu $b$ a~nejvyšší nabídka ostatních je $p$, kterou ovšem neznám.
Vyhraji, právě když $b > p$, a~pak platím $p$, takže můj zisk činí $v - p$. Zbývá projít obě
možné odchylky od pravdivé nabídky.
\begin{itemize}
\item \textbf{Nadhodnotím, tedy nabídnu $b > v$.} Proti volbě $b=v$ se něco změní jedině tehdy,
  když $v < p < b$. Pak sice vyhraji, ale platím $p > v$, takže můj zisk $v - p < 0$ je záporný.
  \emph{Nadhodnocení mi tedy může jen uškodit.}
\item \textbf{Podhodnotím, tedy nabídnu $b < v$.} Proti volbě $b=v$ se něco změní jedině tehdy,
  když $b < p < v$. Pak prohraji, ačkoli bych byl vyhrál se ziskem $v - p > 0$.
  \emph{O~tento zisk se připravím.}
\item V~žádné situaci tedy \emph{nemohu} nabídkou $b \ne v$ získat proti volbě $b = v$ víc.
\end{itemize}
Pravdivá nabídka $b = v$ je proto (slabě) dominantní strategie, a~to bez ohledu na to, co dělají
ostatní. Vickreyova aukce je tím pádem \textbf{odolná vůči manipulaci} (\emph{strategy-proof},
\emph{incentive compatible}) a~zároveň \emph{efektivní}, protože vyhraje agent s~nejvyšším
oceněním.
\end{defbox}

\medskip
\noindent\textbf{Příklad --- srovnání výnosů.} Do dražby jdou tři kupující s~oceněními
$A = 80$~Kč, $B = 60$~Kč a~$C = 30$~Kč. V~ideálním scénáři, tedy bez dodatečné informace
o~ostatních a~bez podvádění, dopadnou jednotlivé aukce takto:
\begin{center}\small
\begin{tabular}{lll}
\toprule
\textbf{Aukce} & \textbf{Vítěz} & \textbf{Cena} \\
\midrule
Anglická & $A$ & $\approx 61$ (přihazuje se, dokud neodstoupí $B$ na 60) \\
Holandská & $A$ & $80$ nebo o~něco méně (79) --- $A$ nesmí riskovat, že přijme až pozdě \\
FPSB & $A$ & $80$ (resp.\ $60+\varepsilon$, zná-li $A$ ocenění ostatních) \\
Vickreyova & $A$ & $60$ (druhá nejvyšší nabídka) \\
\bottomrule
\end{tabular}
\end{center}
Předmět ve všech čtyřech případech získá agent s~nejvyšším oceněním, všechny aukce jsou tedy
\emph{efektivní}. Rozcházejí se až v~ceně, kterou vítěz zaplatí, a~v~množství informace, kterou
dražba cestou odhalí.

\emph{Proč u~FPSB vychází 80, když se má nabízet pod svým oceněním?} Protože scénář předpokládá
\emph{nulovou} informaci o~ostatních. Agent $A$ neví, jak vysoko půjde $B$, a~tak pro něj každý
ústup od 80 znamená riziko prohry. Míra podhodnocení je funkcí přesvědčení o~ostatních, a~právě
proto FPSB ani holandská aukce \textbf{nemají dominantní strategii}.

Jakmile $A$ ocenění ostatních zná, nabídne $60+\varepsilon$ a~výnos prodávajícího spadne na
úroveň Vickreyovy aukce; srov.\ větu o~ekvivalenci výnosů níže. U~anglické a~Vickreyovy aukce je
situace jiná: dominantní strategie tam existuje a~nezávisí na tom, co dělají ostatní.

\begin{defbox}[Věta o~ekvivalenci výnosů (\emph{revenue equivalence theorem})]
Věta se pojí se jmény Vickrey (1961), Myerson a~Riley--Samuelson. Předpokládá se, že kupující
jsou rizikově neutrální, že jejich soukromé hodnoty jsou nezávislé a~pocházejí ze stejného
spojitého rozdělení, že předmět získá kupující s~nejvyšším oceněním a~že kupující s~nejnižším
možným oceněním má nulový očekávaný zisk. Za těchto předpokladů platí, že
\textbf{všechny čtyři základní aukce dávají prodávajícímu stejný očekávaný výnos.}

Rozdíly, které mezi aukcemi pozorujeme v~praxi, tedy pramení z~porušení některého předpokladu.
Jsou-li kupující averzní k~riziku, vede si lépe FPSB a~holandská aukce. Jde-li o~společné hodnoty
s~hrozbou prokletí vítěze, je lepší anglická aukce, protože během vyvolávání odhaluje informaci.
Hrozí-li naopak koluze mezi kupujícími, je anglická aukce zranitelnější. A~tam, kde rozhodují
náklady na komunikaci, vychází levněji obálková metoda.
\end{defbox}
\subsection{Kombinatorické aukce a~VCG}

Dosud se dražil vždy jediný předmět. Jakmile jich jde do dražby několik najednou, mění se úloha
podstatně, protože kupujícího zajímají spíš celé balíky předmětů než jednotlivé kusy.

\begin{defbox}[Kombinatorická aukce]
Draží se \emph{více předmětů současně} a~kupující podávají nabídky rovnou na \emph{podmnožiny}
(\emph{bundles}). Motivem je, že hodnoty předmětů nejsou aditivní. Někdy se předměty doplňují,
což je \emph{komplementarita}: levá a~pravá bota nebo dva sousední frekvenční bloky mají
dohromady vyšší hodnotu než zvlášť, tedy $v(\{x,y\}) > v(x)+v(y)$. Jindy se naopak zastupují,
což je \emph{substituce}: dvě letenky na tentýž let, tedy $v(\{x,y\}) < v(x)+v(y)$.

Aukcionář pak stojí před \textbf{problémem určení vítězů} (\emph{winner determination problem},
WDP): má vybrat takovou množinu vzájemně disjunktních nabídek, která maximalizuje součet cen.
Je to zobecnění množinového pakování, a~tedy problém \textbf{NP-těžký}, navíc i~těžko
aproximovatelný. Řeší se ILP solvery a~metodou větvení a~mezí (CABOB, CPLEX).

Druhá potíž je \emph{reprezentační}. Podmnožin, na které lze nabízet, je $2^m - 1$, a~proto se
nabídky zapisují v~kompaktních jazycích (OR, XOR, OR-of-XOR).
\end{defbox}
\subsection{Další typy aukcí a~alokace zdrojů}

Čtveřicí základních protokolů se repertoár aukcí nevyčerpává. Praxe si vynutila řadu dalších
variant, které mění buď to, kdo nakonec platí, nebo to, kdo vůbec smí nabídku podat.

\begin{itemize}
\item V~\emph{all-pay auction} zaplatí svou nabídku všichni účastníci, ale předmět získá jen ten
  s~nejvyšší nabídkou. Modeluje se tak lobbing, úplatky, sportovní klání i~patentové závody.
\item \emph{Tichá} aukce je písemnou variantou anglické. \emph{Amsterodamská} začíná jako
  anglická, a~jakmile zbudou dva zájemci, přepne se na holandskou s~dvojnásobnou cenou.
  Na \emph{Tsukiji}, tokijském rybím trhu, nabídnou všichni najednou a~remízy se rozhodují
  hrou kámen--nůžky--papír.
\item V~\emph{dvojité aukci} zadávají nabídky kupující i~prodávající a~protokol je páruje;
  tak funguje burza.
\end{itemize}

\subsection{Návrh mechanismů, VCG a~vyjednávání}

\noindent\mimoq{ve slidech není; okruh jmenuje \uv{alokaci zdrojů}, k~níž toto patří ---
kombinatorické aukce slidy zmiňují jednou větou}

\begin{defbox}[Návrh mechanismů (\emph{mechanism design})]
Teorie her se ptá, \emph{jak se agenti zachovají při daných pravidlech}. Návrh mechanismů otázku
obrací, a~proto se mu říká \uv{obrácená teorie her}: ptá se, \emph{jaká pravidla zvolit, aby
racionální sobečtí agenti sami od sebe dělali to, co chceme}, tedy typicky aby mluvili pravdu.

Mechanismus je dvojice pravidel. První je pravidlo výběru výsledku $x(\hat v)$ z~ohlášených
preferencí, druhé platební pravidlo $p_i(\hat v)$; utilita agenta pak vychází jako
$u_i = v_i(x(\hat v)) - p_i(\hat v)$. Od mechanismu se žádá několik vlastností najednou.
\textbf{Kompatibilita s~motivací} znamená, že se agentovi vyplatí ohlásit pravdu; v~nejsilnější
variantě \emph{strategy-proof} je pravdomluvnost dominantní strategií. \textbf{Individuální
racionalita} požaduje $u_i \ge 0$ a~\textbf{efektivita} maximalizaci sociálního blahobytu;
k~tomu přistupuje \textbf{vyrovnaný rozpočet} a~\textbf{výpočetní zvládnutelnost}.
Vše najednou mít nelze.

\textbf{Princip odhalení} říká, že ke každému mechanismu existuje \emph{přímý} mechanismus,
v~němž agenti ohlásí preference pravdivě a~výsledek zůstane stejný. Proto se celá teorie točí
kolem Vickreyho a~VCG. Bez plateb ovšem naráží na Gibbardovu--Satterthwaiteovu větu.
\end{defbox}

\begin{defbox}[Kombinatorické aukce a~VCG]
V~\textbf{kombinatorické aukci} se draží více předmětů současně a~nabídky se podávají na
\emph{podmnožiny}, protože hodnoty nejsou aditivní. Buď se předměty doplňují
(\emph{komplementarita}, $v(\{x,y\}) > v(x)+v(y)$, například dva sousední frekvenční bloky),
nebo se zastupují (\emph{substituce}, $v(\{x,y\}) < v(x)+v(y)$). \textbf{Problém určení vítězů}
(WDP), tedy volba disjunktních nabídek maximalizujících součet cen, je \textbf{NP-těžký}.

\textbf{VCG} (Vickrey--Clarke--Groves) Vickreyovu aukci zobecňuje: zvolí se alokace $x^*$
maximalizující $\sum_i \hat v_i(x)$ a~agent $i$ zaplatí \textbf{externalitu, kterou způsobuje
ostatním}:
\[
p_i \;=\; \max_{x} \sum_{j \ne i} \hat v_j(x) \;-\; \sum_{j \ne i} \hat v_j(x^*) .
\]
Mechanismus je pravdomluvný, efektivní i~individuálně racionální a~pro jediný předmět se
redukuje na Vickreyovu aukci. Nevýhody jsou tři: WDP se musí řešit $n{+}1$-krát, výnos
prodávajícího bývá malý nebo i~nulový a~celý mechanismus je zranitelný vůči koluzi.

\emph{Příklad:} draží se předměty $\{x,y\}$, přičemž agent 1 nabízí $\{x\}\to5$, agent 2 nabízí
$\{y\}\to4$ a~agent 3 nabízí $\{x,y\}\to8$. Optimální alokace rozdělí předměty mezi první dva
agenty, protože $5+4=9 > 8$. Agent 1 zaplatí $p_1 = 8 - 4 = 4$, neboť bez něj by ostatní měli 8,
kdežto se zvolenou alokací mají 4; jeho zisk tedy činí $5-4=1$. Symetricky vychází
$p_2 = 8-5 = 3$ a~prodávající inkasuje jen $7 < 8$.
\end{defbox}

\begin{defbox}[Vyjednávání: monotónní protokol ústupků a~Zeuthenova strategie]
Každá aukce vyžaduje centrálního aukcionáře. Alternativou, která se bez něj obejde, je přímé
\textbf{vyjednávání} mezi agenty. \emph{Vyjednávací množinu} tvoří ty dohody, které jsou
individuálně racionální a~zároveň Pareto-optimální.

V~\textbf{monotónním protokolu ústupků} navrhnou v~každém kole oba agenti dohodu současně.
Vyjednávání skončí, je-li návrh jednoho z~nich pro druhého alespoň tak dobrý jako jeho vlastní.
Jinak musí alespoň jeden \emph{ustoupit}; když neustoupí nikdo, vyjednávání ztroskotá.

\textbf{Zeuthenova strategie} odpovídá na otázku, \emph{kdo} má ustoupit: ten, kdo má menší
ochotu riskovat konflikt.
\[
\mathrm{Risk}_i = \frac{u_i(\text{vlastní návrh}) - u_i(\text{návrh protějšku})}
{u_i(\text{vlastní návrh})} .
\]
Ustoupí agent s~\emph{nižším} $\mathrm{Risk}$, a~to právě tak málo, aby se poměry obrátily.
Dvojice Zeuthenových strategií je v~Nashově ekvilibriu a~výsledná dohoda maximalizuje součin
utilit, čemuž se říká Nashovo vyjednávací řešení.
\end{defbox}

\subsection{Shrnutí ke~zkoušce}

\begin{itemize}
\item Koherence hodnotí výkon celku, kdežto koordinace znamená minimalizaci synchronizační režie.
\item CDPS stojí na benevolenci agentů a~na sdíleném cíli; odlišit ho od PPS, kde jsou homogenní
  procesory a~paralelní algoritmus, i~od MAS, kde mají agenti vlastní cíle, vznikají konflikty
  a~vyjednává se.
\item CDPS probíhá ve třech fázích: dekompozice, řešení podproblémů, syntéza. Sdílení úloh
  vyžaduje dohodu, sdílení výsledků může být proaktivní i~reaktivní a~přináší čtyři přínosy:
  jistotu, úplnost, přesnost a~včasnost.
\item \textbf{CNET} má čtyři kroky: recognition, announcement, bidding a~awarding/expediting.
  Rolemi jsou manažer a~dodavatel, subkontrakty mohou být hierarchické a~protokol je
  standardizován jako FIPA-Contract-Net.
\item \textbf{BBS} tvoří tabule (hierarchie abstrakce, mutex), experti se svými dovednostmi
  (skills) a~arbitr, který experta vybírá. Vazba mezi experty je volná a~úzkým hrdlem je přístup
  k~tabuli.
\item V~DisCSP a~DCOP jsou proměnné a~omezení rozdělené mezi agenty a~nikdo nevidí celý problém;
  ABT si posílá zprávy \texttt{ok?} a~\texttt{nogood}, DCOP navíc maximalizuje sociální blahobyt.
\item Sobecký agent maximalizuje \emph{svoji} očekávanou utilitu za předpokladu, že ostatní
  dělají totéž; společná utilita se tím nemaximalizuje.
\item Ke hře v~normální formě patří rozdíl mezi ryzími a~smíšenými strategiemi, pojem dominance
  a~dominantní strategie a~iterované odstraňování dominovaných strategií.
\item \textbf{Nashovo ekvilibrium} je profil vzájemně nejlepších odpovědí, z~něhož se nikomu
  nevyplatí jednostranně odchýlit. Naivní hledání projde $m^n$ profilů, \textbf{Nashova věta}
  zaručuje existenci ve smíšených strategiích a~samotné hledání je PPAD-úplné.
\item \textbf{Paretova optimalita} říká, že nelze zlepšit jednomu, aniž by se jinému přitížilo,
  zatímco \textbf{sociální blahobyt} je součet utilit; maximum sw je Pareto-optimální,
  obráceně to neplatí.
\item \textbf{Vězňovo dilema} má výplaty $T>R>P>S$, strategie $D$ je dominantní a~$(D,D)$ je
  jediné NE a~zároveň jediný ne-Pareto-optimální profil, kdežto $(C,C)$ maximalizuje sw. Stejnou
  strukturu má tragedy of the commons. U~iterovaného PD se uplatní stín budoucnosti a~Axelrodova
  TFT, která je slušná, odvetná, odpouštějící, nezávistivá a~srozumitelná.
\item Umět spočítat smíšené NE pro hru $2\times2$: pravděpodobnosti se volí tak, aby byl
  indiferentní \emph{protihráč}.
\item Ze sociální volby znát plurality, IRV, sekvenční hlasování, Bordu a~Slatera; Condorcetův
  paradox dává cyklickou kolektivní preferenci; kritérii jsou Pareto, Condorcetův vítěz, IIA,
  univerzalita a~nediktátorství; \textbf{Arrowova věta} říká, že Pareto + IIA $\Rightarrow$
  diktatura, a~Gibbard--Satterthwaite, že každý systém lze manipulovat.
\item Aukce je mechanismus alokace vzácných zdrojů a~efektivní aukce dá zdroj tomu, kdo ho
  oceňuje nejvíc. Klasifikačními dimenzemi jsou private a~common value, open cry a~sealed bid,
  first a~second price a~počet kol.
\item Čtyři základní typy se liší racionální strategií: anglická (přihazuj po $\varepsilon$
  do svého ocenění) a~Vickreyova (nabídni pravdivě) \emph{dominantní} strategii mají, kdežto
  holandská (čekej na ocenění $-\varepsilon$) a~FPSB (nabídni pod svým oceněním) ji
  \textbf{nemají}. U~Vickreyovy aukce umět \textbf{důkaz pravdomluvnosti}.
\item Platí ekvivalence anglická $\leftrightarrow$ Vickreyova a~holandská $\leftrightarrow$ FPSB.
  Znát \textbf{větu o~ekvivalenci výnosů} i~případy, kdy neplatí (averze k~riziku, common values,
  prokletí vítěze, koluze), a~umět příklad s~oceněními 80/60/30.
\item Návrh mechanismů je obrácená teorie her a~patří k~němu princip odhalení; kombinatorické
  aukce vedou na WDP, který je NP-těžký, a~\textbf{VCG} vybere alokaci maximalizující sw
  a~účtuje platbu rovnou externalitě.
\item Vyjednávání se vede monotónním protokolem ústupků doplněným Zeuthenovou strategií, podle
  níž ustupuje ten s~nižším rizikem.
\end{itemize}

%=====================================================================

\section{Metody pro učení agentů a~zpětnovazební učení}
\label{sec:uceni}
%=====================================================================

\noindent\mimoq{celá kapitola} Přednáška řadí učení agentů výslovně mezi témata, která
\emph{nepokrývá}. Okruh ho ale jmenuje, takže se bez něj u~zkoušky obejít nedá. Kapitola je
proto zpracovaná v~rozsahu, který na zkoušce obstojí: přehled metod učení, k~tomu MDP
a~Q-learning podrobně a~zbytek jen rámcově.

\subsection{Proč a~jak se agenti učí}

Od inteligentního agenta v~MAS se čeká, že bude \emph{adaptivní}, tedy že změní své chování
podle toho, co zažil. Přesně to je učení: změna chování na základě zkušenosti. Cest k~ní vede
celá škála a~liší se tím, odkud agent bere učební signál i~tím, co se vlastně učí.

\begin{defbox}[Metody pro učení agentů]
\begin{itemize}
\item \textbf{Podle učebního signálu} se rozlišují tři klasická paradigmata strojového učení.
  Při \emph{učení s~učitelem} dostává agent ke~vjemům i~správné odpovědi, ať už jde o~učení
  modelu protivníka z~pozorovaných tahů, nebo o~klasifikaci vjemů. \emph{Učení bez učitele}
  hledá strukturu v~samotných pozorováních, typicky shlukuje situace nebo objevuje role.
  \textbf{Zpětnovazební učení} má k~dispozici jen skalární odměnu z~prostředí; pro autonomní
  agenty je nejpřirozenější a~věnuje se mu zbytek kapitoly.
\item \textbf{Imitační učení} (\emph{learning from demonstration}) staví na trajektoriích,
  které předvedl expert. Agent je buď přímo kopíruje (\emph{behavioral cloning}), nebo
  z~expertova chování rekonstruuje odměnovou funkci, což je \emph{inverzní RL}.
\item \textbf{Evoluční metody} šlechtí genetickým algoritmem celou populaci kontrolérů, tedy
  parametrů reaktivních chování, vah neuronové sítě nebo pravidel, a~vybírají podle dosažené
  utility. Je to typický způsob, jak se \uv{experimentálně ladí} reaktivní architektury
  z~kap.~\ref{sec:rizeni}; patří sem i~\emph{learning classifier systems} a~neuroevoluce.
\item \textbf{Podle toho, co se agent učí}: může to být model prostředí, tedy přechodová funkce,
  dále model protivníka (\emph{opponent modelling}), hodnotová funkce, přímo politika,
  nebo u~PRS ohodnocení plánů.
\item \textbf{Individuální vs.\ sociální učení} se liší zdrojem zkušenosti. Buď agent těží jen
  z~té vlastní, nebo i~napodobuje ostatní a~zkušenosti s~nimi sdílí, jak to dělá učení v~roji
  a~sdílené zkušenosti v~MARL.
\end{itemize}
\end{defbox}

Nejběžnějším a~teoreticky nejpropracovanějším přístupem k~učení v~MAS je \textbf{zpětnovazební
učení} (\emph{reinforcement learning}, RL). Agent při něm mění chování metodou pokus--omyl
podle \emph{odměn}, které dostává z~prostředí. Do formalismu z~kapitoly \ref{sec:architektura}
zapadá bez potíží: politika $\pi$ není nic jiného než mechanismus výběru akce a~odměna
operacionalizuje utilitu.

Podstatné je, kolik agentů se učí zároveň, protože podle toho se rozlišují dva zcela odlišné
případy:
\begin{itemize}
\item \textbf{Učení jednoho agenta} je snazší, protože si vystačí s~klasickými algoritmy RL,
  především s~Q-learningem.
\item \textbf{Multiagentní RL} (\emph{MARL}) řeší situaci, kdy se agenti učí \emph{současně}.
  Formalizuje se markovskými hrami, u~nichž se ptáme na Nashovu a~Paretovu optimalitu, a~state
  of the art je dnes hluboké MARL (MADRL).
\end{itemize}

\subsection{Markovský rozhodovací proces}

Aby se o~odměnách a~politikách dalo uvažovat přesně, potřebuje RL formální model prostředí.
Tím modelem je markovský rozhodovací proces.

\begin{defbox}[MDP --- \emph{Markov Decision Process}]
MDP je pětice $\langle S, A, T, R, \gamma\rangle$, jejíž složky znamenají toto:
\begin{itemize}
\item $S$ je prostor stavů a~$A$ prostor akcí;
\item $T : S \times A \times S \to [0,1]$ je přechodová funkce,
  $T(s,a,s') = P(s_{t+1}=s' \mid s_t=s, a_t=a)$;
\item $R : S \times A \times S \to \mathbb{R}$ je funkce odměny, tedy okamžitá odměna
  za přechod;
\item $0 \le \gamma < 1$ je diskontní faktor a~vyjadřuje kompromis mezi okamžitou a~budoucí
  odměnou.
\end{itemize}
Platí \textbf{markovská vlastnost}: následující stav a~odměna závisí pouze na aktuálním stavu
a~akci, nikoli na celé historii. Stojí za srovnání s~přechodovou funkcí $\tau : R^A \to 2^E$
z~abstraktní architektury, která historii \emph{má}; MDP je její markovské zjednodušení
obohacené o~pravděpodobnosti a~odměny.

\textbf{Politika} (\emph{policy}) je zobrazení $\pi : S \to A$, případně stochasticky
$\pi(a|s)$. Cílem je najít \emph{optimální politiku} $\pi^*$, která maximalizuje očekávaný
diskontovaný součet budoucích odměn:
\[
\mathbb{E}\Big[ \sum_{t=0}^{\infty} \gamma^t R(s_t, a_t, s_{t+1}) \Big].
\]
\textbf{Hodnotová funkce} $V^\pi : S \to \mathbb{R}$ přiřazuje stavu očekávanou utilitu za
předpokladu, že agent sleduje politiku $\pi$; \textbf{akční hodnotová funkce} $Q^\pi(s,a)$
dělá totéž pro dvojici stav--akce.
\end{defbox}

\begin{defbox}[Bellmanovy rovnice]
Hodnotu stavu lze zapsat rekurzivně přes hodnoty jeho následníků. Pro danou politiku tak
dostaneme \emph{Bellmanovu rovnici pro evaluaci}:
\[
V^\pi(s) = \sum_{s'} T(s,\pi(s),s') \big[ R(s,\pi(s),s') + \gamma V^\pi(s') \big].
\]
Pro optimální politiku platí \emph{Bellmanova rovnice optimality}:
\begin{align*}
V^*(s) &= \max_{a \in A} \sum_{s'} T(s,a,s') \big[ R(s,a,s') + \gamma V^*(s')\big], \\
Q^*(s,a) &= \sum_{s'} T(s,a,s')\big[R(s,a,s') + \gamma \max_{a'} Q^*(s',a')\big].
\end{align*}
Optimální politika se z~$Q^*$ získá jako $\pi^*(s) = \arg\max_a Q^*(s,a)$.
\end{defbox}

\noindent\textbf{Iterace hodnot} (\emph{value iteration}) se použije tehdy, když známe úplný
popis MDP. Bellmanova rovnice optimality slouží rovnou jako předpis pro přepočet hodnot:
\begin{algorithmic}[1]
\State inicializuj $V_0(s) \gets 0$ pro všechna $s$
\Repeat
  \ForAll{$s \in S$}
    \State $V_{k+1}(s) \gets \max_{a} \sum_{s'} T(s,a,s')\big[R(s,a,s') + \gamma V_k(s')\big]$
  \EndFor
\Until{$\max_s |V_{k+1}(s) - V_k(s)| < \varepsilon$}
\State \Return $\pi(s) = \arg\max_a \sum_{s'} T(s,a,s')[R(s,a,s') + \gamma V(s')]$
\end{algorithmic}
Že tento postup skutečně někam dojde, plyne z~toho, že Bellmanův operátor je kontrakce
s~koeficientem $\gamma$; iterace proto konverguje ke $V^*$ geometricky rychle. Alternativou
je \textbf{iterace politiky} (\emph{policy iteration}), která místo hodnot vylepšuje rovnou
politiku a~střídá přitom dva kroky: nejprve aktuální politiku \emph{vyhodnotí} vyřešením
soustavy lineárních rovnic, pak ji hladově podle $V^\pi$ \emph{zlepší}. Konverguje v~konečně
mnoha krocích, protože politik je konečně mnoho.

\subsection{Q-learning a~SARSA}

Předchozí algoritmy předpokládaly, že máme MDP celé před sebou. Ve skutečnosti ovšem $T$ ani
$R$ obvykle \textbf{neznáme}. Agent se proto musí učit ze zkušenosti, kterou získává interakcí
s~prostředím v~diskrétních krocích, a~tomuto přístupu se říká \emph{model-free} RL.

\begin{defbox}[Q-learning (Watkins, 1989)]
Agent si udržuje \emph{tabulku} $Q(s,a)$, tedy odhad diskontovaného součtu budoucích odměn při
provedení akce $a$ ve stavu $s$. Po každém přechodu $s \xrightarrow{a} s'$ s~odměnou $r$
přepíše příslušnou položku:
\[
Q(s,a) \;\leftarrow\; Q(s,a) + \alpha\Big[\, \underbrace{r + \gamma \max_{a'} Q(s',a')}
_{\text{cíl (TD target)}} - Q(s,a) \Big],
\qquad 0 \le \alpha \le 1 \text{ je rychlost učení,}
\]
což se dá zapsat i~jako
$Q(s,a) \leftarrow (1-\alpha)Q(s,a) + \alpha\big(r + \gamma\max_{a'}Q(s',a')\big)$.

\textbf{Off-policy:} cíl obsahuje $\max_{a'}$, takže se algoritmus učí hodnotu \emph{optimální}
politiky bez ohledu na to, jakou politikou agent skutečně jedná. Za předpokladu, že každá
dvojice $(s,a)$ je navštívena nekonečněkrát a~že $\alpha$ vhodně klesá, konverguje $Q \to Q^*$.
\end{defbox}

\begin{defbox}[Dilema průzkumu a~využití; $\varepsilon$-greedy]
Agent stojí před dvěma protichůdnými požadavky. Má \emph{využívat} (\emph{exploit}) to, co už
umí, protože jinak zbytečně přichází o~odměnu, ale zároveň musí \emph{zkoušet} (\emph{explore})
nové akce, protože jinak se o~ničem lepším nedozví. Standardním řešením tohoto dilematu je
$\varepsilon$-greedy politika:
\[
\pi(s) = \begin{cases}
\text{náhodná akce z } A & \text{s pravděpodobností } \varepsilon,\\
\arg\max_a Q(s,a) & \text{s pravděpodobností } 1-\varepsilon .
\end{cases}
\]
Hodnota $\varepsilon$ se typicky v~čase snižuje (\emph{annealing}). Alternativou je
softmax neboli Boltzmannovo rozdělení $\pi(a|s) \propto e^{Q(s,a)/\tau}$, dále optimistická
inicializace $Q$ a~UCB, tedy horní meze spolehlivosti.
\end{defbox}

\medskip
\noindent\textbf{Konkrétní krok Q-learningu.} Nechť $\alpha = 0{,}5$ a~$\gamma = 0{,}9$.
Agent je ve stavu $s_1$, provede akci $a_1$, dostane odměnu $r = 10$ a~přejde do $s_2$.
Aktuální hodnoty jsou $Q(s_1,a_1) = 4$, $Q(s_2,a_1) = 6$ a~$Q(s_2,a_2) = 8$. Výpočet pak
proběhne takto:
\begin{align*}
\max_{a'} Q(s_2,a') &= \max(6,\,8) = 8, \\
\text{TD cíl} &= r + \gamma \max_{a'} Q(s_2,a') = 10 + 0{,}9 \cdot 8 = 17{,}2, \\
\text{TD chyba} &= 17{,}2 - Q(s_1,a_1) = 17{,}2 - 4 = 13{,}2, \\
Q(s_1,a_1) &\leftarrow 4 + 0{,}5 \cdot 13{,}2 = \mathbf{10{,}6}.
\end{align*}

\begin{defbox}[SARSA]
Název je zkratka za \emph{State--Action--Reward--State--Action} a~vyjmenovává přesně to, co
aktualizace potřebuje. Na rozdíl od Q-learningu totiž použije \emph{skutečně zvolenou}
následující akci $a'$, tedy akci vybranou touž politikou, kterou agent jedná:
\[
Q(s,a) \leftarrow Q(s,a) + \alpha\big[r + \gamma\, Q(s',a') - Q(s,a)\big].
\]
\textbf{On-policy:} učí se hodnotu politiky, kterou agent skutečně provádí, včetně jejího
náhodného průzkumu. Důsledkem je, že SARSA je \emph{opatrnější}. V~úloze \uv{chůze po útesu}
(\emph{cliff walking}) se Q-learning naučí optimální cestu těsně u~propasti a~při
$\varepsilon$-greedy do ní občas spadne, zatímco SARSA se naučí bezpečnější cestu dál od
okraje, protože započítá riziko vlastního průzkumu.

Obě metody zobecňuje $\mathrm{TD}(\lambda)$ s~\emph{eligibility traces}, které rozprostírají
kredit na více předchozích kroků.
\end{defbox}

\subsection{Metody založené na politice}
\label{ssec:policy-metody}

Q-learning i~SARSA hledají politiku oklikou přes hodnotovou funkci. Druhá velká rodina metod
tuto okliku vynechává.

\begin{defbox}[Policy gradient a~actor--critic]
Namísto hodnotových funkcí lze optimalizovat \emph{přímo politiku}. Politiku parametrizujeme
vektorem $z$, tedy $\pi(a \mid s, z)$, a~hledáme $z^*$ gradientním výstupem po očekávaném
výnosu. \textbf{REINFORCE} (Williams, 1992) odhaduje gradient z~Monte Carlo simulací celých
epizod:
\[
z_{t+1} = z_t + \alpha\, G_t\, \nabla_z \ln \pi(A_t \mid S_t, z_t),
\]
kde $G_t$ je výnos od času $t$. Intuitivně předpis říká: zvyš pravděpodobnost akcí, po nichž
následoval vysoký výnos. Proti hodnotovým metodám má dvě výhody, totiž že zvládá \emph{spojité}
akční prostory a~přirozeně vytváří \emph{stochastické} politiky, které jsou nutné v~hrách se
smíšenými ekvilibrii. Nevýhodou je vysoký rozptyl odhadu gradientu.

\textbf{Actor--critic} oba přístupy kombinuje: herec (\emph{actor}) reprezentuje politiku
a~kritik (\emph{critic}) se učí hodnotovou funkci a~snižuje rozptyl. Přidá-li se
\textbf{funkce výhody} $A(s,a) = Q(s,a) - V(s)$, škáluje se gradient relativní výhodností akce
místo celého výnosu.
Do rodiny patří varianty A3C, A2C, PPO, SAC a~DDPG.

\textbf{Hluboké RL:} \textbf{DQN} nahrazuje tabulku $Q$ neuronovou sítí. Klíčové jsou přitom
dva triky: \emph{experience replay} rozbíjí korelace mezi vzorky a~\emph{target network}
stabilizuje cíl.
\end{defbox}

\paragraph{Proč hluboké RL potřebuje ty triky: smrtící trojice.} Konvergenční záruky
Q-learningu platí pro \emph{tabulkovou} reprezentaci. Jakmile se sejdou tři věci najednou,
může učení \emph{divergovat}: \textbf{aproximace hodnotové funkce}, kdy tabulku nahradí
neuronová síť, \textbf{bootstrapping}, kdy cíl aktualizace obsahuje vlastní odhad
$\max_{a'} Q(s',a')$, a~\textbf{off-policy učení}, kdy data pocházejí z~jiné politiky, než
kterou se učíme. Sutton a~Barto tuto trojici nazývají \emph{deadly triad}. Každá dvojice sama
o~sobě je zvládnutelná; DQN se s~trojicí vyrovnává právě replay bufferem a~zamrzlou cílovou
sítí, které oslabují druhý a~třetí člen.

\begin{defbox}[Tvarování odměny (\emph{reward shaping})]
Je-li odměna řídká a~agent dostane signál až v~cíli, stává se učení prakticky neproveditelným.
Odměna se proto doplňuje o~průběžné nápovědy $F(s,a,s')$. Naivní doplnění ale mění optimální
politiku: agent se naučí sbírat nápovědy místo řešit úlohu.

\textbf{Ng, Harada, Russell (1999)} ukázali, že doplnění tvaru
\[
F(s,a,s') \;=\; \gamma\,\Phi(s') - \Phi(s)
\]
zachovává množinu optimálních politik pro libovolnou \emph{potenciálovou funkci}
$\Phi : S \to \mathbb{R}$, a~že je to jediný tvar, který to zaručuje pro všechny MDP.
Intuice za tím je prostá: součet takových přírůstků po libovolné cestě teleskopicky závisí
jen na koncových bodech, nikoli na tom, kudy agent šel, takže se cykly nevyplatí.
\end{defbox}

\subsection{Multiagentní zpětnovazební učení}

Dosud se učil jediný agent v~prostředí, které se pod ním nemění. Jakmile se učí víc agentů
naráz, přestává to platit a~formalismus se musí rozšířit.

\begin{defbox}[Markovská hra (\emph{Markov game}, stochastická hra)]
Markovská hra zobecňuje MDP na $N$ agentů a~má tvar
$\langle N, S, (A_i)_{i\in N}, T, (R_i)_{i \in N}, \gamma\rangle$.
\textbf{Sdružený akční prostor} je součin individuálních,
$\mathbf{A} = A_1 \times \cdots \times A_N$. Přechod závisí na akcích \emph{všech} agentů
a~každý agent má \emph{vlastní} odměnu $R_i$ i~vlastní politiku $\pi_i$, ale jeho hodnotová
funkce závisí na politikách všech ostatních.

\textbf{Spojení s~teorií her} je bezprostřední: politika agenta je \emph{nejlepší odpovědí} na
sdružené politiky ostatních, sdružené politiky mohou být v~\textbf{Nashově ekvilibriu} a~mohou
být \textbf{Pareto-optimální}. Zajímavé jsou speciální případy. Pro $N=1$ dostaneme MDP, hra
s~jediným stavem dává hru v~normální formě a~$N=2$ s~nulovým součtem dává \emph{minimax} hru.

\textbf{Minimax-Q} (Littman, 1994) řeší hry dvou agentů s~nulovým součtem. Místo
$\max_{a'} Q(s',a')$ se použije hodnota smíšeného minimaxového ekvilibria matice
$Q(s',\cdot,\cdot)$, což znamená v~každém kroku vyřešit lineární program. \textbf{Nash-Q} dělá
totéž pro obecný součet, konverguje ale jen za silných předpokladů.
\end{defbox}

\begin{defbox}[Proč je MARL těžké]
Potíže jsou čtyři a~objevují se prakticky vždy:
\begin{itemize}
\item Agenti aktualizují politiky \textbf{současně}, takže z~pohledu jednoho z~nich je
  prostředí \textbf{nestacionární}. Tím se porušuje základní předpoklad konvergence RL a~agent
  míří na \uv{pohyblivý cíl}.
\item Sdružený akční prostor roste \textbf{exponenciálně} s~počtem agentů.
\item Prostředí je jen \textbf{částečně pozorovatelné}, což vede na Dec-POMDP, jehož optimální
  řešení je NEXP-úplné.
\item Objevuje se problém \textbf{přiřazení zásluh}: při společné odměně není jasné, který
  agent k~ní přispěl.
\end{itemize}
\textbf{Standardní protiopatření} je \emph{centralizované učení s~decentralizovaným prováděním}
(CTDE): kritik během tréninku vidí stav i~akce všech agentů, zatímco herec při běhu vystačí
jen se svými lokálními pozorováními (MADDPG, COMA, QMIX). Dále se používá sdílení parametrů
mezi homogenními agenty a~\emph{opponent modelling}.
\end{defbox}

\textbf{Self-play.} Zvláštní postavení má situace, kdy se agent učí hraním proti sobě samému,
případně proti svým dřívějším verzím. V~hrách s~nulovým součtem a~úplnou informací vede takový
postup spolehlivě k~silným strategiím, jak dokládají TD-Gammon, AlphaGo Zero a~AlphaZero.
Rizikem jsou \emph{cyklické} preference strategií (kámen--nůžky--papír), u~nichž naivní
self-play osciluje, místo aby se zlepšoval; proto se zavádějí ligy protivníků a~populační
metody.

\subsection{Shrnutí ke~zkoušce}

\begin{itemize}
\item Metody učení agentů se podle učebního signálu dělí na učení s~učitelem, učení bez učitele
  a~učení zpětnovazební. Vedle tohoto dělení stojí imitační učení spolu s~inverzním RL
  a~evoluční metody (LCS, neuroevoluce); dále se metody třídí podle toho, co se agent učí, tedy
  zda model prostředí, model protivníka, nebo přímo politiku, a~podle toho, jestli se učí
  individuálně, nebo sociálně. Nejběžnějším přístupem je RL.
\item \textbf{MDP} je pětice $\langle S,A,T,R,\gamma\rangle$ s~markovskou vlastností; nad ní se
  definuje politika $\pi$ a~hodnotové funkce $V^\pi$ a~$Q^\pi$. U~zkoušky je potřeba napsat
  Bellmanovu rovnici optimality a~vysvětlit value iteration, která konverguje, protože
  Bellmanův operátor je kontrakce, i~policy iteration.
\item \textbf{Q-learning} aktualizuje
  $Q(s,a) \leftarrow Q(s,a) + \alpha[r + \gamma\max_{a'}Q(s',a') - Q(s,a)]$, je off-policy
  a~tabulkový; jeden krok je potřeba umět spočítat s~čísly.
\item \textbf{SARSA} používá aktualizaci $\dots + \alpha[r + \gamma Q(s',a') - Q(s,a)]$, je
  on-policy a~chová se opatrněji, což ukazuje úloha cliff walking. Dilema průzkumu a~využití
  řeší $\varepsilon$-greedy politika.
\item \textbf{Policy gradient} optimalizuje politiku přímo (REINFORCE, tedy $\nabla \ln \pi$
  krát výnos) a~\textbf{actor-critic} k~němu přidává kritika: herec je politika, kritik
  hodnotová funkce a~výhoda je $A = Q - V$. A3C je asynchronní varianta, DQN přináší replay
  buffer a~target network.
\item \textbf{Smrtící trojice} aproximace, bootstrappingu a~off-policy učení může vést
  k~divergenci, a~právě proto potřebuje DQN své triky. U~\textbf{tvarování odměny} platí, že
  optimální politiku zaručeně nemění jen potenciálový tvar $F = \gamma\Phi(s')-\Phi(s)$.
\item \textbf{Markovská hra} je MDP pro $N$ agentů se sdruženým akčním prostorem
  a~individuálními odměnami; ptáme se u~ní na nejlepší odpověď, Nashovo ekvilibrium
  a~Pareto-optimalitu. \textbf{Minimax-Q} řeší hry s~nulovým součtem a~v~každé aktualizaci
  počítá LP, \textbf{Nash-Q} míří na obecný součet.
\item MARL komplikují čtyři věci: nestacionarita prostředí, exponenciální sdružený akční
  prostor, částečná pozorovatelnost a~přiřazení zásluh. Standardní odpovědí na ně je CTDE.
  Za zmínku stojí ještě self-play a~AlphaZero.
\end{itemize}

%=====================================================================

\section{Metodologie návrhu, jazyky a~prostředí MAS}
\label{sec:metodologie}
%=====================================================================

\subsection{Agentově orientované softwarové inženýrství}
Objektově orientované metodologie a~jejich notace UML na multiagentní systém nestačí, a~to
hned ze tří důvodů. Agent je autonomní, takže mu nelze přímo zavolat metodu a~očekávat, že ji
provede. Agent má vlastní cíle, které se nemusejí krýt s~cíli toho, kdo ho volá. A~vztahy mezi
agenty nejsou strukturální, nýbrž \emph{organizační}: popisují se rolemi, protokoly a~závazky,
ne dědičností a~agregací. Právě proto vznikly samostatné metodologie \textbf{AOSE}
(\emph{agent-oriented software engineering}). Popisují systém v~pojmech rolí a~interakcí mezi
nimi, a~ne v~pojmech tříd a~volání metod.

\begin{defbox}[Gaia (Wooldridge, Jennings, Kinny, 2000)]
Jedna z~nejcitovanějších AOSE metodologií. Její klíčovou metaforou je \textbf{organizace}:
multiagentní systém se popisuje tak, jako by šlo o~organizaci, v~níž agenti hrají \textbf{role}.
Gaia postupuje od požadavků přes analýzu k~návrhu, samotnou implementaci nechává na konkrétní
platformě. Analýza a~návrh přitom tvoří dvě oddělené vrstvy popisu: nejprve se systém zachytí
abstraktně, v~rolích a~jejich interakcích, a~teprve pak se tento popis převede na konkrétní
typy agentů a~jimi poskytované služby.

\textbf{Fáze analýzy} staví dva abstraktní modely. \emph{Model rolí} (\emph{roles model})
popisuje každou roli čtyřmi atributy. \emph{Permissions} určují, jaké zdroje smí role číst
a~měnit. \emph{Responsibilities} říkají, co má role zajistit, a~dělí se na dvě skupiny:
\emph{liveness} vlastnosti znějí \uv{něco dobrého se stane} a~popisují se regulárním výrazem
nad aktivitami a~protokoly, kdežto \emph{safety} vlastnosti znějí \uv{nic špatného se nestane}
a~popisují se invarianty. \emph{Activities} jsou výpočty, které role provádí sama,
zatímco \emph{protocols} zachycují její interakce s~jinými rolemi. Druhý model, \emph{model
interakcí}, tyto protokoly mezi rolemi definuje: jejich účel, iniciátora, respondéra, vstupy
a~výstupy.

\textbf{Fáze návrhu} z~obou abstraktních modelů odvozuje tři konkrétní. \emph{Model agentů}
přiřazuje role typům agentů a~stanoví počty instancí. \emph{Model služeb} popisuje funkce,
které role poskytuje, tedy jejich vstupy, výstupy, předpoklady a~efekty. \emph{Model známostí}
(\emph{acquaintance model}) je orientovaný graf zachycující, kdo s~kým komunikuje; slouží ke
kontrole úzkých hrdel a~přílišné provázanosti.

\textbf{Omezení:} původní Gaia počítá s~uzavřeným systémem, kooperativními agenty a~statickou
organizací. Rozšíření \emph{ROADMAP} a~\emph{Gaia v.2} tato omezení uvolňují.
\end{defbox}

\noindent\mimoq{slidy jmenují jen Gaiu a~role; ostatní metodologie jsou orientační}
Gaia zdaleka není jediná. \textbf{Prometheus} (Padgham, Winikoff, 2002) je praktická
metodologie mířená na \textbf{BDI agenty} a~má tři fáze. Ve \emph{specifikaci systému} se
stanoví cíle, \emph{funkcionality} a~rozhraní s~okolím, tedy \emph{percepts} a~\emph{actions}.
\emph{Architektonický návrh} pak určí typy agentů, protokoly a~zprávy. \emph{Detailní návrh}
nakonec sestoupí ke capabilities, plánům a~událostem, což jsou přímo stavební kameny BDI.

\textbf{Tropos} staví na modelování \emph{aktérů a~cílů} pomocí frameworku $i^*$ a~jako jediný
pokrývá i~fázi raných požadavků. Pracuje s~pojmy \emph{actor}, \emph{goal} a~\emph{softgoal},
což je nefunkční požadavek bez ostrého kritéria, a~dále s~pojmy \emph{plan}, \emph{resource}
a~\emph{dependency}. Vedle těchto tří se používají ještě \textbf{MaSE}, \textbf{INGENIAS},
\textbf{ADELFE} a~\textbf{PASSI}.

\subsection{Standard FIPA a~architektura platformy}

Metodologie říká, jak systém navrhnout. Standard FIPA odpovídá na jinou otázku: jak má vypadat
prostředí, ve kterém navržení agenti poběží a~budou spolu mluvit. Popisuje proto referenční
architekturu platformy, tedy komponenty, ze kterých se agentní platforma skládá, a~úlohy,
jež každá z~nich plní.

\begin{defbox}[Referenční architektura agentní platformy podle FIPA]
\begin{itemize}
\item \textbf{AMS} (\emph{Agent Management System}) je povinná komponenta a~funguje jako
  \uv{matrika} platformy: eviduje agenty, přiděluje jim jednoznačné \emph{AID}
  (agent identifier), řídí jejich životní cyklus (create, suspend, resume, terminate)
  a~poskytuje bílé stránky (\emph{white pages}).
\item \textbf{DF} (\emph{Directory Facilitator}) hraje roli žlutých stránek
  (\emph{yellow pages}). Agenti u~něj registrují \emph{služby}, které nabízejí, a~naopak
  v~něm vyhledávají poskytovatele služeb.
\item \textbf{MTS} (\emph{Message Transport Service}) doručuje ACL zprávy jak uvnitř jedné
  platformy, tak mezi platformami; využívá k~tomu protokoly IIOP a~HTTP a~kódování
  bitefficient, XML nebo string.
\item \textbf{ACC} (\emph{Agent Communication Channel}) tvoří rozhraní k~MTS.
\end{itemize}
Standard nekončí u~architektury: definuje dále jazyk ACL, obsahové jazyky (SL), interakční
protokoly a~správu ontologií.
\end{defbox}

\subsection{Jazyky a~prostředí}
\label{ssec:jazyky}

\noindent\mimoq{AGENT-0 a~Concurrent MetateM ve slidech nejsou; JADE a~Jason ano}
Historickým východiskem je \textbf{agentově orientované programování} (Shoham 1993).
V~tomto paradigmatu tvoří stav programu \emph{mentální kategorie}, tedy přesvědčení, závazky
a~schopnosti, a~program sám je množinou \emph{pravidel závazků}. Prvním jazykem tohoto druhu
byl \textbf{AGENT-0}. Jinou cestou jde \textbf{Concurrent MetateM}, kde se specifikace
a~program nerozlišují: agent je \emph{přímo spustitelnou specifikací} v~temporální logice.
Oba jazyky navazují na linii deduktivních agentů (odst.\ 2.1). Prakticky nejúspěšnějším
agentním jazykem je dnes AgentSpeak/Jason a~právě jemu spolu s~dalšími třemi systémy se
věnuje následující přehled.
\begin{center}\small
\begin{tabular}{p{0.15\textwidth} p{0.20\textwidth} p{0.57\textwidth}}
\toprule
\textbf{Systém} & \textbf{Typ} & \textbf{Charakteristika} \\
\midrule
\textbf{JADE}
  & Java, FIPA platforma
  & \emph{Java Agent DEvelopment Framework}, jedna z~nejrozšířenějších platforem pro MAS, plně
    v~souladu s~FIPA. Agent = třída dědící z~\texttt{Agent} s~metodou \texttt{setup()};
    chování se skládá z~objektů \texttt{Behaviour} (\texttt{OneShot}, \texttt{Cyclic},
    \texttt{Ticker}, \texttt{FSMBehaviour}, \texttt{SequentialBehaviour}), které plánuje
    kooperativní (nepreemptivní) plánovač uvnitř jednoho vlákna agenta. Zprávy:
    \texttt{ACLMessage}, \texttt{send()}/\texttt{blockingReceive()}, \texttt{MessageTemplate}.
    Hotové třídy pro protokoly (\texttt{ContractNetInitiator}, \texttt{AchieveREResponder}),
    ladicí nástroje (RMA, Sniffer, Introspector), kontejnery a~mobilita agentů. \\
\textbf{Jason}
  & interpret AgentSpeak(L) v~Javě
  & Referenční implementace \textbf{AgentSpeak}, tedy \textbf{BDI} jazyka (logické programování
    + BDI), přímý potomek PRS. Program agenta = počáteční \emph{beliefs}, \emph{rules}
    a~\emph{knihovna plánů} ve tvaru \texttt{spouštěč : kontext <- tělo}. Spouštěčem je
    \emph{událost} (\texttt{+b} přidání přesvědčení, \texttt{+!g} přidání cíle, \texttt{-!g}
    selhání cíle). Cyklus interpretu: vnímej $\to$ aktualizuj beliefs $\to$ vyber událost $\to$
    najdi \emph{relevantní} plány (podle spouštěče) $\to$ vyber \emph{aplikovatelné}
    (podle kontextu) $\to$ vytvoř záměr $\to$ vykonej krok. S~\texttt{CArtAgO} (prostředí)
    a~\texttt{Moise} (organizace) tvoří \emph{JaCaMo}. \\
\textbf{NetLogo}
  & prostředí pro ABM
  & Jazyk odvozený od Logo pro \emph{agent-based modeling}. Entity: \emph{turtles} (mobilní
    agenti), \emph{patches} (buňky prostředí), \emph{links} a~\emph{observer}. Knihovna modelů
    (Segregation, Wolf Sheep Predation, Ants, Flocking), \emph{BehaviorSpace} pro parametrické
    studie. Určeno pro emergentní chování a~výuku, ne pro FIPA komunikaci. \\
\textbf{SPADE}
  & Python, XMPP
  & \emph{Smart Python Agent Development Environment}. Komunikace přes standardní
    \textbf{XMPP} místo vlastního MTS, agent = třída s~\emph{behaviours}
    (\texttt{CyclicBehaviour}, \texttt{PeriodicBehaviour}, \texttt{FSMBehaviour}), asyncio,
    webové rozhraní; populární díky napojení na Python ekosystém strojového učení. \\
\bottomrule
\end{tabular}
\end{center}

\textbf{Další prostředí.} Vedle čtyř systémů z~tabulky stojí za zmínku tři skupiny nástrojů.
\emph{Jadex}, \emph{JACK} a~\emph{JAM} jsou další implementace BDI/PRS, přičemž Jadex staví
BDI vrstvu nad JADE. \emph{MASON}, \emph{Repast} a~\emph{GAMA} slouží jako simulační platformy
pro agentní modely. \emph{PettingZoo}, \emph{Gymnasium} a~\emph{OpenSpiel} pak nabízejí
standardní prostředí pro multiagentní zpětnovazební učení.

\subsection{Agenti postavení na velkých jazykových modelech}
\label{ssec:llm-agenti}

\noindent\mimoq{nově ve slidech verze 2025 (kapitola \emph{LLM Agents})}

Poslední přírůstek do rodiny mechanismů výběru akcí (\S\ref{ssec:vyber-akci}) nepoužívá ani
dedukci, ani pravidla, ani naučenou $Q$-funkci. Rozhoduje za něj \textbf{velký jazykový model}
(LLM). Agent se pak scvrkne na smyčku, v~níž se stav prostředí zakóduje do textového promptu
a~výstup modelu se interpretuje jako akce. Rozhodovací logika tedy nesídlí v~kódu agenta,
nýbrž v~modelu; kód kolem něj jen sestavuje prompt a~překládá odpověď zpět na akci.

\begin{defbox}[LLM v~roli komponenty agenta]
\begin{itemize}
\item \textbf{LLM jako politika} (\emph{action selection}). Prompt obsahuje popis role, stav
  prostředí a~seznam dostupných akcí a~výstupem je zvolená akce. Model tak plní funkci
  $\mathit{action} : I \to A$ z~abstraktní architektury, přičemž vnitřním stavem $I$ je
  kontextové okno, případně externí paměť.
\item \textbf{LLM jako plánovač.} Model rozloží cíl na podcíle a~generuje kandidátní plány.
  Prohledávání prostoru plánů řídí buď model sám (\emph{chain of thought}, \emph{ReAct}), nebo
  vnější prohledávací algoritmus, jemuž model slouží jako generátor nástupníků a~zároveň jako
  heuristika.
\item \textbf{LLM jako kritik.} Ve schématu \emph{actor--critic} (\S\ref{ssec:policy-metody})
  model hodnotí kvalitu provedené akce nebo celé trajektorie, a~poskytuje tak zpětnou vazbu
  tam, kde prostředí skalární odměnu nedává.
\item \textbf{LLM jako rozhraní k~nástrojům} (\emph{tool use}). Model volí, který externí
  nástroj zavolat a~s~jakými argumenty; tím nástrojem bývá vyhledávání, kalkulátor, API nebo
  jiný agent. Je to přímá obdoba \emph{služeb} registrovaných v~DF a~jejich volání přes ACL.
\item \textbf{Učení s~člověkem ve smyčce} (\emph{human-in-the-loop}). Člověk opravuje
  a~doplňuje zpětnou vazbu a~podle ní se upravuje prompt, paměť nebo přímo váhy modelu.
\end{itemize}
\end{defbox}

\textbf{Typické nasazení: agentové simulace.} Nejčastější aplikací nejsou optimalizační úlohy,
nýbrž \emph{agent-based modeling} (\S\ref{ssec:jazyky}, NetLogo). LLM agent tam nahradí
jednoduchého reaktivního agenta a~dodá mu doménovou znalost i~jazykové chování:
\begin{itemize}
\item \emph{Šíření informací v~sociální síti.} Agenti jsou uzly grafu sociálních interakcí
  a~zprávu si předávají v~přirozeném jazyce; sleduje se, jak rychle a~jakým tvarem se zpráva
  sítí šíří.
\item \emph{Syntetický volební panel.} Populace agentů se vygeneruje podle sociologických
  průzkumů, tedy podle demografie a~postojů, a~použije se k~odhadu výsledků voleb nebo reakce
  na kampaň.
\end{itemize}

\textbf{Zařazení a~kritika.} Z~hlediska okruhu jde o~další rodinu mechanismů výběru akcí,
subsymbolickou stejně jako neuronové sítě, ale s~výrazně širší doménovou znalostí a~bez
nutnosti trénovat model úlohy. Slabiny jsou zrcadlovým obrazem těchto předností. Agent nemá
explicitní model světa, takže o~něm nelze dokazovat tvrzení jako u~deduktivního agenta.
Sémantika jeho \uv{přesvědčení} je \textbf{neverifikovatelná} ve stejném smyslu jako
u~FIPA-ACL (\S\ref{ssec:kqml-acl}). Chování je nedeterministické a~cena jednoho rozhodnutí je
o~několik řádů vyšší než u~pravidlového agenta.

\subsection{Shrnutí ke~zkoušce}

\begin{itemize}
\item Objektově orientované metodologie nestačí proto, že agent je autonomní a~nelze mu jen
  zavolat metodu, že má vlastní cíle a~že vztahy mezi agenty jsou organizační, nikoli
  strukturální.
\item \textbf{Gaia} popisuje systém jako organizaci složenou z~rolí. Analýza dává model rolí
  (permissions, responsibilities s~liveness a~safety vlastnostmi, activities, protocols)
  a~model interakcí, návrh pak model agentů, model služeb a~model známostí. Metodologie
  předpokládá kooperativní agenty a~statickou organizaci.
\item \textbf{Prometheus} cílí na BDI a~postupuje od funkcionalit $\to$ přes typy agentů
  $\to$ ke capabilities a~plánům, kdežto \textbf{Tropos} pracuje s~aktéry, cíli a~softgoals
  a~se závislostmi mezi nimi a~jako jediný pokrývá i~rané požadavky.
\item \textbf{FIPA platforma} stojí na třech komponentách. AMS spravuje bílé stránky, životní
  cyklus agentů a~přidělování AID, DF plní úlohu žlutých stránek s~nabízenými službami
  a~MTS spolu s~rozhraním ACC doručuje ACL zprávy uvnitř platformy i~mezi platformami.
\item \textbf{JADE} je javová FIPA platforma stavěná z~tříd \texttt{Agent},
  \texttt{Behaviour} a~\texttt{ACLMessage}, s~hotovými protokoly a~ladicími nástroji.
  \textbf{Jason/AgentSpeak} je BDI jazyk s~plány tvaru \texttt{spouštěč : kontext <- tělo}
  a~potomek PRS. \textbf{NetLogo} slouží k~ABM simulacím s~entitami turtles, patches a~links
  a~\textbf{SPADE} kombinuje Python s~XMPP.
\item \textbf{LLM agenti} používají jazykový model jako \emph{politiku} (prompt $\to$ akce),
  jako \emph{plánovač}, jako \emph{kritika} ve schématu actor-critic a~jako rozhraní pro
  \emph{volání nástrojů}. Nasazují se hlavně v~agentových simulacích, ať už jde o~šíření zpráv
  sítí, nebo o~syntetický volební panel. Tvoří další rodinu mechanismů výběru akcí: nemají
  explicitní model světa a~sémantika jejich přesvědčení je neverifikovatelná, zato disponují
  širokou doménovou znalostí.
\end{itemize}

\section*{Závěrem: jak na zkoušku}
\addcontentsline{toc}{section}{Závěrem: jak na zkoušku}

Okruh je široký a~zkoušející obvykle ověřuje, nakolik uchazeč rozumí \emph{souvislostem}.
Otázky se točí kolem čtyř os:
\begin{enumerate}
\item \textbf{Reaktivita vs.\ proaktivita.} Tato osa vede od definice inteligentního agenta
  přes strategie odhodlání (bold vs.\ cautious) a~parametry $\gamma/\phi$ v~síti Maes až
  k~vrstvám hybridních architektur.
\item \textbf{Individuální vs.\ kolektivní racionalita.} Sobecký agent dovede systém k~$(D,D)$
  ve vězňově dilematu, k~tragédii obecní pastviny i~k~taktickému hlasování. Odtud plyne
  potřeba mechanism designu (Vickrey, VCG), tedy návrhu \emph{pravidel}, za nichž se
  i~sobeckému agentovi vyplatí být pravdomluvný.
\item \textbf{Symbolická vs.\ subsymbolická reprezentace.} Na jedné straně stojí FOL, STRIPS,
  ontologie a~ACL, na druhé subsumpce, aktivační sítě a~neuronové sítě v~RL. Praxe je
  hybridní.
\item \textbf{Teorie vs.\ praxe.} Proti elegantním, ale nekonstruktivním definicím
  (optimální agent, Nashova věta, Arrowova věta) stojí použitelné algoritmy (A*, CBS,
  Q-learning, CNET, JADE). K~tomu patří i~složitost: hledání NE je PPAD-úplné a~WDP
  i~optimální MAPF jsou NP-těžké.
\end{enumerate}

Zbývá výčet drobností, na které se ptají nejčastěji: rozdíl mezi \emph{desire}
a~\emph{intention}; proč efektem \texttt{inform} \emph{není}, že příjemce uvěří; proč je
$(D,D)$ jediné \emph{ne}-Pareto-optimální řešení vězňova dilematu; proč je pravdomluvnost
dominantní ve Vickreyově aukci; jaký je rozdíl mezi off-policy Q-learningem a~on-policy SARSA;
a~co dělá CBS na vysoké a~co na nízké úrovni.

\end{document}
